\documentclass{cours}
\usepackage{tikzlings}

\title{Le cours d'Ariane Mézard}
\author{Ariane Mézard}

\begin{document}
\begin{center}
    \begin{tikzpicture}
        \panda\thing[tophat, rotate = -1.3]
    \end{tikzpicture}
\end{center}
\part{Fonctions Holomorphes}
\section{Fonctions Analytiques}
\subsection{Séries Entières}    
\begin{définition}{Série Entière}{}
    Une série entière est une série de la forme $\sum_{n \in \N}a_{n}z^{n}$ où $z \in \C$ et $a_{n} \in \C$.\\
    Le domaine de convergence de la série entière est l'ensemble $\Delta$ des nombres complexes $z \in \C$ pour lesquels la série converge.  
\end{définition}
\begin{propositionfr}{Critère de Cauchy}{}
    Soient $a_{n}$ une suite complexe et $0 < r < r_{0}$. S'il existe $M > 0$ tel que 
    \begin{equation*}
        \abs{a_{n}}r_{0}^{n} \leq M, n \geq 0
    \end{equation*}
    alors $a_{n}z^{n}$ converge normalement sur $\overline{D}\left(0, r\right)$.
\end{propositionfr}
\begin{proof}
    Pour tout $n \in \N$ et $z \in \overline{D}\left(0, r\right)$ on a : 
    \begin{equation*}
        \abs{a_{n}z^{n}} \leq \abs{a_{n}}r^{n} \leq M\left(\frac{r}{r_{0}}\right)^{n}
    \end{equation*}
    Comme $0 < r < r_{0}$, $M \left(\frac{r}{r_{0}}\right)^{n}$ est le terme d'une série géométrique convergente.
\end{proof}

\begin{corollaire}{Rayon de Convergence}{}
    Soit $\sum_{n \in \N} a_{n}z^{n}$ une série entière et $R \in \R_{+} \cup \{+ \infty\}$ défini par 
    \begin{equation*}
        R = \sup \left\{r \geq 0 \text{ tel que la suite } \left(\abs{a_{n}}r^{n}\right)_{n \in \N} \text{ soit bornée}\right\}
    \end{equation*}
    Alors le domaine de convergence $\Delta$ de la série vérifie : 
    \begin{equation*}
        D\left(0, R\right) \subseteq \Delta \subseteq \overline{D}\left(0, R\right)
    \end{equation*}
\end{corollaire}

\begin{définition}{Rayon de Convergence}{}
    On appelle le nombre $R$ défini ci-dessus rayon de convergence.
\end{définition}

\begin{propositionfr}{Rayon d'Hadamard}{}
    Le rayon de convergence est donné par 
    \begin{equation*}
        R = \liminf_{n \to \infty}\frac{1}{\abs{a_{n}}^{1/n}}
    \end{equation*}
    Avec la convention $1/0 = \infty$
\end{propositionfr}

\begin{lemme}{Lemme d'Abel}{}
    Soit $u_{n}$ une suite réelle décroissante vers $0$ et $v_{n}$ une suite complexe telle que les sommes partielles $s_{n} = \sum_{k = 0}^{n} v_{k}$ soient bornées. Alors la série $\sum u_{n}v_{n}$ converge.
\end{lemme}

\begin{propositionfr}{Principe des Zéros Isolés}{}
    Soit $f\left(z\right) = \sum a_{n}z^{n}$ la somme d'une série entière de rayon de convergence $R > 0$. Si au moins un des coefficients $a_{n}$ n'est pas nul, il existe $r \in \left] 0, +\infty\right[$ tel que $f$ ne s'annule pas pour $\abs{z} \in \left]0, r\right[$.
\end{propositionfr}
\begin{proof}
    Soit $l = \min \{n \in \N, a_{n} \neq 0\}$, on a : 
    \begin{equation*}
        f\left(z\right) = \sum_{n \geq l} a_{n}z^{n} = z^{l}g\left(z\right)
    \end{equation*}
    avec $g\left(z\right) = a_{l} + a_{l + 1}z + \ldots$ et $g\left(0\right) \neq 0$. 
\end{proof}

\begin{définition}{Dérivée Complexe}{}
    Une fonction $f : U \to \C$ admet une dérivée par rapport à la variable complexe au point $z_{0}$ si
    \begin{equation*}
        \lim_{z \to u} \frac{f\left(z_{0} + u\right)- f\left(z_{0}\right)}{u}
    \end{equation*}
    existe. Cette limite est alors appelée dérivée de $f$ en $z_{0}$.
\end{définition}

\begin{propositionfr}{Dérivée d'une Série Entière}{}
    Soit $f\left(z\right) = \sum a_{n}z^{n}$ une série entière de rayon de convergence $R > 0$. Alors, pour tout $l \in \N^{\star}$, les dérivées $l$-ièmes de $f$ ont pour rayon de convergence $R$ et pour expression : 
    \begin{equation*}
        f^{\left(l\right)}\left(z\right) = \sum_{n \in \N}\frac{\left(n + l\right)!}{n!}a_{n + l}z^{n}  
    \end{equation*}
\end{propositionfr}

\begin{corollaire}{Primitive}{}
    Une série entière $f\left(z\right) = \sum a_{n}z^{n}$ de rayon de convergence $R > 0$ admet sur $D\left(0, R\right)$ une primitive complexe
    \begin{equation*}
        F\left(z\right) = \sum \frac{a_{n}}{n + 1}z^{n + 1}
    \end{equation*}
\end{corollaire}

\begin{propositionfr}
    Soit $f\left(z\right) = \sum a_{n}z^{n}$ une série entière de rayon de convergence $R > 0$. Soit $z_{0} \in D\left(0, R\right)$. La série entière 
    \begin{equation*}
        \sum_{n \in \N}\frac{1}{n!}f^{\left(n\right)}\left(z_{0}\right)\omega^{n}
    \end{equation*}
    a un rayon de convergence supérieur à $R - \abs{z_{0}}$ et pour tout $z \in D\left(z_{0}, R - \abs{z_{0}}\right)$,
    \begin{equation*}
        f\left(z\right) = \sum_{n \geq 0} \frac{1}{n !}f^{\left(n\right)}\left(z_{0}\right)\left(z - z_{0}\right)^{n}
    \end{equation*}
\end{propositionfr}

\subsection{Fonctions Analytiques}
\begin{définition}{Fonction Analytique}{}
    Une fonction $f : U \to \C$ est dite analytique si elle est DSE au voisinage de chaque point de $U$. 
\end{définition}

\begin{propositionfr}{Dérivabilité}{}
    Une fonction analytique sur un ouvert $U$ de $\C$ admet des dérivées de tous ordres qui sont des fonctions analytiques sur $U$. De plus, pour tout $z_{0} \in U$, $f$ est somme de sa série de Taylor en $z_{0}$ sur un voisinage de $z_{0}$.
\end{propositionfr}

\begin{corollaire}{Unicité du DSE}{}
    Une fonction analytique sur $U$ admet un unique développement en série entière au voisinage de chaque point de $U$.
\end{corollaire}

\begin{lemme}{Nullité}{}
    Si $U$ est connexe et $f$ est analytique sur $U$, nulle sur un ouvert non-vide de $U$, alors $f$ est identiquement nulle sur $U$. 
\end{lemme}

\begin{propositionfr}{Zéros Isolés}
    Soit $f$ une fonction analytique sur un ouvert connexe $U$. Si $f$ n'est pas identiquement nulle, ses zéros sont isolés, i.e. si $z_{0} \in U$ avec $f\left(z_{0}\right) = 0$, alors il existe $r > 0$ tel que $z_{0}$ soit le seul $z_{0}$ de $f$ sur $D\left(z_{0}, r\right)$
\end{propositionfr}

\begin{théorème}{Prolongement Analytique}{}
    Soit $U$ un ouvert connexe de $\C$, $f, g$ des fonctions analytiques sur $U$. Si $f, g$ coincident sur une partie $\Sigma$ de $U$ qui a un point d'accumulation dans $U$, alors elles coincident sur $U$. 
\end{théorème}

\begin{définition}{Primitive}{}
    Etant donnée une fonction analytique $f$ sur $U$, une fonction analytique $F$ de $U$ dans $\C$ est dite primitive de $f$ si $F'\left(z\right) = f\left(z\right)$ sur $U$. 
\end{définition}

\subsection{Détermination du Logarithme}
\begin{définition}{Détermination de l'Argument}{}
    Soit $U \subseteq \C^{\star}$ ouvert. Une fonction continue $\arg : U \to \R$ est dite détermination continue de l'argument sur $U$ si pour tout $z \in U$, $\exp\left(i\arg\left(z\right)\right)=\frac{z}{\abs{z}}$
\end{définition}
\begin{définition}{Détermination Principale}{}
    La détermination continue de l'argument 
    \begin{equation*}
        \begin{array}{rl}
            \C - \R_{-} \longrightarrow & \left] - \pi, \pi \right[\\
            z \mapsto & 2\arctan\left(\frac{y}{x + \sqrt{x^{2} + y^{2}}}\right)
        \end{array}
    \end{equation*}
    en prenant la racine carrée de $z$ appartenant au demi-plan $\Re z > 0$ est appelée détermination principale de l'argument.
\end{définition}
\begin{définition}{Logarithme}{}
    Soit $U \subseteq \C^{\star}$ ouvert. Une fonction continue $f : U \to \C$ est dite détermination du logarithme sur $U$ si 
    \begin{equation*}
        \forall z \in U, \exp\left(f\left(w\right)\right) = w
    \end{equation*}
\end{définition}

\begin{définition}{Détermination Principale du Log}{}
    On définit pour $\theta \in \R$ la fonction
    \begin{equation*}
        \log_{\theta} : \C \to \R_{-}e^{i\theta}, \log_{\theta}\left(w\right)=\log \abs{w} + i\arg_{\theta}\left(w\right)
    \end{equation*}
    La fonction $\log_{0}$ est appelée détermination principale du logarithme et notée $\log$. 
\end{définition}

\begin{propositionfr}{DSE du Logarithme}{}
    $\log$ est DSE sur $D\left(1, 1\right)$ et sur $D\left(0, 1\right)$ on a
    \begin{equation*}
        \log\left(1 + z\right) = \sum \frac{\left(-1\right)^{n + 1}}{n}z^{n}
    \end{equation*}
    Par conséquent, sur $D\left(z_{0}, \abs{z_{0}}\right)$, 
    \begin{equation*}
        g\left(z\right) = \log z_{0} + i\theta_{0} + \sum_{n \geq 1}\frac{\left(-1\right)^{n - 1}}{n}\left(\frac{z - z_{0}}{z_{0}}\right)^{n}
    \end{equation*}
    est une détermination analytique du logarithme. 
\end{propositionfr}

\begin{propositionfr}{Analycité des Déterminations}{}
    Il y a équivalence sur un ouvert connexe $U$ de $\C^{\star}$ pour une application continue $l$ entre : 
    \begin{itemize}
        \item $l$ est une détermination du logarithme à l'addition d'une constante près
        \item $l$ est une primitive analytique de $\frac{1}{z}$ sur $U$. 
    \end{itemize}
\end{propositionfr}

\begin{définition}{Détermination}{}
    Soit $U \subseteq \C^{\star}$ et $\alpha \in \C$. Une détermination continue de $z^{\alpha}$ est une application continue $g$ de $U$ dans $\C$ telle qu'il existe une détermination du logarithme $l\left(z\right)$ de $z$ telle que $g\left(z\right) = \exp^{\alpha l\left(z\right)}$.
\end{définition}

\section{Théorie de Cauchy}
\subsection{Homotopie et Simple Connexité}
\begin{définition}{Chemin}{}
    Soit $\left[a, b\right]$ un intervalle de $\R$. Un chemin $\gamma: \left[a, b\right] \to \C$ est une application continue. Le point $\gamma\left(a\right)$ est appelé origine et le point $\gamma\left(b\right)$ est dit extrémité. On orientera par défaut un chemin dans le sens des paramètres croissants. Si $\gamma\left(a\right) = \gamma\left(b\right)$, le chemin est dit lacet d'origine $\gamma\left(a\right)$.
\end{définition}

\begin{définition}{Opérations}{}
    \begin{enumerate}
        \item Si $\gamma$ est constant, son image est réduite à un point. Il est alors appelé chemin (ou lacet) constant.
        \item Soit $\alpha \in \R^{\star}$, $\gamma : t \in \left[0, 1\right] \mapsto e^{2i\pi \alpha t}$ est un chemin dont l'image est une partie du cercle unité $\partial D\left(0, 1\right)$. Si $\alpha = n \in \Z^{\star}$, $\gamma\left(\left[0, 1\right]\right)$ est le cercle tout entier parcouru $n$ fois. 
        \item Si $\gamma : \left[a, b\right] \to \C$ est un chemin, le chemin opposé 
        \begin{equation*}
            \gamma^{0} : t \in \left[a, b\right] \mapsto \gamma\left(a + b - t\right)
        \end{equation*}
        est $\gamma$ parcouru en sens inverse. 
        \item La juxtaposition de $\gamma_{1}, \gamma_{2}$ tels que $\gamma_{1}\left(b\right) = \gamma_{2}\left(c\right)$ est le chemin $\gamma = \gamma_{1} \land \gamma_{2} : \left[a, d + b - c\right] \to \C$
        \begin{equation*}
            \gamma\left(t\right) = \begin{cases}
                \gamma_{1}\left(t\right) & \text{ pour } a\leq t \leq b\\
            \gamma_{2}\left(t - b + c \right) & \text{ pour } b \leq t \leq d + b - c
            \end{cases}
        \end{equation*}
    \end{enumerate}
\end{définition}

\begin{définition}{Homotopie}{}
    Soit $U$ un ouvert de $\C$, $\gamma_{i} : I \to U$, $i \in \{1, 2\}$ deux chemins. Une homotopie de $\gamma_{1}$ à $\gamma_{2}$ dans $U$ est une application continue $\phi$ de $I \times J$ dans $U$ où $I = \left[a, b\right]$ et $J = \left[c, d\right]$ sont deux intervalles de $\R$ telle que : 
    \begin{equation*}
        \phi\left(t, c\right) = \gamma_{1}\left(t\right) \text{ et } \phi\left(t, d\right) = \gamma_{2}\left(t\right), t\in I
    \end{equation*}
\end{définition}

\begin{définition}{Simple Connexité}{}
    Un espace topologique $X$ connexe par arcs est dit simplement connexe si tout lacet dans $X$ est homotope à un point dans $X$. 
\end{définition}

\begin{propositionfr}{}{}
    \begin{itemize}
        \item Un espace topologique est simplement connexe si et seulement si tous les chemins de même extrémités sont homotopes.
        \item Un ouvert étoilé par rapport à un point est simplement connexe. En particulier, dans $\C$, le plan, un demi-plan, un disque ouvert, l'intérieur d'un rectangle ou d'un triangle sont simplement connexes. 
        \item Le demi-plan ouvert $\Im z > 0$ auquel nous ôtons un nombre fini de demi-droites fermées $z = t +i\beta_{k}$, $t\in \left]-\infty, \alpha_{k}\right]$ est simplement connexe non étoilé. 
        \item $\C^{\star}$ n'est pas simplement connexe car le cercle unité n'est pas homotope à un chemin constant. 
    \end{itemize}
\end{propositionfr}

\subsection{Intégrales sur un Chemin}
Dorénavant, les chemins sont supposés $\cont^{1}$ par morceaux. 

\begin{définition}{Equivalence de Chemins}{}
    Deux chemins $\gamma_{i} : I_{i} \to \C$ sont dits équivalents s'il existe une bijection croissante $\phi : I_{2} \to I_{1}$ continue de réciproque continue et $\cont^{1}$ par morceaux telle que : 
    \begin{equation*}
        \gamma_{2}\left(t\right) = \gamma_{1}\left(\phi\left(t\right)\right), t\in I_{2}
    \end{equation*}
\end{définition}

\begin{définition}{Intégrale le long d'un Chemin}{}
    Soit $f : U \to \C$ continue et $\gamma : I = \left[a, b\right] \to \C$ un chemin avec $\gamma\left(I\right) \subseteq U$. Alors, la fonction $t : f\left(\gamma\left(t\right)\right)\gamma'\left(t\right)$ est continue par morceaux dans $\left[a, b\right]$. On appelle intégrale de $f$ le long du chemin $\gamma$ : 
    \begin{equation*}
        \int_{\gamma}f\left(z\right)\d z = \int_{a}^{b}f\left(\gamma\left(t\right)\right)\gamma'\left(t\right)\d t
    \end{equation*}
\end{définition}

\begin{définition}{Longueur}{}
    La longueur d'un chemin est le réel : 
    \begin{equation*}
        long\left(\gamma\right) = \int_{a}^{b}\abs{\gamma^{'}\left(t\right)}\d t
    \end{equation*}
\end{définition}

\begin{propositionfr}{Propriétés}{}
    \begin{itemize}
        \item Si $F$ est une primitive de $f$, pour tout chemin $\gamma$ : 
        \begin{equation*}
            \int_{\gamma}f\left(z\right)\d z = F\left(\gamma\left(b\right)\right) - F\left(\gamma\left(a\right)\right)
        \end{equation*}
        \item Si $\gamma_{1} \sim \gamma_{2}$ alors 
        \begin{equation*}
            \int_{\gamma_{1}} f = \int_{\gamma_{2}} f
        \end{equation*}
        \item Si $\left[Z_{0}, z_{1}\right] \subseteq U$, nous notons $\int_{\left[z_{0}, z_{1}\right]} f\left(z\right)\d z = \int_{\gamma}f\left(z\right)\d z$ où $\gamma : t\in \left[0, 1\right] \mapsto \left(1 - t\right)z_{0} + tz_{1}$.
        \item Si $\partial D\left(z_{0}, r\right) \subseteq U$, soit le lacet $\gamma : \theta \in \left[0, 2\pi\right] \mapsto z_{0} + re^{i\theta}$. On a : 
        \begin{equation*}
            \int_{\gamma}f\left(z\right) \d z = \int_{\partial D\left(z_{0}, r\right)}f\left(z\right)\d z = \int_{0}^{2\pi}f\left(z_{0} + re^{i\theta}\right)ire^{i\theta}\d\theta
        \end{equation*}
        \item En séparant parties réelles et imaginaires, $f = P + iQ$ et $\gamma = u + iv$, on a :
        \begin{equation*}
            \begin{aligned}
                \int_{\gamma}f\left(z\right) \d z = & \int_{a}^{b}\left(\left(P \circ \gamma\right)u' - \left(Q \circ \gamma\right)v'\right)\d t + i \int_{a}^{b}\left(\left(Q \circ \gamma\right)u' + \left(P \circ \gamma\right)u'\right)\d t\\
                = &\int_{\gamma}\left(P\d x - Q\d y\right) + i \int_{\gamma}\left(P\d y + Q \d x\right)
            \end{aligned}
        \end{equation*}
        \item On a :
        \begin{equation*}
            \int_{\gamma} f\left(z\right) \d z = - \int_{\gamma^{0}}f\left(z\right) \d z
        \end{equation*}
        \item On a :
        \begin{equation*}
            \abs{\int_{\gamma}f\left(z\right)\d z} \leq long\left(\gamma\right)\max_{\gamma}\abs{f}
        \end{equation*} 
    \end{itemize}
\end{propositionfr}

\subsection{Théorème de Cauchy}
\begin{théorème}{de Cauchy}{cauchy}
    Soit $U \subseteq \C$ un ouvert connexe et $f$ une fonction analytique dans $U$. Si $\gamma_{1}, \gamma_{2}$ sont deux lacets homotopes dans $U$, alors
    \begin{equation*}
        \int_{\gamma_{1}}f\left(z\right) \d z = \int_{\gamma_{2}}f\left(z\right)\d z
    \end{equation*}
    En particulier, si $U$ est simplement connexe, l'intégrale sur un lacet de $f$ est nulle. 
\end{théorème}

\begin{théorème}{Primitives Analytiques}{primitive-analytique}
    Soit $U \subseteq \C$ un ouvert simplement connexe. 
    \begin{enumerate}
        \item Toute fonction analytique dans $U$ admet une primitive.
        \item Si $f : U \to \C^{\star}$ est analytique, alors il existe $g : U \to \C$ analytique tel que $\exp\left(g\right) = f$ sur $U$. 
    \end{enumerate}
\end{théorème}

\subsection{Formule de Cauchy}
\begin{lemme}{Intégrité de l'Indice}{}
    Soit $\gamma : I = \left[c, d\right] \to \C$ un lacet et $a\notin \gamma\left(I\right)$. Alors
    \begin{equation*}
        j\left(a, \gamma\right) = \frac{1}{2i\pi}\int_{\gamma}\frac{\d z}{z - a} \in \Z
    \end{equation*}
\end{lemme}
\begin{proof}
    Pour $t \in \left[c, d\right]$ on pose 
    \begin{equation*}
        h\left(t\right) = \int_{c}^{t}\frac{\gamma'\left(s\right)\d s}{\gamma\left(s\right) - a}
    \end{equation*}
    On a $h'\left(t\right) = \frac{\gamma'\left(t\right)}{\gamma\left(t\right)-a}$, sauf en un nombre fini de points de $I$. \\
    Remarquons que $g\left(t\right) = e^{-h\left(t\right)}\left(\gamma\left(t\right) - a\right)$ a pour dérivée
    \begin{equation*}
        g'\left(t\right)= - h'\left(t\right)e^{-h\left(t\right)}\left(\gamma\left(t\right)- a\right) + \gamma'\left(t\right)e^{-h\left(t\right)} = 0
    \end{equation*}
    sauf en un nombre fini de points de $I$. Comme $g$ est continue, elle est constante et $g\left(c\right) = g\left(d\right)$. \\
    Or, $h\left(c\right) = 0$ donc $g\left(c\right) = \gamma\left(c\right) - a = g\left(d\right) = e^{-h\left(d\right)}\left(\gamma\left(d\right) - a\right)$. Mais $\gamma$ est un lacet, donc $\gamma\left(c\right) = \gamma\left(d\right)$. Donc $h\left(d\right) = 2in\pi$. Donc $j\left(a, \gamma\right) = n \in \Z$. 
\end{proof}
\begin{définition}{Indice}{}
    L'entier $j\left(a, \gamma\right)$ est appelé indice de $a$ par rapport au lacet $\gamma$ et s'interprète comme le nombre de fois que le lacet tourne autour de $a$ lorsque $a$ est intérieur au lacet.
\end{définition}

\begin{propositionfr}{Propriétés}{}
    \begin{enumerate}
        \item Soit $\gamma, \gamma_{1}, \gamma_{2}$ des lacets de même origine dont les lacets ne contiennent pas $a$. Alors,
        \begin{equation*}
            j\left(a, \gamma^{0}\right) = -j\left(a, \gamma\right) \text{ et } j\left(a, \gamma_{1} \land \gamma_{2}\right) = j\left(a, \gamma_{1}\right) + j\left(a, \gamma_{2}\right)
        \end{equation*}
        \item En appliquant le théorème de Cauchy à la fonction analytique $1/\left(z - a\right)$ dans $\C - \{a\}$, nous obtenons $j\left(a, \gamma_{1}\right) = j\left(a, \gamma_{2}\right)$ si $\gamma_{1}, \gamma_{2}$ sont homotopes dans $\C - \{a\}$.
        \item Soit $U \subset \C$ un ouvert simplement connexe et $\gamma \subset U$. Si $a \notin U$, alors $j\left(a, \gamma\right) = 0$.
        \item Si $\gamma$ set un lacet dans $\C$, pour tout ouvert connexe $U$ de $\C - \gamma\left(I\right)$, la fonction $z \mapsto j\left(z, \gamma\right)$ est constante dans $U$.
        \item Soit $\gamma_{n} : t \mapsto e^{int}$, on a : 
        \begin{equation*}
            j\left(z_{0}, \gamma_{n}\right) = \begin{cases}
                n & si \abs{z_{0}} < 1\\
                0 & si \abs{z_{0}} > 1
            \end{cases}
        \end{equation*}
    \end{enumerate}
\end{propositionfr}
\begin{proof}[Démonstration du point iv.]
    Soit $z \in D\left(z_{0}, r\right) \subseteq U$, 
    \begin{equation*}
        j\left(z, \gamma\right) = \frac{1}{2i\pi}\int_{\gamma}\frac{\d u}{u - z} = \frac{1}{2i\pi}\int_{\gamma_{1}}\frac{\d u}{u - z} = \frac{1}{2i\pi}\int_{\gamma}\frac{\d u}{u - z_{0}} = j\left(z_{0}, \gamma\right)
    \end{equation*}
    pour $\gamma_{1} : t \mapsto \gamma\left(t\right) + \left(z - z_{0}\right)$ qui est homotopie à $\gamma$ via 
    \begin{equation*}
        \phi\left(t, s\right) = \gamma\left(t\right) + s\left(z - z_{0}\right), 0 \leq s \leq 1
    \end{equation*}
    Donc $j\left(\cdot, \gamma\right)$ est localement constante donc constante sur $U$ connexe. 
\end{proof}

\begin{théorème}{Formule de Cauchy}{}
    Soit $U \subseteq \C$ un ouvert simplement connexe, $\gamma : I \to U$ un lacet dans $U$. Soit $f$ analytique sur $U$. Pour tout $w \in U \setminus \gamma\left(I\right)$
    \begin{equation*}
        j\left(w, \gamma\right)f\left(w\right) = \frac{1}{2i\pi}\int_{\gamma}\frac{f\left(z\right)}{z - w}\d z
    \end{equation*}
\end{théorème}
\begin{proof}
    La fonction
    \begin{equation*}
        g : z \in U \mapsto \begin{cases}
            \frac{f\left(z\right)-f\left(w\right)}{z-w} & \text{ si } z \neq w\\
            f'\left(w\right) & \text{ si } z = w
        \end{cases}
    \end{equation*}
    est analytique sur $U$. En effet pour $r > 0$ assez petit, $f$ admet un développement de Taylor sur $D\left(w, r\right) \subseteq U$ et donc pour $z \in D\left(w, r\right)$ :
    \begin{equation*}
        g\left(z\right) = f'\left(w\right) + \frac{f''\left(w\right)}{2!}\left(z - w\right) + \cdot + \frac{f^{\left(n\right)}\left(w\right)}{n!}\left(z - w\right)^{n - 1} + \cdot
    \end{equation*}
    Comme $U$ est simplement connexe, le théorème de Cauchy donne $\int_{\gamma} g = 0$ et comme $w \notin \gamma\left(I\right)$, $\int_{\gamma}\frac{f\left(z\right) - f\left(w\right)}{z - w}\d z = 0$ c'est à dire : 
    \begin{equation*}
        \int_{\gamma}\frac{f\left(z\right)\d z}{z - w} = f\left(w\right)\int_{\gamma}\frac{\d z}{z - w} = 2i\pi j\left(w, \gamma\right)f\left(w\right)
    \end{equation*}
\end{proof}

\begin{corollaire}{Valeur en un point}{}
    On a : 
    \begin{equation*}
        f\left(w\right) = \frac{1}{2i\pi}\int_{\partial D\left(z_{0}, r\right)}\frac{f\left(z\right)}{z-w}\d z, w \in D\left(z_{0}, r\right)
    \end{equation*}
\end{corollaire}

\begin{propositionfr}{Continuité sur un Lacet}{cont_lacet}
    Soit $\gamma : I = \left[c, d\right] \to \C$ un lacet et $g : \gamma\left(I\right) \to \C$ une fonction définie et continue sur $\gamma\left(I\right)$. Alors : 
    \begin{equation*}
        f\left(z\right) = \int_{\gamma}\frac{g\left(u\right)\d u}{u - z}
    \end{equation*}
    est définie et analytique dans $\C \setminus \gamma\left(I\right)$.\\
    Précisément, pour tout $w \in \C \setminus \gamma\left(I\right)$ pour tout $n \in \N$ et 
    \begin{equation*}
        c_{n} = \int_{\gamma} \frac{g\left(u\right)\d u}{\left(u - w\right)^{n + 1}}
    \end{equation*}
    nous avons un développement en série entière convergente 
    \begin{equation*}
        f\left(z\right) = \sum_{n \geq 0}c_{n}\left(z - w\right)^{n}
    \end{equation*}
    dans tout disque ouvert de centre $w$ et de rayon $r = d\left(w, \gamma\left(I\right)\right)$ et 
    \begin{equation*}
        f^{\left(n\right)}\left(w\right) = n!c_{n} = n!\int_{\gamma}\frac{g\left(u\right)\d u}{\left(u - w\right)^{n + 1}}
    \end{equation*}
\end{propositionfr}
\begin{proof}
    Pour tout $u \in \gamma\left(I\right), z \in D\left(w, qr\right), q \in \left[0, 1\right]$, la série 
    \begin{equation*}
        \frac{1}{u - z} = \frac{1}{u - w} \frac{1}{1 - \frac{z-w}{u-w}} = \sum_{n = 0}^{+\infty}\frac{\left(z - w\right)^{n}}{\left(u - w\right)^{n + 1}}
    \end{equation*}
    est convergente. Comme $\left(g \circ \gamma\right)\gamma'$ est continue par morceaux sur $\left[c, d\right]$ il existe $M$ tel que 
    \begin{equation*}
        \abs{g\left(\gamma\left(t\right)\right)\gamma'\left(t\right)} \leq M
    \end{equation*}
    Donc :
    \begin{equation*}
        \abs{g\left(\gamma\left(t\right)\right)\gamma'\left(t\right)\frac{\left(z - w\right)^{n}}{\left(\gamma\left(t\right) - w\right)^{n + 1}}} \leq M \frac{q^{n}}{r}, t\in \left[c, d\right]
    \end{equation*}
    Finalement, la série sous l'intégrale est normalement convergente et : 
    \begin{equation*}
        f\left(z\right) = \int_{c}^{d}\frac{g\left(\gamma\left(t\right)\right)\gamma'\left(t\right)\d t}{\gamma\left(t\right) - z} = \int_{c}^{d}g\left(\gamma\left(t\right)\right)\gamma'\left(t\right)\left(\sum_{n = 0}^{+\infty}\frac{\left(z - w\right)^{n}}{\left(\gamma\left(t\right) - w\right)^{n + 1}} \right)\d t
    \end{equation*}
    et donc $f\left(z\right) = \sum_{n = 0}^{+ \infty} c_{n}\left(z - w\right)^{n}$
\end{proof}

\begin{propositionfr}{Dérivée $n$-ième}{}
    Soit $f$ analytique sur $U$ et $\gamma$ le bord de $\overline{D}\left(w, r\right) \subseteq U$. D'après la formule de Cauchy : 
    \begin{equation*}
        f^{\left(n\right)}\left(w\right) = \frac{n!}{2\pi r^{n}}\int_{0}^{2\pi}\frac{f\left(w + re^{it}\right)}{e^{nit}}\d t
    \end{equation*}
\end{propositionfr}

\begin{corollaire}{}{taylor_ser}
    \begin{enumerate}
        \item Soit $f$ analytique sur $U$. Pour tout $a \in U$, la série de Taylor de $f$ au voisinage de $a$ est convergente et a pour somme $f\left(z\right)$ dans le plus grand disque ouvert de centre $a$ contenu dans $U$
        \item Si $f$ est analytique sur $\C$, sa série de Taylor en tout point de $\C$ est convergente sur $\C$. 
    \end{enumerate}
\end{corollaire}
\begin{proof}
    On applique la formule de Cauchy sur le contour $\gamma$ d'un disque $D\left(a, r\right)$ contenu dans $U$. Pour $z \in D\left(a, r\right), j\left(z, \gamma\right) = 1$ et 
    \begin{equation*}
        f\left(z\right) = \frac{1}{2i\pi}\int_{\gamma}^{}\frac{f\left(w\right)}{w -z}\d z
    \end{equation*}
    La proposition \ref{propo:cont_lacet} donne un développement en série entière de $f$ en $z - a$ convergeant sur $D\left(a, r\right)$. Par unicité du développement, il s'agit de la série de Taylor. En faisant tendre $r$ vers $d\left(a, \C - U\right)$, nous obtenons le résultat annoncé. 
\end{proof}

\begin{corollaire}{Constance Locale}{constance_locale}
    Supposons $U$ connexe, $a \in U$ et $f : U \to \C$ analytique. Si pour tout $k > 0, f^{\left(k\right)}\left(a\right) = 0$, alors $f$ est constante sur $U$. 
\end{corollaire}
\begin{proof}
    D'après le corollaire \ref{cor:taylor_ser}, $f$ est localement somme de sa série de Taylor. Donc $f$ est constante sur un ouvert contenant $a$. Soit $\Omega = \{w \in U, \forall k > 0, f^{\left(k\right)}\left(w\right) = 0\}$. Cet ensemble est ouvert, non vide, et fermé. Par connexité de $U$, $\Omega = U$, $f' = 0$ sur $U$ et $f$ est constante sur $U$.
\end{proof}

\begin{théorème}{Multiplicité}{multiplicité}
    Soit $f : U \to \C$ analytique non constante au voisinage de $a \in U$. Si $f\left(a\right) = 0$, il existe un unique entier $m \geq 1$ et $g : V \to \C$ analytique sur un voisinage $V$ de $a$ tels que 
    \begin{equation*}
        f\left(z\right) = \left(z - a\right)^{m}g\left(z\right), g\left(a\right) \neq 0, z\in V
    \end{equation*}
    En particulier, le point $a$ possède un voisinage dans lequel il est l'unique zéro de $f$.
\end{théorème}
\begin{proof}
    D'après le corollaire \ref{cor:constance_locale}, si $f$ n'est pas constante dans un voisinage de $a$, il existe $m \geq 1$ tel $f^{\left(m\right)}\left(a\right) \neq 0$ et $f'\left(a\right) = \ldots = f^{\left(m-1\right)}\left(a\right) = 0$.\\
    Comme $f\left(a\right) = 0$, on peut alors factoriser $\left(z -a\right)^{m}$ dans le développement en série de Taylor de $f$ en $a$.
\end{proof}

\begin{définition}{Ordre}{}
    L'entier $m$ du théorème précédent est dit ordre de $f$ en $a$, noté $ord\left(f, a\right)$.
\end{définition}

\subsection{Inégalités de Cauchy, Premières Applications}
\begin{propositionfr}{Inégalités de Cauchy}{}
    Soit $f : U \to \C$ analytique, $\overline{D}\left(w, r\right) \subset U, r > 0$. On a, pour $n \in \N$ : 
    \begin{equation*}
        \abs{f^{\left(n\right)}\left(w\right)}\leq \frac{n!}{r^{n}}\sup_{z \in \partial D\left(w, r\right)}\abs{f\left(z\right)}
    \end{equation*}
\end{propositionfr}
\begin{proof}
    On a :
    \begin{equation*}
        f^{\left(n\right)}\left(w\right) = \frac{n!}{2\pi r^{n}}\int_{0}^{2\pi}\frac{f\left(w + re^{it}\right)}{e^{nit}} \d t
    \end{equation*}
    On en déduit immédiatement le résultat.
\end{proof}
\begin{lemme}{Bornitude et Polynomialité}{}
    Soit $f$ analytique sur $\C$. Supposons qu'il existe $A, B \geq 0$ tels que 
    \begin{equation*}
        \forall z \in \C, \abs{f\left(z\right)}\leq A\left(1 + \abs{z}\right)^{B}
    \end{equation*} 
    Alors $f$ est un polynôme de degré $\leq B$. 
\end{lemme}
\begin{proof}
    Soit $n \geq \lfloor B \rfloor + 1 > B$. Par les inégalités de Cauchy, puisque 
    \begin{equation*}
        \sup_{\partial D\left(z, r\right)} \abs{f\left(z\right)} \leq A\left(1 + \abs{z} + r\right)^{B}
    \end{equation*}
    on a : 
    \begin{equation*}
        \abs{f^{\left(n\right)}\left(z\right)} \leq \frac{n!}{r^{n}}A\left(1 + \abs{z} + r\right)^{B}
    \end{equation*}
    En faisant tendre $r$ vers $+\infty$, par croissance comparée, $f^{\left(n\right)}\left(w\right) = 0$ pour $n \geq B$. Localement, $f$ étant somme de sa série de Taylor, c'est localement un polynôme de degré au plus $B$, ce qui est donc le résultat.
\end{proof}

\begin{théorème}{Liouville}{}
    Une fonction analytique bornée sur $\C$ est constante.
\end{théorème}

\begin{théorème}{d'Alembert-Gauss}{}
    Tout polynôme $P \in \C\left[z\right]$ de degré $\geq 1$ admet une racine dans $\C$.
\end{théorème}
\begin{proof}
    Par l'absurde, si $P\left(z\right) = \sum_{i = 0}^{d}a_{i}z^{i}$ ne s'annule pas, $f = 1/P$ est analytique sur $\C$ et $\abs{f\left(z\right)} \sim \frac{1}{\abs{a_{d}}\abs{z}^{d}}$ tend vers $0$ quand $\abs{z}$ tend vers $+\infty$. En particulier, $f$ est bornée sur $\C$ donc constante d'après le théorème de Liouville. Ainsi, $P = 1/f$ est constant, ce uqui est absurde.
\end{proof}

\begin{théorème}{Topologie}{}
    Les ouverts $\C$ et $D\left(0, 1\right)$ sont homéomorphes mais pas isomorphes. 
\end{théorème}


\section{Fonctions Holomorphes}
\subsection{Définitions}
\begin{définition}{Holomorphie}{}
    Une fonction $f : U \to \C$ est dite holomorphe en $z_{0} \in U$ si la limite \begin{equation*}\lim_{h \in \C \to 0} \frac{f\left(z_{0} + h\right) - f\left(z_{0}\right)}{h}\end{equation*} existe. On la note $f'\left(z_{0}\right)$.\\
    On définit $\O\left(U\right)$ l'ensemble des fonctions holomorphes.
\end{définition}

\begin{propositionfr}{Exemples Holomorphe}{}
    \begin{itemize}
        \item Si $f$ est constante, $f$ est holomorphe et $f' = 0$
        \item Si $f$ est un polynôme, $f$ est holomorphe
        \item Si $f$ est analytique, $f$ est holomorphe
        \item $\sin, \cos, \exp, \tan$ sont holomorphes sur $\C$.
        \item $z \mapsto \bar{z}$ n'est pas holomorphes en aucun point : 
        \begin{equation*}
            \frac{f\left(z + h\right) - f\left(z\right)}{h} = \frac{\bar{h}}{h}
        \end{equation*}
        n'a pas de limite en $0$. 
        \item $f\left(z\right) = \abs{z}^{2}$ n'est holomorphe que pour $z = 0$ : 
        \begin{equation*}
            \frac{\left(z + h\right)\left(\bar{z} + \bar{h}\right) - z\bar{z}}{h} = \frac{h\bar{z} + \bar{h}z + h\bar{h}}{h}
        \end{equation*}
        n'a une limite que si $z = 0$. 
    \end{itemize}
\end{propositionfr}

\subsection{Différentiabilité}
\begin{définition}{Forme Différentielle}{}
    Une $1$-forme différentielle sur $\Omega$ est une application $\alpha : \Omega \to Hom_{\R}\left(\R^{n}, \C\right)$.\\
    En particulier, les $\d x_{i} \in Hom_{\R}\left(\R^{n}, \C\right)$ qui à $a \mapsto \d x_{i}\left(a\right) = a_{i}$ permettent d'écrire : 
    \begin{equation*}
        \alpha\left(x\right) = \sum_{i = 1}^{n} \alpha_{i}\left(x\right)\d x_{i}
    \end{equation*}
    où $\alpha_{i} : \Omega \to \C$. On a alors : 
    \begin{equation*}
        \alpha\left(x\right)\left(a\right) = \sum_{i = 1}^{n}\alpha_{i}\left(x\right)\d x_{i}\left(a\right)
    \end{equation*}
    On dit que $\alpha$ est de classe $\cont^{k}$ si et seulement si tous les $\alpha_{i}$ sont de classe $\cont^{k}$.
\end{définition}

\begin{définition}{$\R$ Différentiabilité}{}
	Une fonction $f$ d'un ouvert connexe $\Omega$ de $\R^{n}$ est dite $\R$-différentiable sur $ \Omega$ si et seulement si il existe une $1$-forme différentielle $\d f : \Omega \to Hom_{\R}\left(\R^{n}, \C\right)$ telle que 
	\begin{equation*}
		f\left(z + h\right) = f\left(z\right) + \d f\left(z\right)\left(h\right) + o\left(h\right)
	\end{equation*}
	On pose $\d_{x}f = \sum_{i = 1}^{n}\frac{\partial f}{\partial x_{i}}\left(x\right)\d x_{i}$.
\end{définition}

Dans la suite on travaille dans $\C \simeq \R^{2}$ et pour $h \in \C$, on note $h = k + il = \left(k, l\right)$ et pour $z \in U = \Omega, z = x+ iy = \left(x, y\right)$.

\begin{propositionfr}{Différentielle dans une base}{}
	Soit $f : U \to \C$ différentiable de différentielle $\d f = \frac{\partial f}{\partial x}\d x + \frac{\partial f}{\partial y} \d y$. On a $\forall z \in U$ : 
	\begin{equation*}	
		\d_{z}f = \partialderiv{f}{x}\left(z\right)\d x + \partialderiv{f}{y}\left(z\right) \d y
	\end{equation*}
	En $h = k + il$ : 
	\begin{equation*}
		\d_{z}f\left(h\right) = \partialderiv{f}{x}\left(z\right)k + \partialderiv{f}{y}\left(z\right)l
	\end{equation*}
	On définit 
	\begin{equation*}
		\d z = \d x + i\d y \text{ et } \d\bar{z} = \d x - i\d y
	\end{equation*}
	On a alors : 
	\begin{equation*}	
		\d f = \frac{1}{2}\left(\partialderiv{f}{x} - i\partialderiv{f}{y}\right)\d z + \frac{1}{2}\left(\partialderiv{f}{x} + i\partialderiv{f}{y}\right)\d\bar{z}
	\end{equation*}
	ce qu'on écrit aussi : 
	\begin{equation*}	
		\d f = \partialderiv{f}{z}\d z + \partialderiv{f}{\bar{z}}\d\bar{z}
	\end{equation*}
	On a par ailleurs 
	\begin{equation*}
		\overline{\left(\partialderiv{f}{z}\right)} = \partialderiv{\bar{f}}{\bar{z}}
	\end{equation*}
\end{propositionfr}

\begin{propositionfr}{Exemples}{}
	\begin{enumerate}
		\item Si $f\left(z\right) = z$, $\partialderiv{z}{z} = \partialderiv{f}{z} = 1$ et $\partialderiv{z}{\bar{z}} = 0$. A l'inverse, $\partialderiv{\bar{z}}{z} = 0$.
		\item Pour $P\left(x, y\right) = \sum_{0 \leq \alpha, \beta\leq d}c_{\alpha, \beta}x^{\alpha}y^{\beta}$. En notant $x = \frac{z + \bar{z}}{2}$ et $y = \frac{z - \bar{z}}{2i}$, on a : 
			\begin{equation*}
				P\left(z\right) = \sum_{\alpha, \beta}a_{\alpha, \beta}z^{\alpha}\bar{z}^{\beta}
			\end{equation*}
			où on a 
			\begin{equation*}	
				a_{\alpha, \beta} = \frac{1}{\alpha!\beta!}\frac{\partial^{\alpha + \beta}}{\partial z^{\alpha}\partial z^{\beta}}P\left(0\right)
			\end{equation*}
			On retrouve que $P$ est holomorphe si on a $a_{\alpha, \beta} = 0$ pour $\beta \geq 1$. 
	\end{enumerate}
\end{propositionfr}

\begin{théorème}{Lien $\C$-dérivabilité et $\R$-différentiabilité}{}
	Soit $f : U \to \C$. On a équivalence entre : 
	\begin{enumerate}
		\item $f \in \O\left(U\right)$
		\item $f$ est $\R$-différentiable sur $U$ et $\d_{z} f$ est $\C$-linéaire pour tout $z \in U$
		\item $f$ est $\R$-différentiable sur $U$ et $\partialderiv{f}{\bar{z}} = 0$ pour tout $z \in U$
	\end{enumerate}
\end{théorème}

\begin{proof}
	\begin{description}
		\item[$i \Rightarrow ii$] $f\left(z + h\right) = f\left(z\right) + hf'\left(z\right) + o\left(h\right) \Longrightarrow f$ est $\R$-différentiable en $z$ et $\d_{z}f : \R^{2}\to \C$ qui à $h \mapsto hf'\left(z\right)$ est $\C$-linéaire.
		\item[$ii \Rightarrow iii$] On a : 
			\begin{equation*}
				\begin{aligned}
					\d_{z}f\left(h\right) = & \partialderiv{f}{z}\left(z\right)\d z\left(h\right) + \partialderiv{f}{\bar{z}}\left(z\right)\d \bar{z}\left(h\right)\\
					= & \partialderiv{f}{z}\left(z\right)h + \partialderiv{f}{\bar{z}}\left(z\right)\bar{h}\\
					\d_{z}f\left(h\right) =& \partialderiv{f}{z}h + \partialderiv{f}{\bar{z}}\bar{h}
				\end{aligned}
			\end{equation*}
			On a alors : $\d_{z}f\left(ih\right) = \partialderiv{f}{z}ih -i\partialderiv{f}{\bar{z}}\bar{h}$. Mais $\d_{z}f$ est $\C$-linéaire par hypothèse. Donc : 
			\begin{equation*}
				\d_{z}f\left(ih\right) = i\d_{z}f\left(h\right) = i\partialderiv{f}{z}h + i\partialderiv{f}{\bar{z}}\bar{h}
			\end{equation*}
			Ainsi : $\partialderiv{f}{\bar{z}} = - \partialderiv{f}{\bar{z}} = 0$.
		\item[$iii \Rightarrow i$] On a : 
			\begin{equation*}
				\d_{z}f\left(h\right) = \partialderiv{f}{z}h
			\end{equation*}
			D'où 
			\begin{equation*}
				f\left(z + h\right) = f\left(z\right) + \partialderiv{f}{z}h + o\left(h\right)
			\end{equation*}
			Ainsi : 
			\begin{equation*}	
				\lim_{h \to 0}\frac{f\left(z + h\right) - f\left(z\right)}{h} = \partialderiv{f}{z}
			\end{equation*}
			et $f$ est holomorphe en $z$. 
	\end{description}
\end{proof}

\begin{propositionfr}{Équations de Cauchy-Riemann}{}
	On note $f\left(x +iy\right) = P\left(x, y\right) + iQ\left(x, y\right)$ où $P, Q : \R^{2} \to \R$. Si $f$ est holomorphe, on a :  
	\begin{equation*}	
		\partialderiv{f}{x} + i\partialderiv{f}{y} = 0
	\end{equation*}
	i.e. 
	\begin{equation*}
		\begin{cases}
			\partialderiv{P}{x} - \partialderiv{Q}{y} =& 0\\
			\partialderiv{Q}{x} + \partialderiv{P}{y} =& 0
		\end{cases}
	\end{equation*}
	Ce sont les équations de Cauchy-Riemann.
\end{propositionfr}
\begin{proof}
	On a : 
	\begin{equation*}
		\d_{z}f\left(h\right) = \partialderiv{f}{x}\left(z\right)\d x\left(h\right) + \partialderiv{f}{y}\left(z\right)\d y\left(h\right) = f'\left(z\right)h = f'\left(z\right)\left(k + il\right)
	\end{equation*}
	On obtient 
	\begin{equation*}	
		f'\left(z\right)\left(k + il\right) = \partialderiv{f}{x}\left(z\right)k + \partialderiv{f}{y}\left(z\right)l
	\end{equation*}
	et donc : 
	\begin{equation*}
		\begin{cases}
			f'\left(z\right) = \partialderiv{f}{x}\left(z\right)&\\
			if'\left(z\right) = \partialderiv{f}{y}\left(z\right)
		\end{cases}
	\end{equation*}
	On obtient ainsi la première égalité en identifiant.\\
	On réécrit ceci avec $\partialderiv{f}{\bar{z}} = 0$ : 
	\begin{equation*}
		\begin{cases}
			\partialderiv{P}{x} - \partialderiv{Q}{y} =& 0\\
			\partialderiv{Q}{x} + \partialderiv{P}{y} =& 0
		\end{cases}
	\end{equation*}
\end{proof}

\begin{propositionfr}{Constance sur un Connexe}{}
	Si $f$ est holomorphe sur $U$ connexe on a équivalence entre : 
	\begin{itemize}
		\item $f$ est constante sur $U$
		\item $\Re f$ l'est
		\item $\Im f$ l'est
		\item $\abs{f}$ l'est
		\item $\bar{f}$ est holomorphe
	\end{itemize}
\end{propositionfr}

\subsection{Intégrale sur le bord d'un Compact}
\begin{définition}{Classe du Bord d'un Compact}{}
	Soit $K$ un compact de $\R^{2}$. $K$ est dit à bord compact de classe $\cont^{1}$ par morceaux si pour tout élément $z_{0} \in \partial K$, il existe des coordonnées $\left(u, v\right)$ associées à un repère affine de $\R^{2}$ d'origine $z_{0}$ orienté positivement par rapport à l'orientation canonique de $\R^{2}$ et un rectangle ouvert $R = \left\{ -\delta < u < \delta\right\} \times \left\{-\eta < v < \eta\right\}$ tel que : 
	\begin{equation*}	
		K \cap R = \left\{\left(u, v\right) \in R, v \geq h\left(u\right)\right\}
		\end{equation*}
	où $h$ est une fonction réelle $\cont^{1}$ par morceaux sur $\left[-\delta, \delta \right]$ avec $h\left(0\right) = 0$ et $\sup \abs{h} < \eta$.
\end{définition}

\begin{définition}{Orientation du Bord}{}
	Soit $K$ un compact à bord de classe $\cont^{1}$ par morceaux. On appelle orientation canonique du bord l'orientation donnée par les arcs $u \mapsto \left(u, h\left(u\right)\right)$ avec $u$ croissant.
\end{définition}

\begin{lemme}{Existence de l'Orientation}{lemme:exist-orient}
	La définition a du sens.
\end{lemme}
\begin{lemme}{Recoupement de Rectangles}{}
	Soit $R, R'$ des rectangles ouverts tels que $\partial K \cap R \cap R' \neq \emptyset$.\\
	On définit 
	\begin{equation*}
		K \cap R = \left\{\left(u, v\right) \in R, v\geq h\left(u\right)\right\}
	\end{equation*}
	et 
	\begin{equation*}
		K \cap R' = \left\{\left(u', v'\right)\in R', v' \geq l\left(u'\right)\right\}
	\end{equation*}
	Alors, les orientations sur $\partial K \cap R \cap R'$ coïncident.
\end{lemme}
\begin{proof}
	Soit $z_{0}  \in \partial K \cap R \cap R'$. $h$ et $l$ sont $\cont^{1}$ par morceaux. 
	En évitant un nombre fini de points de $\partial K \cap R\cap R'$ on peut supposer $h$ et $l$ $\cont^{1}$ en $z_{0}$. 
	Autrement dit, le bord admet une tangente en $z_{0}$. On a deux repères affines orientés $\left(z_{0}, e_{1}, e_{2}\right)$ et $\left(z_{0}, e_{1}', e_{2}'\right)$ qui génèrent des coordonnées $\left(u, v\right)$ et $\left(u', v'\right)$.
	Quitte à remplacer $h$ par $h\left(u\right) - h'\left(0\right)u$ on peut supposer que $h'\left(0\right) = l'\left(0\right) = 0$. \\
	Ainsi, $e_{1}$ et $e_{1}'$ sont colinéaires. Puisqu'on a supposé que $\left(e_{1}, e_{2}\right)$ et $\left(e_{1}', e_{2}'\right)$ sont orientés positivement par rapport à l'orientation canonique de $\R^{2}$ et puisque $e_{2}$ et $e_{2}'$ doivent être dans le même sens (i.e. à l'intérieur du compact), on a bien le fait que $e_{1}$ et $e_{1}'$ sont dans le même sens.\\
	Finalement, les orientations sur $\partial K \cap R \cap R'$ coïncident. 
\end{proof}

\subsection{Formule de Green-Riemann}
Soit $p, n \in \N^{*}$. On note $\Lambda_{\R}^{p}\left(\R^{n}\right)$ le $\R$-ev des formes $p$-linéaires alternées sur $\R^{n}$. Toute forme $S \in \Lambda_{\R}^{p}\left(\R^{n}\right)$ s'écrit de manière unique 
\begin{equation*}
    S = \sum_{1 \leq i_{1} < \cdots <i_{p} \leq n}c_{i_{1},\ldots, i_{p}}\d x_{i_{1}} \land \cdots \land \d x_{i_{p}}
\end{equation*}
\begin{définition}{Produit Extérieur}{}
    Pour $S \in \Lambda_{\R}^{p}\left(\R^{n}\right), T\in \Lambda_{\R}^{q}\left(\R^{n}\right)$ on définit le produit extérieur de $S$ et $T$ noté $S \land T \in \Lambda_{\R}^{p + q}\left(\R^{n}\right)$ comme : 
    \begin{equation*}
        S\land T\left(v_{1}, \ldots, v_{p + q}\right) = \frac{1}{p!q!}\sum_{\sigma}\sgn\left(\sigma\right)S\left(v_{\sigma\left(1\right)},\ldots, v_{\sigma\left(p\right)}\right)T\left(v_{\sigma\left(p + 1\right)}, \ldots, v_{\sigma\left(p + q\right)}\right)
    \end{equation*}
\end{définition}

\begin{propositionfr}{Exemples}{}
   La paire $\left(\d x, \d y\right)$ forme une base de $\Lambda_{\R}^{1}\left(\C\right)$. Les seuls produits extérieurs à considérer sont : 
        \begin{equation*}
            \d x \land \d x = \d y \land \d y = 0, \ \ \d x \land \d y = -\d y \land \d x
        \end{equation*}
    De plus, $\d x \land \d y$ est la forme bilinéaire alternée déterminant dans la base canonique. 
\end{propositionfr}

\begin{définition}{$2$-forme différentielle}{}
    Une $2$-forme différentielle $\beta$ sur un ouvert $U$ de $\C$ est une application continue de $U$ dans $\Lambda_{\R}^{2}\left(\C\right)$ : $\beta = w\left(x, y\right)\d x \land \d y$ pour $w$ continue. 
\end{définition}

\begin{définition}{Intégrale d'une $2$-forme}{}
    Soit $\beta = w\left(x, y\right)\d x \land \d y$ une $2$-forme différentielle sur $U$. On définit : 
    \begin{equation*}
        \int_{U}\beta = \int_{U}w\left(x, y\right)\d x\d y
    \end{equation*}
\end{définition}
\begin{définition}{Différentielle d'une Différentielle}{}
    Soit $\alpha = u\left(x, y\right)\d x + v\left(x, y\right)\d y$ une $1$-forme différentielle $\cont^{1}$ sur $U$. La différentielle $\d \alpha$ de $\alpha$ est la $2$-forme différentielle 
    \begin{equation*}
        \d \alpha = \d u\land \d x + \d v \land \d y = \left(\partialderiv{v}{x} - \partialderiv{u}{y}\right)\d x \land \d y 
    \end{equation*}
\end{définition}

\begin{propositionfr}{Exemples}{}
    \begin{enumerate}
        \item Soit $\alpha$ une $1$-forme différentielle sur $U$ et $f : U\to \C$ de classe $\cont^{1}$. Alors 
        \begin{equation*}
            \d\left(f\alpha\right) = \d f \land \alpha + f\d \alpha
        \end{equation*}
        \item Si on écrit la $1$-forme différentielle $\cont^{1}$ $\alpha = f\d z + g\d \bar{z}$ on a : 
        \begin{equation*}
            \d\left(f\d z + g \d\bar{z}\right) = \left(\partial_{z}g - \partial_{\bar{z}}f\right) \d z \land \d\bar{z} = -2i\left(\partial_{z}-\partial_{\bar{z}}f\right)\d x \land \d y
        \end{equation*}
        En particulier, si $\alpha = f\d z$ avec $f$ holomorphe, alors $\partial_{\bar{z}}f = 0$ et $\d \alpha = 0$. 
    \end{enumerate}
\end{propositionfr}

\begin{lemme}{Formule de Green-Riemann sur un rectangle}{Green-Riemann}
    Soit $K$ un compact de $\C$ à bord de classe $\cont^{1}$ par morceaux orienté canoniquement. Soit $\alpha = u\left(x, y\right)\d x + v\left(x, y\right)\d y$ une forme $1$-différentielle de classe $\cont^{1}$ sur un ouvert de $K$ à support dans un rectangle $R = \left[- \delta, \delta\right] \times \left[-\eta, \eta\right] \subseteq K$. Alors 
    \begin{equation*}
        \int_{\partial K} \alpha = \int_{K}\d \alpha \text{ i.e. } \int_{\partial K}u\left(x, y\right)\d x + v\left(x, y\right) \d y = \int_{K}\left(\partialderiv{v}{x} - \partialderiv{u}{y}\right)\d x\d y
    \end{equation*}
\end{lemme}
\begin{proof}
    Comme $K$ est compact, $R\subseteq \mathring{K}$ ou bien $\partial{K \cap R}$ est le graphe d'une fonction $h$ et $K \cap R$ est la partie située à l'intérieur du graphe.\\
    Supposons $R \subseteq \mathring{K}$, alors $u = v = 0$ sur $\partial R$ et 
    \begin{equation*}
        \int_{-\delta}^{\delta}\partialderiv{v}{x}\left(x, y\right)\d x = v\left(\delta, y\right) - v\left(-\delta, y\right) = 0
    \end{equation*}
    \begin{equation*}
        \int_{-\eta}^{\eta}\partialderiv{u}{y}\left(x, y\right)\d y = u\left(x, \eta\right) - u\left(x, -\eta\right) = 0  
    \end{equation*}
    et 
    \begin{equation*}
        \int_{K}\d \alpha = \int_{R}\d \alpha = \int_{-\delta \leq x \leq \delta, -\eta \leq y \leq \eta} \left(\partialderiv{v}{x} - \partialderiv{u}{y}\right)\d x\d y = 0
    \end{equation*}
    Puisque le support de $\alpha$ ne rencontre pas $\partial K$ on a $\int_{\partial K} \alpha = 0 = \int_{K}\d \alpha$. \\
    Si $K \cap R = \{\left(x, y\right) \in R \mid y \leq h\left(x\right)\}$, alors 
    \begin{gather*}
        \int_{K}\d \alpha = \int_{K \cap R}\d \alpha = \int_{- \delta}^{\delta}\d x \int_{h\left(x\right)}^{\eta}\left(\partialderiv{v}{x} - \partialderiv{u}{y}\right)\d y\\
        = \int_{- \delta}^{\delta}\left(\frac{\partial}{\partial x}\left(\int_{h\left(x\right)}^{\eta}v\left(x, y\right)\d x\right)+ v\left(x, h\left(x\right)\right)h'\left(x\right) - \left(u\left(x, \eta\right) - u\left(x, h\left(x\right)\right)\right)\right) \d x\\
        = \int_{h\left(\delta\right), \eta}v\left(\delta, y\right)\d y - \int_{h\left(-\delta\right)}^{\eta} v\left(-\delta, y\right)\d y + \int_{-\delta}^{\delta} \left(v\left(x, h\left(x\right)\right)h'\left(x\right) + u\left(x, h\left(x\right)\right)\right)\d x\\
        = \int_{-\delta}^{\delta} \left(v\left(x, h\left(x\right)\right)h'\left(x\right) + u\left(x, h\left(x\right)\right)\right)\d x
    \end{gather*}
    car $u\left(x, \eta\right) = v\left(\delta, y\right) = v\left(-\delta, y\right) = 0$. Par ailleurs, comme $\partial K \cap R$ est paramétré par $y = h\left(x\right)$,
    \begin{equation*}
        \int_{\partial{K}\cap R}\alpha = \int_{\partial K \cap R} u\left(x, y\right)\d x + v\left(x, y\right)\d y = \int_{- \delta}^{\delta} \left(v\left(x, h\left(x\right)\right)h'\left(x\right) + u\left(x, h\left(x\right)\right)\right)\d x
    \end{equation*}
    Le support de $\alpha$ est inclus dans $R$. Nous concluons donc :
    \begin{equation*}
        \int_{K}\d \alpha = \int_{\partial K} \alpha 
    \end{equation*}
\end{proof}

\begin{définition}{Partition de l'unité}{}
    Soit $K$ un compact de $\C$ recouvert par un nombre fini d'ouverts $U_{i}$. Une partition de l'unité de classe $\cont^{1}$ subordonnée au recouvrement $U_{i}$ est une famille $\phi_{i}$ de fonctions de $K$ dans $\left[0, 1\right]$ de classe $\cont^{1}$ à support dans $U_{i}$ telles que $\sum \phi_{i}\left(x\right) = 1$ pour tout $x \in K$. 
\end{définition}

\begin{lemme}{Unité sur un Voisinage}{unit_voisin}
    Soit $z \in U$. Il existe $V$ un voisinage de $z$ avec $\overline{V}\subseteq U$ et une fonction $\cont^{1}$ $\phi_{U}$ à support dans $U$ valant $1$ sur $V$. 
\end{lemme}
\begin{proof}
    Soit $r > r'>0$ tel que $D\left(z, r'\right) \subset D\left(z, r\right) \subset U$. On définit les fonctions $\cont^{\infty}$ 
    \begin{equation*}
        f_{r} : \R \to \R, f_{r}\left(t\right) = \begin{cases}
            e^{\frac{1}{t^{2}-r^{2}}} & \text{ si } \abs{t} < r\\
            0 & \text{sinon}
        \end{cases}
    \end{equation*}
    et 
    \begin{equation*}
        g_{r} : \R \to \R, g_{r}\left(s\right) = \frac{\int_{-\infty}^{s}f_{r}\left(t\right)\d t}{\int_{- \infty}^{\infty}f_{r}\left(t\right)\d t}
    \end{equation*}
    En particulier : 
    \begin{equation*}
        g_{r}\left(s\right) = \begin{cases}
            0 & \text{si } s \leq -r\\
            1 & \text{si } s \geq r
        \end{cases}
    \end{equation*}
    Alors, $V = D\left(z, r'\right)$ et $\phi_{U}\left(w\right) = f_{r}\left(r + \frac{2r}{r-r'}\left(r'-\abs{w - z}\right)\right)$ conviennent. 
\end{proof}

\begin{lemme}{Existence d'une Partition}{}
    Soit $K \subseteq \C$ un compact et $\left(U_{i}\right)$ un recouvrement fini par des ouverts de $K$. Il existe une partition de l'unité $\cont^{1}$ subordonnée au recouvrement $U_{i}$
\end{lemme}
\begin{proof}
    Pour tout $z \in K \setminus U_{j}$ il existe $i$ tel que $z \in U_{i}$. Par le lemme \ref{lemme:unit_voisin} on constuit $\psi_{z}^{j}$ de classe $\cont^{1}$ qui vaut $1$ sur un voisinage ouvert $W_{z}^{j}$ de $z$ et dont le support est dans l'ouvert $U_{i} \cap \left(K \setminus U_{j}\right)$. Le support de $\psi_{z}^{j}$ est un fermé de $K$ donc est compact. \\
    On obtient donc un recouvrement ouvert $W_{z}^{j}$ du compact $K \setminus U_{j}$ donc on extrait un sous-recouvrement fini $\left\{W_{z_{1}}^{j}, \ldots, W_{z_{j_{l}}}^{j}\right\}$. \\
    On procède de même pour tout $j \leq n$. En réindexant on obtient une famille finie $\left(\psi_{l}\right)_{1 \leq l \leq N}$ de fonctions dont l'union des supports recouvre $K$, i.e. pour tout $z \in K$, il existe $l$ tel que $\psi_{l}\left(z\right) > 0$. On pose alors 
    \begin{equation*}
        \psi = \sum \psi_{l} \text{ et pour } l\leq N, \rho_{l} = \frac{\psi_{l}}{\psi}
    \end{equation*}
    Ainsi, $\rho_{l}$ est une partition de l'unité de classe $\cont^{1}$ de $K$ telle que pour tout $l$ il existe $1 \leq i \leq n$ tel que le support de $\rho_{l}$ soit inclus dans $U_{i}$.
\end{proof}

\begin{théorème}{Formule de Green-Riemann}{Green-Riemann}
    Soit $K$ un compact de $\C$ à bord de classe $\cont^{1}$ par morceaux orienté canoniquement. Soit $\alpha$ une $1$-forme différentielle de classe $\cont^{1}$ sur un ouvert de $K$. On a alors 
    \begin{equation*}
        \int_{\partial K}\alpha = \int_{K}\d \alpha
    \end{equation*}
\end{théorème}
\begin{proof}
    Comme $K$ est compact, il est recouvert par un nombre fini de rectangles ouverts $R_{j}$ qui vérifient $R_{j} \subseteq \mathring{K}$ ou $\partial\left(K \cap R_{j}\right)$ est le graphe d'une fonction $h_{j}$ et $K\cap R_{j}$ est la partie située à l'intérieur du graphe. Soit $\left(\chi_{j}\right)$ une partition de l'unité subordonnée au recouvrement $R_{j}$. Écrivons $\alpha = \sum \alpha_{j}$ où les $1$-formes différentielles $\alpha_{j} = \chi_{j}\alpha$ sont de classes $\cont^{1}$ à support dans $R_{j}$. On se ramène alors au cas du lemme \ref{lemme:Green-Riemann}
\end{proof}

\begin{théorème}{Cauchy}{}
    Soit $U$ un ouvert de $\C$, $K$ un compact à bord de classe $\cont^{1}$ par morceaux inclus dans $U$, avec l'orientation canonique du bord. Alors pour toute fonction holomorphe de classe $\cont^{1}$ sur $K$ nous avons 
    \begin{equation*}
        \int_{\partial K}f\left(z\right)\d z = 0
    \end{equation*}
\end{théorème}
\begin{proof}
    On applique la formule de Green-Riemann \ref{thm:Green-Riemann} à $\alpha = f\left(z\right)\d z$, $1$-forme différentielle de classe $\cont^{1}$. On a $\d \alpha = -\partial_{\bar{z}}f\d z \land \d \bar{z} = 0$.
\end{proof}
\begin{corollaire}{Analycité Holomorphe $\cont^{1}$}{}
    Soit $f$ holomorphe de classe $\cont^{1}$ sur un ouvert $U$. Alors $f$ est analytique sur $U$. 
\end{corollaire}
\begin{proof}
    Soit $\overline{D}\left(w, r\right) \subseteq U$ et $\gamma$ le lacet $t \mapsto w + re^{it}$. Pour $\lambda \leq 1$, on pose
    \begin{equation*}
        g\left(\lambda\right) = \frac{1}{2i\pi}\int_{\gamma}\frac{f\left(z + \lambda\left(u - z\right)\right)}{u - z} \d u = \frac{r}{2\pi}\int_{0}^{2\pi}\frac{f\left(z +\lambda\left(w + re^{it} - z\right)\right)}{w + re^{it} - z}e^{it}\d t
    \end{equation*}
    Ainsi, $g$ est continue sur $\left[0, 1\right]$, dérivable sur $\left]0, 1\right[$ de dérivée 
    \begin{equation*}
        g'\left(\lambda\right) = \frac{r}{2\pi}\int_{0}^{2\pi}f'\left(z + \lambda\left(w + re^{it} - z\right)\right)e^{it}\d t = \left[\frac{1}{2i\pi\lambda}f\left(z + \lambda\left(w + re^{it} - z\right)\right)\right]_{t = 0}^{2\pi} = 0
    \end{equation*}
    Donc $g$ est constante avec 
    \begin{equation*}
        g\left(1\right) = \frac{1}{2i\pi}\int_{\gamma}^{}\frac{f\left(u\right)\d u}{u - z} \text{ et } g\left(0\right) = \frac{1}{2i\pi}\int_{\gamma}\frac{f\left(z\right)\d u}{u - z} = f\left(z\right)
    \end{equation*}
    D'où, par la proposition \ref{propo:cont_lacet}, $f$ est analytique 
\end{proof}
\subsection{Analycité des Fonctions Holomorphes}
\begin{lemme}{Goursat}{Goursat}
    Soient $U \subseteq \C$ un ouvert et $T$ un triangle inclus dans $U$. Pour tout fonction holomorphe sur $U$
    \begin{equation*}
        \int_{\partial T}f\left(z\right)\d z = 0
    \end{equation*}
\end{lemme}
\begin{proof}
    Nous découpons $T$ en quatre triangles $T_{i}$ dont les sommets sont ceux de $T$ et les milieux des côtés de $T$. Nous orientons les arêtes opposées des triangles $T_{k}$ de telle façon que 
    \begin{equation*}
        I = \int_{\partial T}f\left(z\right)\d z = \sum_{k = 1}^{4}\int_{\partial T_{k}}f\left(z\right)\d z
    \end{equation*}
    Il existe donc un indice $k$ avec $\abs{\int_{\partial T_{k}}f\left(z\right)\d z} \geq \abs{I}/4$. De cette façon, nous construisons une suite de triangles emboîtés $T_{0}' = T, T_{1}' = T_{k}$ avec $diam T_{n}' = diam T/2^{n}$ et $\abs{\int_{\partial T_{n}'} f\left(z\right)\d z} \geq \abs{I}/4^{n}$.
	L'intersection des triangles emboîtés $T_{n}'$ est réduite à un point $z_{0}$. Comme $f$ est holomorphe en $z_{0}$ : 
	\begin{equation*}
		f\left(z\right) = f\left(z_{0}\right) + \left(z - z_{0}\right)f'\left(z_{0}\right) + \left(z - z_{0}\right)\epsilon\left(z\right)
	\end{equation*}
	avec $\epsilon\left(z\right)$ qui tend vers $0$ quand $z$ tend vers $z_{0}$.  
	On a ainsi :
	\begin{equation*}
		\abs{\int_{\partial T_{n}'} f\left(z\right)\d z} = \abs{\int_{\partial T_{n}'} \left(z - z_{0}\right) \epsilon\left(z\right) \d z} \leq {\rm long}\left(\partial T_{n}'\right) \sup_{\partial T_{n}'} \abs{z - z_{0}}\abs{\epsilon\left(z\right)}
	\end{equation*}
	et donc 
	\begin{equation*}
		\abs{\int_{\partial T_{n}'} f\left(z\right)\d z}\leq 3 \left({\rm diam}T_{n}'\right)^{2}\sup_{\partial T_{n}'}\abs{\epsilon\left(z\right)}
	\end{equation*}
	Donc $\abs{I} \leq 4^{n}\abs{\int_{\partial T_{n}'} f\left(z\right) \d z} \leq 3 \left({\rm diam}T_{n}\right)^{2}\sup_{\partial T_{n}'}\abs{\epsilon\left(z\right)}$ et donc $I = 0$.
\end{proof}

\begin{théorème}{Goursat}{Goursat}
	Soit $U \subseteq \C$ ouvert et $K$ un compact à bord de classe $\cont^{1}$ avec l'orientation canonique du bord. Pour toute fonction holomorphe sur $U$ on a : 
	\begin{equation*}
		\int_{\partial K}f\left(z\right)\d z = 0
	\end{equation*}
\end{théorème}
\begin{proof}
    On approche $K$ par des compacts à bords polygonaux. Notons $\delta = d\left(K, \C\setminus U\right) > 0$. 
    Paramétrons $\delta K$ par un nombre fini d'arcs $\cont^{1}$ par morceaux. Pour chaque tel arc $\gamma : \left[a, b\right] \to U$, soit une subdivision $a = \tau_{0} < \tau_{1} < \ldots < \tau_{n} = b$ telle que $\abs{\gamma\left(\tau_{j + 1}\right) - \gamma\left(\tau_{j}\right)} \leq \epsilon \leq \delta/2$. 
    Chaque segment $\left[\gamma\left(\tau_{j + 1}\right), \gamma\left(\tau_{j}\right)\right] \subset U$. 
    Pour $\epsilon$ assez petit, la réunion de ces segments constitue le bord d'un compact $K_{\epsilon}$ à bord polygonal. 
    $K_{\epsilon} = \cup_{i} T_{i}$ est réunion de triangles adjacents et le lemme de Goursat \ref{lemme:Goursat} implique 
    \begin{equation*}
        \int_{\partial K_{\epsilon}}f\left(z\right)\d\left(z\right) = \sum_{i}\int_{\partial T_{i}} f\left(z\right) \d z = 0
    \end{equation*}
    D'après la proposition , on a bien : 
    \begin{equation*}
        \lim_{\epsilon \to 0}\int_{\partial K_{\epsilon}} = \int_{\partial K}f\left(z\right)\d z
    \end{equation*}
    D'où le résultat. 
\end{proof}

\begin{théorème}{Formule de Cauchy}{formule_cauchy}
    Soit $f$ holomorphe sur un ouvert $U \subseteq \C$ et $K$ un compact à bord orienté $\cont^{1}$ par morceaux inclus dans $U$. Alors, pour tout $z \in K$ 
    \begin{equation*}
        f\left(z\right) = \frac{1}{2i\pi}\int_{\partial K}\frac{f\left(\omega\right)}{\omega -  z}\d \omega
    \end{equation*}
\end{théorème}

\begin{proof}
    Soit $r > 0$ tel que $\overline{D\left(z, r\right)} \subset \mathring{K}$. On note $K_{r} = K \setminus D\left(z, r\right)$. $K_{r}$ est un compact à bord orienté $\cont^{1}$ par morceaux dont le bord est $\partial K_{r} = \partial K \cup \partial D^{-}\left(z, r\right)$ où $\partial D^{-}$ signifie que ce cercle a l'orientation opposée à celle obtenue comme bord de $\overline{D\left(z,r\right)}$. La fonction $g\left(\omega\right) = f\left(\omega\right)/\left(\omega - z\right)$ et holomorphe sur $U \setminus \{z\}$. Le théorème de Goursat \ref{thm:Goursat} appliqué à $g$ sur le compact $K_{r} \subseteq U \setminus \{z\}$ donne 
    \begin{equation*}
        \int_{\partial K}\frac{f\left(\omega\right)}{\omega - z}\d \omega - \int_{\partial D\left(z, r\right)} \frac{f\left(\omega\right)}{\omega - z}\d \omega = 0
    \end{equation*}
    En posant $\omega = z + re^{it}$ on a : 
    \begin{equation*}
        \int_{\partial D\left(z, r\right)}\frac{f\left(\omega\right)}{\omega - z}\d \omega = \int_{0}^{2\pi}\frac{f\left(z + re^{it}\right)}{re^{it}}ire^{it}\d t = i \int_{0}^{2\pi}f\left(z + re^{it}\right) \d t
    \end{equation*}
    et cette dernière intégrale tend vers $2i\pi f\left(z\right)$ lorsque $r$ tend vers $0$ par continuité de $f$ au point $z$. 
\end{proof}

\begin{théorème}{Équivalence Holomorphie-Analycité}{equiv_holo_analy}
    Soit $f : U \to \C$. $f$ est holomorphe sur $U$ si et seulement si elle est analytique.
\end{théorème}
\begin{proof}
    On a déjà l'implication analycité holomorphie. Supposons $f$ holomorphe sur $U$ et $\overline{D}\left(z_{0}, r\right) \subset U$. Pour $z \in D\left(z_{0}, r\right)$, la formule de Cauchy \ref{thm:formule_cauchy} donne 
    \begin{equation*}
        f\left(z\right) \frac{1}{2i\pi}\int_{\partial D\left(z_{0}, r\right)}\frac{f\left(\omega\right)}{\omega - z}\d \omega
    \end{equation*}
    Or 
    \begin{equation*}
        \frac{1}{\omega - z} = \sum_{n = 0}^{+ \infty}\frac{\left(z - z_{0}\right)^{n}}{\left(\omega - z_{0}\right)^{n + 1}}
    \end{equation*}
    De plus, pour $\omega = z_{0} + re^{it}$,
    \begin{equation*}
        \abs{\frac{\left(z - z_{0}\right)^{n}}{\left(\omega - z_{0}\right)^{n + 1}}} = \frac{1}{r}\left(\frac{\abs{z - z_{0}}}{r}\right)^{n}, \text{ avec } \abs{z - z_{0}}/r < 1
    \end{equation*}
    Par convergence normale pour $t \in \left[0, 2\pi\right]$, on obtient :
    \begin{equation*}
        \begin{aligned}
            f\left(z\right) &= \frac{1}{2i\pi}\int_{\partial D\left(z_{0},r \right)}f\left(\omega\right)\sum_{n = 0}^{+ \infty}\frac{\left(z - z_{0}\right)^{n}}{\left(\omega - z_{0}\right)^{n + 1}}\d \omega \\ &= \sum_{n = 0}^{+ \infty}\int_{\partial D\left(z_{0}, r\right)}\frac{f\left(\omega\right)}{2i\pi\left(\omega - z_{0}\right)^{n + 1}}\d \omega\left(z - z_{0}\right)^{n} \\ &= \sum a_{n}\left(z - z_{0}\right)^{n}
        \end{aligned}
    \end{equation*}
    et la série entière ci-dessus converge normalement sur les compacts de $D\left(z_{0}, r\right)$.
\end{proof}
\begin{corollaire}{Classe des Dérivées}{}
    Soit $U$ un ouvert de $\C$. Toute fonction holomorphe sur $U$ est de classe $\cont^{\infty}$ sur $U$.\\
    Précisément, pour tout $K \subset U$ compact à bord de classe $\cont^{1}$ par morceaux et pour tout $z \in \mathring{K}$ nous avons : 
    \begin{enumerate}
        \item $\forall n \geq 0, \frac{\partial^{n}f}{\partial z^{n}}\left(z\right) = f^{\left(n\right)}\left(z\right) = \frac{n!}{2i\pi}\int_{\partial K}\frac{f\left(\omega\right)}{\left(\omega - z\right)^{n + 1}}\d \omega$
        \item $\forall n \geq 0, \forall m \geq 0, \frac{\partial^{n + m}f}{\partial z^{n}\partial \bar{z}^{m}}\left(z\right) = 0$.
    \end{enumerate}
    En particulier, une fonction holomorphe $f$ admet des dérivées complexes $f^{\left(n\right)}$ d'ordre $n$ arbitraire et les dérivées $f^{\left(n\right)}$ sont holomorphes.
\end{corollaire}

\begin{théorème}{Morera}{morera}
    Soit $f$ une fonction continue sur un ouvert $U$ de $\C$. Nous supposons que $\int_{\partial T}f\left(z\right) \d z = 0$ pour tout triangle $T$ inclus dans $U$. Alors $f$ est holomorphe sur $U$.
\end{théorème}
\begin{proof}
    Soit $z_{0} \in U$ et $r > 0$ tel que $\overline{D}\left(z_{0}, r\right) \subset U$. Pour $z \in D\left(z_{0}, r\right)$, on pose 
    \begin{equation*}
        F\left(z\right) = \int_{\left[z_{0}, z\right]} f\left(\omega\right) \d \omega
    \end{equation*}
    Soit $z \in D\left(z_{0}, r\right)$ et $h \neq 0$ tel que $z + h \in D\left(z_{0}, r\right)$. Comme le triangle de sommets $z_{0}, z, z + h$ est inclus dans $D\left(z_{0}, r\right)$, nous avons 
    \begin{equation*}
        \frac{F\left(z + h\right)-F\left(z\right)}{h} = \frac{1}{h}\int_{\left[z, z + h\right]}f\left(\omega\right)\d \omega = \int_{0}^{1} f\left(z + th\right)\d t
    \end{equation*}
    Comme $f$ est continue au point $z$, 
    \begin{equation*}
        \lim_{h \in \C^{*}, h \to 0} \frac{F\left(z + h\right) - F\left(z\right)}{h} = f\left(z\right)
    \end{equation*}
    Ainsi $F$ est holomorphe sur $D\left(z_{0}, r\right)$ donc analytique d'après le théorème \ref{thm:equiv_holo_analy} et sa dérivée $f = F'$ l'est donc aussi.
\end{proof}

\begin{corollaire}{$\Gamma$}{gamma}
    La fonction $\Gamma$
    \begin{equation*}
        \Gamma\left(s\right) = \int_{0}^{+\infty}e^{-t}t^{s - 1}\d t
    \end{equation*}
    est holomorphe pour $\Re s > 0$.
\end{corollaire}

\begin{proof}
    L'intégrale converge en $t = 0$ car $\abs{t^{s - 1}e^{-t}} \leq t^{\Re s  - 1}$. À $s$ fixé pour $t \in \R_{+}$ grand, 
    \begin{equation*}
        \abs{t^{s - 1}e^{-t}} = t^{\Re s - 1}e^{-t}\leq e^{t/2}e^{-t} = e^{-t/2}
    \end{equation*}
    Donc $\Gamma\left(s\right)$ est bien définie pour $\Re s > 0$. Soit $\gamma : \left[0, 1\right] \to \left\{s, \Re s > 0\right\}$ la courbe décrivant un triangle. Alors, d'après le théorème de Fubini 
    \begin{equation*}
        \int_{\gamma}\Gamma\left(s\right) \d s = \int_{\gamma}\int_{0}^{+\infty}t^{s - 1}e^{-t}\d t\d s = \int_{0}^{+ \infty}\left(\int_{\gamma}t^{s-1}\d s\right)e^{-t}\d t = 0
    \end{equation*}
    Ainsi, en appliquant le théorème de Morera \ref{thm:morera}, la fonction $\Gamma$ est holomorphe sur le demi-plan $\Re s > 0$.
\end{proof}

\section{Propriétés Éléméntaires des Fonctions Holomorphes}
\subsection{Théorème d'inversion locale}
\begin{théorème}{Inversion Locale}{inv_loc}
    Si $f\in \O\left(U\right)$, $a \in U, f'\left(a\right) \neq 0$, alors, $\exists V$ voisinage ouvert de $a$ inclus dans $U$ sur lequel $f$ est biholomorphe sur $f\left(V\right)$ ouvert. 
\end{théorème}
\begin{proof}
    Comme $f\in \O\left(U\right)$, $f$ est $\R$-différentiable. Donc il existe un voisinage $V$ ouvert de $U$ contenant $a$ sur lequel $f_{\mid V} : V \to f\left(V\right)$ est un difféomorphisme. Alors, $\d_{f\left(z\right)}\left(f^{-1}\right) = \left(d_{z} f\right)^{-1}$ et donc $f^{-1} \in \O\left(U\right)$.
\end{proof}
\begin{proof}[Idée des Séries Majorantes]\label{thm:inv_loc_serie_maj}
    \begin{itemize}
        \item On suppose d'abord $a = 0, f\left(a\right) = 0, f'\left(a\right) = 1$. On a 
        \begin{equation*}
            f\left(z\right) = z - \sum_{n \geq 2}a_{n}z^{n}, z\in D\left(0, r\right)
        \end{equation*}
        On veut résoudre $f\left(z\right) = \omega = z - \sum_{n \geq 2} a_{n}z^{n}$ i.e. $z = \omega + \sum_{n \geq 2}a_{n}z^{n}$. Mais, $\sum_{n \geq 2}a_{n}z^{n} = \O\left(w^{2}\right)$ : 
        \begin{equation*}
            z = \omega + \sum_{n \geq 2}a_{n}\left(\omega + \O\left(\omega^{2}\right)\right)^{n} = \omega + a_{2}\omega^{2} + \O\left(\omega^{3}\right)
        \end{equation*}
        On peut alors réinjecter : 
        \begin{equation*}
            z = \omega + a_{2}\omega^{2} + \left(2a_{2}^{2} + a_{3}\right)\omega^{3} + \O\left(\omega^{4}\right)
        \end{equation*}
        et ainsi de suite : 
        \begin{equation*}
            z = \omega + \sum_{n = 2}^{N}P_{n}\left(a_{2}, \ldots, a_{n}\right)\omega^{n} + \O\left(\omega^{N+1}\right)
        \end{equation*}
        où les $P_{n} \in \N\left[X_{2}, \ldots, X_{n}\right]$.
        \item Montrons maintenant que cette série converge lorsque $N \to \infty$. On sait que la série $\sum a_{n}z^{n}$ converge sur $D\left(0, r\right)$. Pour $r' < r$, $\abs{a_{n}r'^{n}} \to 0$. Donc il existe $M > 0$ tel que $\abs{a_{n}}\leq M^{n}$. Or, \begin{equation*}z = \omega + \sum_{n = 2}^{+ \infty}P_{n}\left(M^{2}, \ldots, M^{n}\right)\omega^{n}\end{equation*} est solution de :
        \begin{equation*}
            \begin{aligned}
                \omega =& z - \sum_{n \geq 2}M^{n}z^{n}\\
                =& z - \left(\frac{1}{1 - Mz} - 1 - Mz\right)
            \end{aligned}
        \end{equation*}
        Donc 
        \begin{equation*}
            \begin{aligned}
                \left(1 - Mz\right)\omega =& z\left(1 - Mz\right) - 1 + 1 - Mz + Mz\left(1 - Mz\right)
            \end{aligned}
        \end{equation*}
        C'est à dire : 
        \begin{equation*}
            z^{2}\left(M + M^{2}\right) + z\left(-M\omega - 1\right) + \omega = 0
        \end{equation*}
        ou 
        \begin{equation*}
            z = \frac{\left(M\omega + 1\right) - \sqrt{\left(1 + M\omega\right)^{2} - 4\omega\left(M + M^{2}\right)}}{2\left(M + M^{2}\right)}
        \end{equation*}
        On prend ici pour $\sqrt{\cdot}$ la détermination holomorphe de $\left(\right)^{1/2}$ qui existe sur $D\left(1, 1\right)$ et pour laquelle $\sqrt{1} = 1$ de sorte que pour $\omega = 0$, $z = 0$.\\
        La série définissant $\sqrt{\cdot}$ converge alors sur $D\left(0, R\right)$ où $R = \frac{1}{\left(1 + \sqrt{2}\right)M + 4M^{2}}$. En effet, alors, on a 
        \begin{equation*}
            \abs{M^{2}\omega^{2}} \leq M^{2}\abs{\omega}R \leq \frac{M^{2}\abs{\omega}}{\left(1 + \sqrt{2}\right)M} = \left(\sqrt{2}- 1\right)M\abs{\omega}
        \end{equation*}
        et donc 
        \begin{equation*}
            \abs{\left(2M + 4M^{2}\right)\omega - M^{2}\omega^{2}} \leq \left(2M + 4M^{2}\right)\abs{\omega} + \abs{M^{2}\omega^{2}}\leq \left(\left(1 + \sqrt{2}\right)M + 4M^{2}\right)\abs{w} < 1
        \end{equation*}
        D'où la convergence de $g\left(\omega\right) = \omega + \sum_{n \geq 2} P_{n}\left(a_{2}, \ldots, a_{n}\right)\omega^{n}$ sur $D\left(0, R\right)$. et $g\left(D\left(0, R\right)\right) \subset D\left(0, 1/M\right)$.
        \item Par identification de la série entière en zéro et principe du prolongement analytique, nous avons $f \circ g\left(\omega\right) = \omega$ pour $\omega \in D\left(0, R\right)$. De plus, par construction, $g$ est injective sur $W = D\left(0, R\right)$ et l'image $\omega = f\left(z\right)$ atteint surjectivement $W$ sur $g\left(W\right) \subseteq D\left(0, 1/M\right) \cap f^{-1}\left(W\right)$. Prenons $V$ la composante connexe de $0$ dans $D\left(0, 1/M\right) \cap f^{-1}\left(W\right)$. Alors $f\left(V\right) \subset W$ et $g\left(W\right) \subset V$. $V, W$ sont ouverts et $f_{\mid V} \circ g_{\mid W} = id_{W}$. Par connexité de $V$ et prolongement analytique, $g_{\mid W} \circ f_{\mid V} = id_{V}$.
    \end{itemize}
\end{proof}
\subsection{Théorème de l'Application Ouverte}
\begin{théorème}{Pré-Application Ouverte}{pre_open_app}
    Soit $f \in \O\left(U\right)$ non constante au voisinage de $a \in U$, $f\left(a\right) = 0$ et \begin{equation*}m = \min\{k \in \N^{*}\mid f^{\left(k\right)}\left(a\right) \neq 0\}\end{equation*}
    Il existe alors un voisinage ouvert $V$ de $a$, un voisinage ouvert $W$ de $0$ et un biholomorphisme $\phi : V\to W$ tel que $\phi$ envoie $a$ sur $0$ et $f\left(z\right) = f\left(a\right) + \phi\left(z\right)^{m}$.
\end{théorème}
\begin{proof}
    D'après le théorème \ref{thm:multiplicité} il existe $U'\subseteq U$ un voisinage de $a$ et $g\in \O\left(U'\right)$ tels que pour tout $z \in U'$
    \begin{equation*}
        f\left(z\right) - f\left(a\right) = \alpha\left(z - a\right)^{m}g\left(z\right)
    \end{equation*}
    avec $\alpha \in \C^{*}$ et $g\left(a\right) = 1$.\\
    Soit $ V = \left\{z \in U'\mid \abs{g\left(z\right) - 1} < 1\right\}$. C'est un voisinage de $a$ sur lequel $\exp \frac{1}{m}\log\left(g\left(z\right)\right)$ existe.\\
    On a alors
    \begin{equation*}
        \forall z \in V', f\left(z\right) = f\left(a\right) + \left(\phi\left(z\right)\right)^{m}
    \end{equation*}
    où 
    \begin{equation*}
        \phi\left(z\right) = \alpha_{m}\left(z - a\right)\exp\left(\frac{1}{m}\log\left(g\left(z\right)\right)\right)
    \end{equation*}
    où $\alpha_{m}^{m} = \alpha$. Alors, $\phi \in \O\left(V'\right)$ avec $\phi\left(a\right) = 0$ et $\phi'\left(a\right) = 1$. Par théorème d'inversion locale \ref{thm:inv_loc}, on a un voisinage $V \subset V'$ de $a$ sur lequel $\phi$ est un biholomorphisme.
\end{proof}

\begin{corollaire}{Solutions d'une Équation}{nb_sol_eq}
    Soit $f\in \O\left(U\right)$ non constante au voisinage de $a \in U$ et \begin{equation*}m = \min\{k \in \N^{*}\mid f^{\left(k\right)}\left(a\right) \neq 0\}\end{equation*}. Alors, $\exists r,\rho \in \R_{+}^{*}$ tels que $\forall \omega \in D\left(f\left(a\right), \rho\right) \setminus \{f\left(a\right)\}$ l'équation $f\left(z\right) = \omega$ a exactement $m$ solutions dans $D\left(a, r\right)$.
\end{corollaire}
\begin{proof}
    On écrit par le théorème \ref{thm:pre_open_app} précédent $f\left(z\right) = \omega = f\left(a\right) + \phi\left(z\right)^{m}$ où $\phi : V \to W$ est tel que $\phi\left(a\right) = 0$. On suppose $\phi\left(z\right) = \left(\omega - f\left(a\right)\right)^{1/m}$ pour une certaine détermination de l'exponentielle. 
    On prend $r$ tel que $D\left(a, r\right) \subset V$. $\phi\left(D\left(a, r\right)\right)$ est un ouvert de $W$ voisinage de $0$. Il existe un $\rho'$ tel que $D\left(0, \rho'\right)$ est inclus dans $\phi\left(D\left(a, r\right)\right)$. 
    Alors, pour tout $\omega \in D\left(f\left(a\right), \rho'^{m}\right)$, $\left(\omega - f\left(a\right)\right)^{1/m} \in D\left(0, \rho'\right)$. Mézalor, $e^{2ik\pi/m}\left(w -f\left(a\right)\right)^{1/m}$ sont dans $D\left(0, \rho'\right)$.
    On obtient alors \begin{equation*}z_{k} = \phi^{-1}\left(e^{2ik\pi/m}\left(\omega - f\left(a\right)\right)^{1/m}\right) \in D\left(a, r\right)\end{equation*}.\\
    Les $z_{k}$ sont solutions de $f\left(z\right) = \omega$ et donc il y en a bien exactement $m$. \\
    De même, l'équation $f\left(z\right) = f\left(a\right)$ n'a qu'une solution $z = a$ dans $D\left(a, r\right)$ de multiplicité $m$.
\end{proof}

\begin{théorème}{Application Ouverte}{open_app}
    Une fonction holomorphe non constante sur un ouvert $U$ connexe est une application ouverte.
\end{théorème}
\begin{proof}
    Par le corollaire \ref{cor:nb_sol_eq}, tout point $z_{0} \in U$ admet un voisinage $V_{z_{0}} \subset U$ tel que $f\left(V_{z_{0}}\right) = D\left(f\left(z_{0}\right), \rho\left(z_{0}\right)\right)$. Ainsi, $f\left(U\right) = \cup D\left(f\left(z_{0}\right), \rho\left(z_{0}\right)\right)$ est ouvert. 
\end{proof}
\begin{théorème}{Théorème d'Inversion Gloable}{global_inv}
    Soit $U$ un ouvert connexe et $f \in \O\left(U\right)$ injective. Alors : 
    \begin{enumerate}
        \item $f\left(U\right)$ est un ouvert de $\C$
        \item $f'$ ne s'annule pas sur $U$
        \item $f : U \to f\left(U\right)$ est un biholomorphisme
    \end{enumerate}
\end{théorème}
\begin{proof}
    \begin{enumerate}
        \item D'après le théorème de l'application ouverte \ref{thm:open_app}, $f$, injective donc non constante, est ouverte donc $f\left(U\right)$ est ouverte et $f$ est une bijection continue ouverte de $U$ dans $f\left(U\right)$, i.e., un homéomorphisme. 
        \item Supposons qu'il existe $z_{0}$ pour lequel $f'\left(z_{0}\right) = 0$. Dans le théorème \ref{cor:nb_sol_eq}, on a un entier $m \geq 2$ et donc $f$ n'est pas injective au voisinage de $z_{0}$ ce qui est absurde. Donc $f'$ ne s'annule pas sur $U$. 
        \item D'après les deux premiers points et le théorème \ref{thm:inv_loc} d'inversion locale, $f^{-1}$ est holomorphe sur $f\left(U\right)$ et $f: U \to f\left(U\right)$ est un biholomorphisme.
    \end{enumerate}
\end{proof}

\subsection{Lemme de Schwarz}
\begin{théorème}{Principe du Maximum}{principe_max}
    Soit $f \in \O\left(U\right)$
    \begin{enumerate}
        \item Si $\abs{f}$ admet une maximum local en un point $a \in U$, alors $f$ est constante sur la composante connexe contenant $a$. 
        \item Pour tout $K \subset U$
        \begin{equation*}
            \max_{K}\abs{f} = \max_{\partial K}\abs{f}, \max_{K}\Re f = \max_{\partial K}\Re f, \max_{K}\Im f = \max_{\partial K} \Im f
        \end{equation*}
    \end{enumerate}
\end{théorème}
\begin{proof}
    \begin{enumerate}
        \item Supposons $f$ non constante sur la composante connexe $U_{0}$ de $U$ contenant $a$ avec $\abs{f\left(a\right)} = \sup_{U}\abs{f} = \sup_{U_{0}}\abs{f}$. D'après le théorème de l'application ouverte, $f$ est ouverte sur $U_{0}$. L'image $f\left(U_{0}\right)$ est un voisinage de $f\left(a\right)$ donc contient des points de module strictement supérieur à $\abs{f\left(a\right)}$.
        \item Si $\max_{\partial K}\Re f < \max_{K}\Re f$, il existe $z_{0} \in \mathring{K}$ avec $\Re f\left(z_{0}\right) = \max_{K}\Re f$. Soit $U_{0}$ une composante connexe de $z_{0}$ dans $\mathring{K}$ et $f$ non constante dans $U_{0}$. Alors $f\left(U_{0}\right)$ est un ouvert qui contient $f\left(z_{0}\right)$ et qui est contenue dans le demi-plan $\left\{w\mid \Re w \leq \Re f\left(z_{0}\right)\right\}$. Donc $f$ est constante sur $U_{0}$ et $\Re f_{\mid \partial U_{0}} = \Re f\left(z_{0}\right)$ par continuité de $f$. Or
        \begin{equation*}
            \emptyset \neq \partial U_{0} \subset \partial \mathring{K} = \overline{\mathring{K}} \setminus \mathring{K} \subseteq \overline{K} \setminus \mathring{K} = \partial K
        \end{equation*}
        ainsi, $\max_{K}\Re f = \Re f\left(z_{0}\right)$ est atteint sur $\partial K$. Les cas $\abs{f}$ et $\Im f$ sont analogues. 
    \end{enumerate}
\end{proof}

\begin{théorème}{Lemme de Schwarz}{lemme_schwarz}
    Soit $f$ holomorphe sur $D\left(0, 1\right)$ avec $f\left(0\right) = 0$ et $\abs{f\left(z\right)} \leq 1$ sur $D\left(0, 1\right)$. Alors, on a 
    \begin{equation*}
        \forall z \in D\left(0, 1\right), \abs{f\left(z\right)} \leq \abs{z} \text{ et } \abs{f'\left(0\right)}\leq 1
    \end{equation*}
    De plus si 
    \begin{equation*}
        \exists z_{0}\in D\left(0, 1\right)\setminus \{0\}, \abs{f\left(z_{0}\right)} = \abs{z_{0}} \text{ ou } \abs{f'\left(0\right)} = 1
    \end{equation*}
    alors $f$ est une rotation : $f\left(z\right) = az$ pour $a \in \C, \abs{a} = 1$.
\end{théorème}
\begin{proof}
    L'application $g\left(z\right) = f\left(z\right)/z$ est holomorphe sur $D\left(0, 1\right)$ et vérifie $g\left(0\right) = f'\left(0\right)$. En appliquant le principe du maximum à la fonction $g$ sur $D\left(0, r\right), r < 1$, on obtient : 
    \begin{equation*}
        \sup_{D\left(0, r\right)}\abs{g} = \sup_{\partial D\left(0, r\right)}\abs{g} = \sup_{\partial D\left(0, r\right)}\abs{f}/r \leq 1/r
    \end{equation*}
    En faisant tendre $r$ vers $1$, nous obtenons
    \begin{equation*}
        \sup_{D\left(0, 1\right)}\abs{g} \leq 1 \Longleftrightarrow \begin{cases}
            \abs{f\left(z\right)} \leq \abs{z} & \forall z \in D\left(0, 1\right)\setminus \{0\}\\
            \abs{f'\left(0\right)} \leq 1 &
        \end{cases}
    \end{equation*}
    De plus, si $\exists z_{0} \in D\left(0, 1\right)\setminus \{0\}$, $\abs{f\left(z_{0}\right)} = \abs{z_{0}}$ ou si $\abs{g\left(0\right)} = \abs{f'\left(0\right)} = 1$, alors $\abs{g}$ admet un maximum local en $z_{0}$ ou $0$ donc $g = a$ est constante avec $\abs{a} = 1$.
\end{proof}

\begin{corollaire}{Point Fixe}{point_fixe}
    Soit $f : D\left(0, 1\right) \to D\left(0, 1\right)$ holomorphe, $f\neq {\rm Id}$. Alors $f$ a au plus un point fixe. 
\end{corollaire}
\begin{proof}
    Supposons que nous ayons $a \neq b \in D\left(0, 1\right)$ avec $f\left(a\right) = a$ et $f\left(b\right) = b$.\\
    Si $a = 0, f\left(0\right) = 0, f\left(b\right) = b$, le lemme de Schwarz montre que $f = e^{i\theta}z$ et comme $f\left(b\right) = b$, $f = \rm Id$.\\
    On peut donc supposer $a \neq 0$. Posons $\phi_{\alpha}\left(z\right)\footnote{C'est la fonction du jour !} = \frac{z - a}{1 - \bar{a}z}$ et 
    \begin{equation*}
        g = \phi_{\alpha} \circ f \circ \phi_{\alpha}^{-1} \text{ et } \lambda = \phi_{\alpha}\left(b\right) \neq 0
    \end{equation*}
    Ainsi, $g : D\left(0, 1\right) \to D\left(0, 1\right)$ avec deux points fixes $g\left(0\right) = 0, g\left(\lambda\right) = \lambda$. Donc $g = {\rm Id} = f$. Absurde.
\end{proof}

\subsection{Disque Unité et Inversion Locale Effective}
\begin{définition}{Automorphisme}{}
    Un automorphisme de $U$ est une bijection holomorphe de $U$ dans $U$. On note $\Aut U$ l'ensemble des automorphismes de $U$. 
\end{définition}

On rappelle que le groupe unitaire de signature $\left(n, m\right)$ est le groupe des matrices préservant une forme hermitienne de signature $\left(n, m\right)$. Il est isomorphe au groupe $U\left(n, m\right)$ préservant la forme hermitienne diagonale de signature $\left(n, m\right)$ : 
\begin{equation*}
    U\left(n, m\right) = \left\{\begin{pmatrix}
        A & B \\ C & D
    \end{pmatrix} \in {\rm GL}_{n + m}\left(\C\right) \ \middle| \ \begin{aligned}&AA* - BB* = I_{n}\\
    &DD* - CC* = I_{m}\\ &AC* = BD* \end{aligned} \right\} 
\end{equation*}

\begin{théorème}{Automorphisme du Disque Unité}{}
    \begin{equation*}
        \Aut D\left(0, 1\right) = \left\{\phi_{\alpha} : z \longmapsto e^{i\theta}\frac{z - a}{1 - \bar{a}z}\ \middle|\ \theta \in \R, a \in \C, \abs{a} < 1\right\} = {\rm PSU}\left(1, 1\right)
    \end{equation*}
\end{théorème}
\begin{proof}
    \begin{itemize}
        \item Nous vérifions que $\phi_{\alpha} \in \Aut D$. Soit $f \in \Aut D$. On note $a = f\left(0\right)$. On a $\phi_{\alpha} \circ f\left(0\right) = 0$ et pour tout $z \in D$, $\abs{\phi_{\alpha} \circ f\left(z\right)} < 1$. Le lemme de Schwarz \ref{thm:lemme_schwarz} implique alors que $g = \phi_{\alpha} \circ f$ vérifie $\abs{g\left(z\right)} \leq \abs{z}$ et de même $\abs{g^{-1}\left(z\right)} \leq \abs{z}$ donc $\abs{g\left(z\right)} = \abs{z}$ pour $z\in D$. Donc il existe $\theta \in \R$ tel que $g\left(z\right) = e^{i\theta}z$.
        \item L'application : 
        \begin{equation*}
            \phi_{\alpha}\longmapsto \frac{1}{1- \abs{a}^{2}}\begin{pmatrix}
                e^{i\theta/2} & -\alpha e^{i\theta/2} \\
                -\bar{\alpha}e^{-i\theta/2} & e^{-i\theta/2}
            \end{pmatrix}
        \end{equation*}
        induit un automorphisme 
        \begin{equation*}
            \Aut D \rightarrow {\rm PSU}\left(1, 1\right) = \left\{\begin{pmatrix}
                A & B \\ \overline{A} & \overline{B}
            \end{pmatrix}\ \middle| \ \abs{A}^{2} - \abs{B}^{2} = 1\right\}/\pm {\rm Id}
        \end{equation*}
    \end{itemize}
\end{proof}
\begin{corollaire}{Schwarz-Pick}{schwarz-pick}
    Soit $f : D\left(0, 1\right) \to D\left(0, 1\right)$ holomorphe. Pour $z \in D\left(0, 1\right)$ on a :
    \begin{equation*}
        \frac{\abs{f'\left(z\right)}}{1 - \abs{f\left(z\right)}^{2}} \leq \frac{1}{1 - \abs{z}^{2}}
    \end{equation*}
    avec égalité si et seulement si $f \in \Aut D\left(0, 1\right)$.
\end{corollaire}
\begin{proof}
    La dérivée de $\phi_{\alpha}\left(z\right) = \frac{z - \alpha}{1 - \bar{\alpha}z}\in \Aut D$ vérifie : 
    \begin{equation*}
        \phi_{\alpha}'\left(z\right) = \frac{1 - \abs{\alpha}^2}{\left(1 - \bar{\alpha}z\right)^{2}}, \phi_{\alpha}'\left(0\right) = 1 - \abs{\alpha}^{2}, \abs{\phi_{\alpha}'\left(\alpha\right)} = \frac{1}{1 - \abs{\alpha}^{2}}
    \end{equation*}
    Soit $z_{0} \in D$ et $g = \phi_{f\left(z_{0}\right)} \circ f \circ \phi_{-z_{0}}$. On a $g\left(0\right) = 0$ et $g : D\to D$. Le lemme de Schwarz \ref{thm:lemme_schwarz} implique que $\abs{g\left(z\right)} \leq \abs{z}$ et $\abs{g'\left(0\right)} \leq 1$. L'inégalité attendue résulte de 
    \begin{equation*}
        g'\left(0\right) = \phi_{f\left(z_{0}\right)}\left(f\left(z_{0}\right)\right)f'\left(z_{0}\right)\phi_{-z_{0}}'\left(0\right) = \frac{1}{1 - \abs{f\left(z_{0}\right)}^{2}}f'\left(z_{0}\right)\left(1 - \abs{z_{0}}\right)^{2}
    \end{equation*}
    Le cas d'égalité correspond à $\abs{g'\left(0\right)} = 1$. Alors $g$ et donc $f$ sont des automorphismes de $D$. 
\end{proof}

\begin{théorème}{Bloch-Landau}{bloch-landau}
    Soit $f \in \O\left(D\left(z_{0}, r\right)\right)$ telle que $f'\left(z_{0}\right) \neq 0$. Alors il existe $U \subset D\left(z_{0}, r\right)$ tel que $f_{\mid U}$ est un biholomorphisme de $U$ sur $f\left(U\right) = D\left(\omega_{0}, R\right)$ disque de rayon $R \geq \frac{r}{12}\abs{f'\left(z_{0}\right)}$.
\end{théorème}
\begin{proof}
    Quitte à considérer la restriction de $f$ et à remplacer $f$ par $f\left(z_{0} + rz\right)$ on peut considérer définie au voisinage du disque fermé $\overline{D}\left(0, 1\right)$. On pose
    \begin{equation*}
        m = \sup_{z \in \overline{D}\left(0, 1\right)}\left(1 - \abs{z}^{2}\right)\abs{f'\left(z\right)} \geq \abs{f'\left(0\right)}
    \end{equation*}
    La valeur $m$ est atteinte en $\alpha \in D\left(0, 1\right)$ et alors $m = \left(1 - \abs{\alpha}^{2}\right)\abs{f'\left(\alpha\right)}$. Posons $h = f\circ\phi_{-\alpha}$. On a : 
    \begin{equation*}
        h\left(0\right) = f\left(\alpha\right), h'\left(0\right) = f'\left(\alpha\right)\phi_{-\alpha}\left(0\right) = \left(1 - \abs{\alpha}^{2}\right)f'\left(\alpha\right), \abs{h'\left(0\right)} = m \geq \abs{f'\left(0\right)}
    \end{equation*}
    De plus, pour $\abs{z} < 1$, d'après le corollaire de Schwarz-Pick \ref{cor:schwarz-pick} : 
    \begin{equation*}
        \left(1-\abs{z}^{2}\right)\abs{h'\left(z\right)} = \frac{\left(1 - \abs{z}^{2}\right)\abs{\phi_{-a}'\left(z\right)}}{1 - \abs{\phi_{-\alpha}\left(z\right)}^{2}}\left(1 - \abs{\phi_{-\alpha}\left(z\right)}^{2}\right)\abs{f'\left(\phi_{-\alpha}\left(z\right)\right)} \leq m
    \end{equation*}
    Quitte à remplacer $h$ par $\frac{1}{h'\left(0\right)}\left(h-h\left(0\right)\right)$, on peut supposer $h\left(0\right) = 0$ et $m = h'\left(0\right) = 1$. Donc : 
    \begin{equation*}
        \abs{h'\left(z\right)} \leq \frac{1}{1 - \abs{z}^{2}}, z \in D\left(0, 1\right)
    \end{equation*}
    Le rayon de convergence du développement en série entière\footnote{DONC DE TAYLOR}
    \begin{equation*}
        h\left(z\right) = z + \sum_{n = 2}^{+ \infty}a_{n}z^{n}
    \end{equation*}
    est supérieur à $1$ et les inégalités de Cauchy appliquées à $h'$ sur le disque $D\left(0, \rho\right)$ donnent 
    \begin{equation*}
        h^{\left(n\right)}\left(\omega\right) = \frac{\left(n- 1\right)!}{2\pi\rho^{n-1}}\int_{0}^{2\pi}\frac{h'\left(\omega + re^{it}\right)}{e^{i\left(n - 1\right)t}}\d t \text{ et } n\abs{a_{n}} \leq \frac{1}{\left(1 - \rho^{2}\right)\rho^{n - 1}}
    \end{equation*}
    Or, $\rho \mapsto \left(1 - \rho^{2}\right)\rho^{n - 1}$ atteint son maximum en $\rho = \sqrt{\frac{n-1}{n + 1}}$. Pour ce $\rho$, on obtient : 
    \begin{equation*}
        \abs{a_{n}} \leq \frac{n + 1}{2n}\left(1 + \frac{2}{n - 1}\right)^{\left(n-1\right)/2}
    \end{equation*}
    et pour $M = 3^{3/4}/2$, $\abs{a_{n}} \leq M^{n}$ si $n \geq 2$.\\
    En reprenant la preuve du théorème d'inversion locale par la méthode des séries majorantes \ref{thm:inv_loc_serie_maj}, on obtient que $h$ est un biholomorphisme d'un ouvert $U$ sur le disque $D\left(0, R\right)$ avec : 
    \begin{equation*}
      R = \frac{1}{\left(1 + \sqrt{2}\right)M + 4M^{2}}>\frac{1}{12}
    \end{equation*}
\end{proof}

\section{Espaces de Fonctions Holomorphes}
\subsection{Convergence de Suites de Fonctions Holomorphes}
On dit qu'une suite converge localement de manière uniforme lorsque 
\begin{equation*}
    \forall x \in U, \exists r_{x} > 0, f_{n} \xrightarrow[D\left(x, r_{x}\right)]{\rm CU} f
\end{equation*}
La convergence uniforme localement est équivalente à la convergence uniforme sur tout compact.
\begin{proof}
    \begin{description}
        \item[$\Leftarrow$] $K$ étant compact, on le recouvre par un nombre fini de disques. 
        \item[$\Rightarrow$] En considérant un disque fermé inclus dans l'ouvert, on a convergence sur ce disque.   
    \end{description}
\end{proof}
On rappelle que si $\sum f_{n}$ converge uniformément sur un compact et que les $f_{n}$ sont holomorphes, $\sum f_{n} \in \O\left(K\right)$.

\begin{théorème}{Weierstrass}{}
    Soit $f_{n} \in \O\left(U\right)$, $f_{n} \xrightarrow[\forall K \subset U]{\rm CU} f$. Alors, $f \in \O\left(U\right)$ et $\forall l \in \N, f_{n}^{\left(l\right)} \xrightarrow[\forall K \subset U]{\rm CU} f^{\left(l\right)}$.
\end{théorème}
\begin{proof}
    On considère $\overline{D\left(a, r\right)}\subset U$. On a : 
    \begin{equation*}
        f_{n}\left(z\right) = \frac{1}{2i\pi}\int_{\partial D\left(a, r\right)}\frac{f_{n}\left(\omega\right)}{\omega - z}\d \omega
    \end{equation*}
    On échange limite et intégrale par convergence uniforme et on obtient les formules de Cauchy pour $f$. 
\end{proof}

\subsection{Théorèmes de Runge}
On cherche une suite de polynôme convergeant vers une fonction holomorphe.
\begin{théorème}{Runge - 1}{runge}
    Soit $f \in \O\left(D\left(z_{0}, R\right)\right)$, $K \subseteq D\left(z_{0}, R\right)$. Alors, si $\epsilon > 0$, il existe $P_{\epsilon} \in \C\left[z\right]$ tel que 
    \begin{equation*}
       \max_{K}\abs{f - P_{\epsilon}} < \epsilon
    \end{equation*} 
\end{théorème}
\begin{proof}
    On suppose d'abord $z_{0} = 0$. On pose $f\left(z\right) = \sum a_{n}z^{n}$ sur $D\left(0, R\right)$, on note $R'$ son rayon de convergence. Par la formule d'Hadamard : 
    \begin{equation*}
        \frac{1}{R^{'}} = \limsup \abs{a_{n}}^{1/n}
    \end{equation*}
    Pour $\delta > 0$, il existe $N\left(\delta\right)$ tel que $\abs{a_{n}}^{1/n} \leq 1/R' + \delta$. On choisit $r$ assez petit pour que $0 < R - r < d\left(K, \C \setminus D\left(0, R\right)\right)$.
\end{proof}

\begin{lemme}{Approximation Rationnelle}{approx-rationnelle}
    Soit $K$ un compact inclus dans $U$ ouvert de $\C$. Pour $\gamma : \left[0, 1\right] \to U \setminus K$ un chemin de classe $\cont^{1}$. Nous notons 
    \begin{equation*}
        F_{\gamma}\left(z\right) = \frac{1}{2i\pi}\int_{\gamma}\frac{f\left(w\right)}{w-z}\d w
    \end{equation*}
    Pour tout $\epsilon > 0$, il existe $P, Q \in \C\left[Z\right]$ avec $Q$ sans zéro dans $K$ tels que 
    \begin{equation*}
        \max_{z \in K} \abs{F_{\gamma}\left(z\right) - \frac{P\left(z\right)}{Q\left(z\right)}} \leq \epsilon
    \end{equation*}
\end{lemme}
\begin{proof}
    Idée : La fonction 
    \begin{equation*}
        H : \begin{array}{rcl}
            K \times \left[0, 1\right] & \longrightarrow & \C\\
            \left(z, t\right) & \longmapsto & \frac{1}{2i\pi}\frac{f\left(\gamma\left(t\right)\right)}{\gamma\left(t\right) - z}\gamma'\left(t\right)
        \end{array}
    \end{equation*}
    est continue sur le compact $K ra\times \left[0, 1\right]$ donc uniformément continue. Par conséquent
    \begin{equation*}
        \forall \epsilon > 0, \exists \delta\left(\epsilon\right) > 0, \abs{t_{1} - t_{2}} \leq \delta\left(\epsilon\right) \Longrightarrow \abs{H\left(z, t_{1}\right) - H\left(z, t_{2}\right)} \leq \epsilon, \forall z \in K
    \end{equation*}
    Ainsi, pour $n$ tel que $1/n \leq \delta\left(\epsilon\right)$, 
    \begin{equation*}
        \abs{F_{\gamma\left(z\right)} - \sum_{k = 0}^{n - 1} H\left(z, k/n\right)\frac{1}{n}} \leq \sum_{k = 0}^{n - 1}\int_{k/n}^{\left(k + 1\right)/n}\abs{H\left(z, t\right) - H\left(z, k/n\right)}\d t \leq \epsilon
    \end{equation*}
    En posant 
    \begin{equation*}
        \frac{P}{Q} = \sum_{k = 0}^{n - 1}H\left(z, k/n\right)\frac{1}{n} = \sum_{k = 0}^{n - 1}\frac{1}{2i\pi}\frac{f\left(\gamma\left(k/n\right)\right)}{\gamma\left(k/n\right) - z}\gamma'\left(k/n\right)1/n
    \end{equation*}
    on obtient une fraction rationnelle qui approche $F_{\gamma}$ sur $K$ avec $Q\left(z\right) = \prod_{k = 0}^{n - 1}\left(\gamma\left(k/n\right) - z\right) \in \C\left[z\right]$.
\end{proof}

\begin{théorème}{Runge - 2 - Rationelle Edition}{runge-rationnel}
    Soit $K \subseteq U$ un compact. Toute fonction $f \in \O\left(U\right)$ peut-être approchée uniformément sur $K$ par des fractions rationnelles sans pôle sur $K$ : 
    \begin{equation*}
        \forall f \in \O\left(U\right), \forall \epsilon > 0, \exists P, Q \in \C\left[z\right] \mathrm{ tels que } \max_{z \in K}\abs{f\left(z\right) - \frac{P\left(z\right)}{Q\left(z\right)}} \leq \epsilon
    \end{equation*}
\end{théorème}
\begin{proof}
    \begin{itemize}
        \item Construction d'un contour de $K$. Soit $\alpha \in \left]0, \frac{1}{\sqrt{2}}\right[$ et 
        \begin{equation*}
            \delta = \alpha d\left(K, \C \setminus U\right) > 0
        \end{equation*}
        Pour $k, l \in \Z$, nous notons le carré plein fermé de côté $\delta$: 
        \begin{equation*}
            \overline{\Pi}_{k, l} = \left\{x + iy \in \C \middle| k \delta \leq x \leq \left(k + 1\right)\delta, l\delta \leq y \leq \left(l + 1\right)\delta\right\}
        \end{equation*}
        et $\Pi_{k, l} = \mathring{\overline{\Pi}}_{k, l}$ le carré plein ouvert. Nous avons le pavage
        \begin{equation*}
            \C = \bigcup_{k \in \Z} \bigcup_{l \in \Z} \overline{\Pi}_{k, l}
        \end{equation*}
        Nous notons $\overline{\Pi}_{1}, \ldots, \overline{\Pi}_{J}$ la collection de tous les $\overline{\Pi}_{k, l}$ d'intersection non nulle avec $K$ et $\gamma_{1}, \ldots, \gamma_{H}$ les segments de longueur $\delta$ qui appertiennent au bord d'un unique $\overline{\Pi}_{j}$ sans appartenir à deux bords de deux carrés adjacents. Quitte à translater $K$ et réduire $\delta$, on peut supposer que $K$ contient au moins un point à l'intérieur de chaque carré. \\
        Pour tout $j$, il existe $z_{j} \in \Pi_{j} \cap K$. On a : 
        \begin{equation*}
            \abs{z_{j} - w} \leq \sqrt{2}\delta \Rightarrow w \in \overline{\Pi}_{j}
        \end{equation*}
        Or, $\delta < \alpha d\left(K, \C \setminus U\right)$ donc 
        \begin{equation*}
            \sqrt{2}\delta < d\left(K, \C \setminus U\right)
        \end{equation*}
        Par conséquent, $\overline{\Pi}_{j} \subseteq \overline{D}\left(z_{j}, \sqrt{2}\delta\right) \subseteq U$ et 
        \begin{equation*}
            \bigcup_{j} \overline{\Pi}_{j} \subseteq U \text{ et } \bigcup_{1 \leq h \leq H} \gamma_{h} \subseteq U
        \end{equation*}
        De plus, pour tout $1 \leq h \leq H$, $\gamma_{h} \cap K = \emptyset$. Sinon, si $z \in \gamma_{h} \cap K$, il existe deux carrés adjacents tel que $\gamma_{h} = \overline{\Pi}_{k\left(h\right), l\left(h\right)} \cap \overline{\Pi}_{k'\left(h\right), l'\left(h\right)}$. Comme $z$ appartient à chacun d'eux, les carrés $\overline{\Pi}_{k\left(h\right), l\left(h\right)}$ et $\overline{\Pi}_{k'\left(h\right), l'\left(h\right)}$ appartiennent tous deux à la collection des $\overline{\Pi}_{i}$. Donc $\gamma_{h}$ est un bord commun à deux carrés de la collection ce qui est exclu. 
        \item Pour tout $z \in K \setminus \bigcup\limits_{j = 1}^{J} \delta\Pi_{j}$, il existe un unique $j\left(z\right)$ avec $z \in \Pi_{j\left(z\right)}$. Pour tout $j' \neq j\left(z\right)$, 
        \begin{equation*}
            \frac{1}{2i\pi} \int_{\delta\Pi_{j'}}\frac{f\left(w\right)}{w-z}\d w = 0
        \end{equation*}
        et 
        \begin{equation*}
            f\left(z\right) = \sum_{j = 1}^{J}\frac{1}{2i\pi}\int_{\delta\Pi_{j}}\frac{f\left(w\right)}{w-z}\d w
        \end{equation*}
        \item Si $z \in K \setminus \bigcup_{j = 1}^{J}\delta \Pi_{j}$, on a, par définitions des $\gamma_{h}$, les contributions des autres bords des $\delta \Pi_{j}$ s'annulant deux à deux : 
        \begin{equation*}
            f\left(z\right) = \sum_{j = 1}^{J}\frac{1}{2i\pi} \int_{\delta \Pi_{j}}\frac{f\left(w\right)}{w - z}\d w = \sum_{h = 1}^{H}\frac{1}{2i\pi}\int_{\gamma_{h}}\frac{f\left(w\right)}{w-z} \d w
        \end{equation*}
        Comme $\gamma_{h} \cap K = \emptyset$, les intégrales $\int_{\gamma_{h}}\frac{f\left(w\right)}{w-z} \d w$ sont continues par rapport à $z \in K$. L'égalité précédente est donc encore satisfaite pour $z \in K \cap \bigcup_{j = 1}^{J} \delta \Pi_{j}$.\\
        D'après le lemme précédent \ref{lemme:approx-rationnelle}, pour tout $1 \leq h \leq H$, il existe $P_{h}, Q_{h} \in \C\left[z\right]$, $Q_{h}$ sans zéros dans $K$ avec
        \begin{equation*}
            \sup_{z \in K}\abs{F_{\gamma_{h}}\left(z\right) - \frac{P_{h}\left(z\right)}{Q_{h}\left(z\right)}} \leq \frac{\epsilon}{H}
        \end{equation*}
        Posons alors $Q = \prod_{h = 1}^{H}Q_{h} \in \C\left[z\right]$ et 
        \begin{equation*}
            P = Q\prod_{h = 1}^{H}\frac{P_{h}}{Q_{h}} \in C\left[z\right]
        \end{equation*}
        On obtient alors : 
        \begin{equation*}
            \abs{f\left(z\right) - \frac{P\left(z\right)}{Q\left(z\right)}} \leq \epsilon
        \end{equation*}
    \end{itemize}
\end{proof}

\begin{lemme}{Approximation Polynômiale des Fractions Rationnelles}{rationnelle-polynomes}
	Si $K \subseteq \C$ compact a un complémentaire connexe, alors pour tout $z_{0} \notin K$, la fonction rationnelle $\frac{1}{z - z_{0}}$ peut être approchée uniformément sur $K$ par des polynômes de $\C\left[z\right]$. 
\end{lemme}
\begin{proof}
	\begin{itemize}
		\item Comme $K$ est compact, il est borné donc inclus dans un disque $D\left(0, R\right)$. Pour $z_{1} \in \C \setminus D\left(0, R\right)$, on a : 
			\begin{equation*}
				\frac{1}{z - z_{1}} = - \frac{1}{z_{1}}\frac{1}{1 - \frac{z}{z_{1}}} = \sum_{n = 1}^{+ \infty}-\frac{z^{n}}{z_{1}^{n + 1}}
			\end{equation*}
			série uniformément convergente pour tout $z \in K$. Donc la fraction $\frac{1}{z - z_{1}}$ peut être approximée uniformément par des polynômes sur $K$. De même pour toute puissance. 
		\item Comme $\C \setminus K$ est connexe de $\C$ donc connexe par arcs, il existe un chemin $\gamma : \left[0, 1\right] \to \C \setminus K$ avec $\gamma\left(0\right) = z_{0}$ et $\gamma\left(1\right) = z_{1}$. L'image $\gamma\left(\left[0, 1\right]\right)$ est compacte disjointe de $K$ donc $\rho = \frac{1}{2}d\left(\gamma\left(\left[0, 1\right]\right), K\right) > 0$. Choississons une subdivisions de $\gamma\left(\left[0,1\right]\right)$ de pas au plus $\rho$. Pour toute paire de points $w, w' \in \gamma\left(\left[0, 1\right]\right)$ avec $\abs{w - w'} \leq \rho$, pour $z \in K$, on a : 
			\begin{equation*}
				\frac{1}{z-w} = \frac{1}{z - w'}\sum_{n = 0}^{+\infty}\left(\frac{w - w'}{z - w'}\right)^{n}
			\end{equation*}
			et cette comme est uniformément convergente car $\abs{\frac{w - w'}{z - w'}}\leq 1/2$. Ainsion a une approximation de $\frac{1}{z - w}$ par des polynômes en $\frac{1}{z - w'}{z - w}$ par des polynômes en $\frac{1}{z - w'}$.
	\end{itemize}
\end{proof}

\begin{théorème}{Runge - 3 - Polynomial Edition}{runge-polynomes}
	Soit $K \subseteq U$ compact de complémentaire connexe. Alors, les fonctions holomorphes dans $U$ sont uniformément approximables sur $K$ par des polynômes: 
	\begin{equation*}
		\forall f \in \O\left(U\right), \forall \varepsilon > 0, \exists P_{\varepsilon} \in \C\left[z\right] \mathrm{ tel que } \max_{z\in K} \abs{f\left(z\right) - P_{\epsilon}\left(z\right)}\leq \varepsilon
	\end{equation*}
\end{théorème}
\begin{proof}
	Ceci découle du lemme \ref{lemme:rationnelle-polynomes} que l'on combine à la version rationnelle \ref{thm:runge-rationnel} du théorème de Runge
\end{proof}

\subsection{Théorème de Montel}
\begin{définition}{Bornitude Uniforme}{born-uni}
	Une famille $\mathfrak{F} \subseteq \O\left(U\right)$ est uniformément bornée sur les compacts de $U$ si pour tout $K \subseteq U$ compact, il existe $M_{K} \geq 0$ telle que 
	\begin{equation*}
		\abs{f\left(z\right)} \leq M_{K}, z\in K, f \in \mathfrak{F}
	\end{equation*}
\end{définition}

\begin{définition}{Uniforme Équicontinuité}{born-equicont}
	Une famille $\mathfrak{F}\subseteq \O\left(U\right)$ est uniformément équicontinue sur les compacts de $U$ si pour tout $K \subseteq U$ compact, pour tout $\varepsilon > 0$, il existe $\delta = \delta\left(K, \varepsilon\right) > 0$ tel que 
	\begin{equation*}
		\left(z, w \in K \text{ avec } \abs{z - w} \leq \delta\right) \Longrightarrow \abs{f\left(z\right) - f\left(w\right)} \leq \varepsilon, f\in \mathfrak{F}
	\end{equation*}

	
\end{définition}


\begin{propositionfr}{Exemples}{}
		\begin{enumerate}
			\item Sur $\overline{D}\left(0, 1\right)$, la famille des $\left(z^{n}\right)$ est uniformément bornée mais pas uniformément équicontinue. Elle l'est sur $D\left(0, 1\right)$.
			\item Sur $\left[0, 1\right[$, la famille $\left(\sin\left(2\pi nx\right)\right)$ est uniformément bornée par $1$ mais n'est uniformément équicontinue sur aucun compact car les dérivées sont non bornées en $n$. 
	\end{enumerate}

\end{propositionfr}

\begin{définition}{Normalité}{}
	Une famille $\mathfrak{F} \subseteq \O\left(U\right)$ est dite normale si pour toute suite de fonctions $f_{n}$, il existe une suite extraite uniformément convergente sur tout compact de $U$. 
	
\end{définition}

\begin{définition}{Exhaustivité}{}
	Une suite croissante de compacts $\left(K_{n}\right)$ de $U$ est exhaustive si tout compact de $U$ est contenu dans un $K_{n}$. 
\end{définition}

\begin{théorème}{Suite Exhaustive}{suite-exhaust}
	Pour tout ouvert $U$ de $\C$, il existe une suite croissante de compacts exhaustive de réunion $U$. 	
\end{théorème}
\begin{proof}
Considérons la suite $B_{n}$ des disques fermés contenus dans $U$ dont les coordonnées du centre et le rayon sont rationnels. Ils forment un ensemble dénombrable. La suite des $K_{i} = \cup_{n = 1}^{i} B_{n}$ convient
\end{proof}

\begin{théorème}{Mantel}{mantel}
	Soit $\mathfrak{F} \subseteq \O\left(U\right)$ une famille de fonctions uniformément bornées sur les compacts de $U$. Alors
	\begin{enumerate}
		\item $\mathfrak{F}$ est uniformément équicontinue sur les compacts de $U$. 
		\item $\mathfrak{F}$ est une famille normale. 
	\end{enumerate}
\end{théorème}
\begin{proof}
	\begin{enumerate}
		\item Soit $K$ un compact de $U$, $r < \frac{1}{3}d\left(K, \C - U\right)$. 
        Pour $w \in K$, la famille $\mathfrak{F}$ est uniformément bornée sur $L = \left\{u \in U, d\left(u, K\right) \leq 2r\right\}$.
        Pour $z \in D\left(w, r\right)$ la formule de Cauchy donne : 
        \begin{equation*}
            f\left(z\right) - f\left(w\right) = \frac{1}{2i\pi}\int_{\partial D\left(w, 2r\right)}f\left(u\right)\left(\frac{1}{u - z} - \frac{1}{u - w}\right)\d u
        \end{equation*}
        Or, $\abs{u - z} > r$ et 
        \begin{equation*}
            \abs{\frac{1}{u - z} - \frac{1}{u - w}} = \frac{\abs{z - w}}{\abs{u - z}2r} \leq \frac{\abs{z - w}}{2r^{2}}
        \end{equation*}
        Ainsi
        \begin{equation*}
            \abs{f\left(z\right) - f\left(w\right)} \leq \frac{1}{2\pi} M_{L}\frac{\abs{z - w}}{2r^{2}}2\pi 2r \propto\abs{z - w}
        \end{equation*}
        Pour tout $\epsilon > 0$, on peut choisir convenablement $\delta < r$ de telle sorte qu'on a bien l'uniforme équicontinuité sur $K$.
        \item Soit $r < \frac{1}{2}d\left(K, \C - U\right)$. On pose
        \begin{equation*}
            K_{r} = \left\{z \in \C, d\left(z, K\right) \leq r\right\} \subseteq K_{2r} \subseteq U
        \end{equation*}
        Par hypothèses, $M_{0} = \sup_{n \in \N, z \in K_{2r}}\abs{f_{n}\left(z\right)} < + \infty$. Par l'inégalité de Cauchy, 
        \begin{equation*}
            M_{1} = \sup_{n \in \N, z \in K_{r}} \abs{f_{n}'\left(z\right)} \leq \frac{M_{0}}{r} < + \infty
        \end{equation*}
        Soient $z_{1}, z_{2} \in K$. Si $d\left(z_{1}, z_{2}\right) \leq r$, alors $\abs{z_{1}, z_{2}} \subseteq K_{r}$ et par théorème des accroissements finis
        \begin{equation*}
            \abs{f_{n}\left(z_{1}\right) - f_{n}\left(z_{2}\right)} \leq M_{1} d\left(z_{1}, z_{2}\right)
        \end{equation*}
        À l'inverse, si $d\left(z_{1}, z_{2}\right) \geq r$, on a : 
        \begin{equation*}
            \abs{f_{n}\left(z_{1}\right) - f_{n}\left(z_{2}\right)} \leq \frac{2M_{0}}{r}d\left(z_{1}, z_{2}\right)
        \end{equation*}
        Par théorème d'Ascoli, on peut extraire des $f_{n\mid_{K}}$ une suite uniformément convergente sur $K$. \\
        Pour $\nu \in \N^{*}$, on considère une suite croissante exhaustive de compacts $\left(K_{\nu}\right)$. D'après ce qui précède, il existe une partie infinie $S_{1} \subseteq \N$ telle que la suite $\left(f_{n}\right)_{n \in S_{1}}$ converge uniformément sur $K_{1}$. Par récurrence sur $\nu$, on construit une suite de parties infinies emboîtées $S_{\nu} \subseteq \cdots \subseteq S_{1} \subseteq \N$ telle que la sous-suite $\left(f_{n}\right)_{n \in S_{\nu}}$ converge uniformément sur $K_{\nu}$. Par extraction diagonale, on construit une sous-suite $\left(f_{n_{i}}\right)_{i \in \N}$ qui converge uniformément sur tous les compacts $K_{\nu}$.
	\end{enumerate}
\end{proof}

\subsection[Compacts de O(U)]{Compacts de $\O\left(U\right)$}
On munit $\O\left(U\right)$ de la topologie de la convergence uniforme sur tout compact. C'est une topologie invariante par translation dans laquelle les 
\begin{equation*}
    V\left(K, \epsilon\right) = \left\{f \in \O\left(U\right)\middle| \sup\limits_{z \in K}\abs{f\left(z\right)} < \epsilon\right\}
\end{equation*}
pour $K$ un compact de $U$ et $\epsilon>0$ constituent un système de voisinage de $0$. L'addition et l'homothétie sont continues pour cette topologie. \\
On définit de plus :
\begin{equation*}
    p_{K_{n}} : \begin{array}{rcl}
        \O\left(U\right) & \longrightarrow & \R^{+}\\
        f & \longmapsto & \sup_{z \in K_{n}}\abs{f\left(z\right)}
    \end{array}
\end{equation*}
et 
\begin{equation*}
    d\left(f, g\right) = \sum_{i \geq 1}2^{-i}\min\left(1, p_{K_{i}}\left(f - g\right)\right)
\end{equation*}
où les $K_{i}$ forment une suite croissante exhaustive de compacts de $U$. 
\begin{lemme}{Métrisabilité}{metris-odeu}
    La topologie de la convergence sur tout compact est métrisable par la distance $d$ sur $\O\left(U\right)$.
\end{lemme}
\begin{proof}
    \begin{itemize}
        \item Soit $K$ un compact. Soit $i \in \N$ tel que $K \subseteq K_{i}$ Alors, $d\left(f, 0\right) < 2^{-i} \epsilon$ implique $p_{K_{i}}\left(f\right) < \epsilon$. Donc : $B\left(0, 2^{-i}\epsilon\right) \subseteq V\left(K_{i}, \epsilon\right) \subseteq V\left(K, \epsilon\right)$. Donc $V\left(K, \epsilon\right)$ contient une boule centrée en $0$ pour $d$. 
        \item Soit $B\left(0, \epsilon\right)$ une boule pour $d$. Soit $i \in \N$ tel que $2^{-i} < \frac{\epsilon}{2}$. Alors, si $f \in V\left(K_{i}, \epsilon/2\right)$
        \begin{equation*}
            d\left(f, 0\right) \leq \sum_{j = 1}^{i}2^{-j}p_{K_{i}}\left(f\right) + \sum_{j = i + 1}^{+\infty}2^{-j} < \epsilon/2 + 2^{-i} < \epsilon
        \end{equation*}
        Donc $V\left(K_{i}, \epsilon/2\right) \subseteq B\left(0, \epsilon\right)$.
    \end{itemize}
\end{proof}
L'espace vectoriel $\O\left(U\right)$ localement convexe, métrisable et complet est un espace de Fréchet. 

\begin{théorème}{Compacts de $\O\left(U\right)$}{compact-odeu}
    Les compacts de $\O\left(U\right)$ sont les fermés bornés.
\end{théorème}
\begin{proof}
    \begin{itemize}
        \item Soit $A \subseteq \O\left(U\right)$ une partie compacte. 
        Comme $\O\left(U\right)$ est métrisable, il est séparé, et donc $A$ est fermée. 
        Soit $K \subseteq U$ compact. 
        L'application $p_{K}$ est continue sur $\O\left(U\right)$. Donc comme $A$ est compacte, $p_{K}\left(A\right)$ est compacte donc bornée dans $\R$. 
        Les applications $f\in A$ sont donc uniformément bornées sur $K$. 
        Pour tout compact $K$ de $U$, il existe donc $\epsilon_{A} > 0$ tel que $A \subseteq V\left(K, \epsilon_{A}\right)$. 
        Ainsi pour tout $V\left(K, \epsilon\right)$ d'un système fondamental de voisinages de $0$, il existe $\lambda \in \R^{+}$ avec $A \subseteq V\left(K, \lambda \epsilon\right)$.
        Donc $A$ est bornée.
        \item Soit $A$ une partie fermée bornée de $\O\left(U\right)$. 
        Comme $A$ est bornée, pour tout $K$ compact de $U$, les applications $f \in A$ sont uniformément bornées sur $K$. Soit $\left(f_{n}\right)$ une suite d'éléments de $A$. 
        Par théorème de Mantel \ref{thm:mantel}, cette suite est uniformément bornée donc admet une sous-suite convergente uniformément sur tout compact de $U$ de limite $f \in \O\left(U\right)$.
        Comme $A$ est fermée, $f \in A$ donc, $A$ est séquentiellement compact pour une topologie métrisable, donc par théorème de Bolzano-Weierstrass est compact. 
    \end{itemize}
\end{proof}

\section{Aparté Réel}
\subsection{Méthode de Laplace}
\begin{théorème}
	{Méthode de Laplace Réelle}{laplace-reelle}
Soient $g, h$ deux fonctions réelles continues par morceaux dans $\left]0, +\infty\right[$ vérifiant
	\begin{enumerate}
		\item $\int_{0}^{+\infty} \abs{g\left(x\right)}e^{h\left(x\right)}\d x < +\infty$
		\item Il existe $\delta_{0} > 0$ tel que $\forall \delta < \delta_{0}$ et tout $x \geq \delta$, on a $h\left(x\right) \leq h\left(\delta\right)$.
		\item Au voisinage de $0$, on a les développements asymptotiques
			\begin{equation*}
				g\left(x\right) \sim Ax^{\alpha}, h\left(x\right) = a - cx^{\beta} + o\left(x^{\beta}\right)
			\end{equation*}
			où $A \neq 0$, $\alpha > -1, \beta > 0, c > 0$.
			Alors, 
			\begin{equation*}
				I\left(t\right) \sim \frac{A}{\beta}\Gamma\left(\frac{\alpha + 1}{\beta}\right)e^{\alpha t}\left(ct\right)^{-\left(\alpha + 1\right)/\beta}
			\end{equation*}
	\end{enumerate}
\end{théorème}
\begin{proof}
	Calcul.
\end{proof}


\begin{corollaire}
	{}{}
	Soit $g, h : \R \to \R$, des fonctions de classes $\cont^{2}$ telles que 
	\begin{itemize}
		\item $\int_{\R}\abs{g\left(x\right)}e^{h\left(x\right)}\d x$ converge
		\item $h'$ ne change de signe qu'en un seul point $c \in \R$ où $h$ atteint un maximum
		\item $g\left(c\right) \neq 0$ et $h''\left(c\right) < 0$
	\end{itemize}
	Alors, pour $t$ au voisinage de $+ \infty$, 
	\begin{equation*}
		\int_{\R}g\left(x\right)e^{th\left(x\right)}\d x \sim{x \to + \infty} \sqrt{pi}g\left(c\right)e^{th\left(c\right)}\sqrt{\frac{2}{-th''\left(c\right)}}
	\end{equation*}
\end{corollaire}


\subsection{Comportement Asymptotique d'une Intégrale}
\begin{théorème}
	{Laplace Complexe}{laplace}
	Supposons que \begin{equation*}
		\int_{0}^{+\infty} \abs{g\left(\gamma\left(s\right)\right)}e^{\Re h\left(\gamma\left(s\right)\right)}\abs{\gamma'\left(s\right)}\d s
	\end{equation*}
	est convergente. Pour $t$ au voisinage de $+\infty$, on a: 
	\begin{equation*}
		\int_{\gamma}g\left(z\right)e^{th\left(z\right)}\d z \sim \frac{A}{m}\Gamma\left(\frac{n + 1}{m}\right)\left(\rho t\right)^{-\frac{n + 1}{m}}e^{at + \left(n+1\right)i\omega_{k}}
	\end{equation*}
\end{théorème}
\begin{proof}
	Calcul. Encore.
\end{proof}

\subsection{Fonctions Harmoniques}
Ici, on identifie $\C$ à $\R^{2}$
\begin{définition}
	{Fonctions Harmonique}{}
	Soit $f : U \to \R$ de classe $\cont^{2}$. Le \emph{Laplacien} de $f$ est 
	\begin{equation*}
		\Delta f = \partial_{xx}f'' + \partial_{yy}f''
	\end{equation*}
	La fonction $f$ est dite $\emph{harmonique}$ si $\Delta f = 0$.
\end{définition}

\begin{propositionfr}
	{Harmonicité et Holomorphie}{harmo-holo}
	Si $U$ est simplement connexe, toute fonction harmonique sur $U$ est la partie réelle d'une fonction holomorphe, unique à une constante imaginaire près. 
\end{propositionfr}
\begin{proof}
	Soit $f: U \to \R$ une fonction harmonique. 
	Alors $g = \partial_{y}f + i\partial_{x}f$ est une fonction de classe $\cont^{1}$ sur $U$. 
	La fonction $g$ est holomorphe sur $U$ car $g = P + iQ$ avec $\partial_{x}P = \partial_{y}Q = \partial_{xy}f$ et $\partial_{y}P = \partial_{yy}f = -\partial_{xx}f = -\partial_{x}Q$. 
	Notons $G = A +iB$ une primitive de $g$ sur $U$. 
	Ainsi $G'\left(z\right) = \partial_{x}A + i\partial_{x}B = \partial_{y}B -i\partial_{y}A = g\left(z\right) = \partial_{y}f + i\partial_{x}f$.
	Les équations de Cauchy-Riemann montrent alors que $\phi = f - iA$ est holomorphe. Si $f = \Re\left(\phi\right) = \Re\left(\psi\right)$ avec $\phi, \psi \in \O\left(U\right)$, alors $\phi - \psi = c$ est constante sur $U$ avec $\Re\left(c\right) = 0$. 
\end{proof}

\begin{corollaire}
	{Harmonicité Locale}{local-harmo}
	Toute fonction harmonique est localement la partie réelle d'une fonction holomorphe.
\end{corollaire}

\begin{corollaire}
	{Classe Harmonique}{classe-harmo}
	Toute fonction harmonique est de classe $\cont^{\infty}$.
\end{corollaire}

\begin{corollaire}
	{Composition Harmonique}{}
	Soit $\phi: U \to V$ holomorphe et $f : \to V$ harmonique, alors $f \circ \phi : U \to \R$ est harmonique. 
\end{corollaire}

\begin{corollaire}
	{Harmonicité sur l'Intérieur}{harmo-inté}
	Si deux fonctions harmoniques sur $U$ connexe coïncident sur un sous-ensemble $V \subseteq U$ d'intérieur non vide, alors elles sont égales sur $U$
\end{corollaire}
\begin{proof}
	Soit $u = f_{1} - f_{2}$ la différence de deux fonctions satisfaisant les hypothèses du corollaire et 
	\begin{equation*}
		F = \{z \in U, u\left(z\right) = 0\}
	\end{equation*}
	Comme $u$ est continue, $F$ est fermé et d'intérieur non vide. 
	Montrons que $F = \mathring{F}$. Soit $z' \in \partial F$ et $\left(z_{n}\right)$ une suite à termes dans $\mathring{F}$ convergeant vers $z'$. 
	Comme $z' \in U$ ouvert, il existe $r' > 0$ avec $D\left(z', r'\right) \subseteq U$ et quitte à supprimer les premiers termes de la suite, $z_{n} \in D\left(z', r'\right)$ pour $n \in \N$. 
	De plus, pour tout $n \in \N$, il existe $r_{n} > 0$ avec $D\left(z_{n}, r_{n}\right) \subseteq \mathring{F} \cap D\left(z', r'\right)$.\\
	Par \ref{propo:harmo-holo}, $u = \Re f$ avec $f \in \O\left(D\left(z', r'\right)\right)$. Donc $f$ constante sur tout $D\left(z_{n}, r_{n}\right)$, $f$ est constante sur $D\left(z', r'\right), u = 0$ sur $D\left(z', r'\right)$ et enfin $z' \in \mathring{F}$.\\
	L'ensemble $F$ est un ouvert fermé non vide de $U$ conxnexe donc $F = U$ et $u = 0$ sur $U$.
\end{proof}

\begin{définition}
	{Moyenne}{moyenne}
	La fonction $f: U \to \C$ vérifie la propriété de la moyenne si pour tout $\overline{D\left(a, r\right)} \subseteq U$, 
	\begin{equation*}
		f\left(a\right) = \frac{1}{2\pi}\int_{0}^{2\pi} f\left(a+re^{i\theta}\right)\d \theta
	\end{equation*}
\end{définition}

\begin{propositionfr}
	{Moyenne d'une Fonction Harmonique}{moyenne-harmo}
	Si $f$ est harmonique, alors elle vérifie la propriété de la moyenne.
\end{propositionfr}
\begin{proof}
	D'après la proposition \ref{propo:harmo-holo}, sur l'ouvert simplement connexe $D\left(a, r + \varepsilon\right)$, $f = \Re \phi$ avec $\phi : D\left(a, r + \varepsilon\right) \to \C$ holomorphe. D'après la formule de Cauchy
	\begin{equation*}
		\phi\left(a\right) = \frac{1}{2i\pi}\int_{\partial D\left(a, r\right)} \frac{\phi\left(u\right)}{u - a} = \frac{1}{2\pi}\int_{0}^{2\pi}\phi\left(a + re^{i\theta}\right) \d \theta
	\end{equation*}
	En passant à la partie réelle des deux côtés, on a le résultat. 
\end{proof}

\begin{théorème}
	{Harmonicité et Moyenne}{}
	Soit $f : U \to \R$ une fonction continue qui vérifie la propriété de la moyenne. Alors $f$ est harmonique.
\end{théorème}
\begin{proof}
	Si $f$ est de classe $\cont^{\infty}$, écrivons son développement de Taylor au point $a$: 
	\begin{equation*}
		\begin{aligned}
		f\left(a + re^{i\theta}\right) =& f\left(a\right) + r\cos\left(\theta\right)\partial_{x}f\left(a\right) + r \sin\theta \partial_{y}f\left(a\right) +\\
		& +\frac{1}{2}\left(r^{2}\cos^{2}\left(\theta\right)\partial_{x, x}f\left(a\right) + 2r^{2}\cos \theta \sin \theta \partial_{x, y}f\left(a\right) + r^{2}\sin^{2}\theta \partial_{y, y}f\left(a\right)\right) + \epsilon\left(r, \theta\right)
		\end{aligned}
	\end{equation*}
	avec $\varepsilon\left(r, \theta\right) = \O\left(r^{3}\right)$ pour $0 < r \leq R$ si $\overline{D\left(a, R\right)} \subset U$. En prenant la moyenne sur $\partial D\left(a, r\right)$, on a: 
	\begin{equation*}
		f\left(a\right) = f\left(a\right) + \frac{r^{2}}{4}\left(\partial_{x, x}f\left(a\right) + \partial_{y, y}f\left(a\right)\right) + \frac{1}{2\pi}\int_{0}^{2\pi} \epsilon\left(r, \theta\right) \d \theta
	\end{equation*}
	Donc $\Delta f= \frac{-2}{r^{2}\pi}\int_{0}^{2\pi}\epsilon\left(r, \theta\right) \d \theta$ donc $\Delta f = 0$ et $f$ est harmonique. \\
	Pour conclure, il suffit donc de montrer que la propriété de la moyenne implique la classe $\cont^{\infty}$
	Notons
	\begin{equation*}
		U_{\epsilon} = \left\{z \in U \middle| d\left(z, U^{\complement}\right) > \sqrt{\epsilon}\right\}
	\end{equation*}
	Soit $\psi_{\epsilon} : \R_{+} \to \R_{+}$ une fonction de classe $\cont^{\infty}$ de support inclus dans $\left[0, \epsilon\right]$ et $\int_{0}^{+\infty}\psi_{\epsilon} = 1$. Enfin, pour $z \in U_{\epsilon}$, posons
	\begin{equation*}
		F\left(z\right) = \iint_{\R^{2}} f\left(u\right)\psi_{\epsilon}\left(\abs{u - z}^{2}\right)\d u
	\end{equation*}
	où $\d u$ désigne la mesure de Lebesgue. D'après le théorème de dérivation sous le signe somme, la fonction $F$ est de classe $\cont^{\infty}$ sur $U$ et
	\begin{equation*}
		F\left(z\right) = \iint_{\R^{2}} f\left(u - z\right) \psi_{\epsilon}\left(\abs{u}^{2}\right)\d u = \int_{0}^{\sqrt{\epsilon}} \psi_{\epsilon}\left(r^{2}\right) \left(\int_{0}^{2\pi} f\left(z + re^{i\theta}\right) \d \theta\right)r \d r
	\end{equation*}
	D'après la propriété de la moyenne, 
	\begin{equation*}
		F\left(z\right) = 2\pi f\left(z\right) \int_{0}^{\sqrt{\epsilon}} \psi_{\epsilon}\left(r^{2}\right) r \d r = \pi f\left(z\right)
	\end{equation*}
	Donc $f$ est de classe $\cont^{\infty}$ comme $F$. 
\end{proof}

\begin{corollaire}
	{Fermeture des Fonctions Harmoniques}{}
	L'ensemble des fonctions harmoniques sur $U$ est un fermé de $\cont^{0}\left(U, \R\right)$ muni de la topologie de la convergence uniforme sur tout compact inclus dans $U$. 
\end{corollaire}

\begin{théorème}
	{Principe du Maximum}{maximum-moyenne}
	Soit $U \subseteq \R$ un ouvert et $f : U \to \C$ une fonction continue satisfaisant la propriété de la moyenne. 
	Si $\abs{f}$ admet un maximum local en $a \in U$, alors $f$ est constante sur la composante connexe de $U$ contenant $a$. 
\end{théorème}
\begin{proof}
	Si $f\left(a\right) = 0$, le résultat est clair. Supposons $f\left(a\right) \neq 0$. Quitte à multiplier $f$ par une constante complexe, nous pouvons supposer $f\left(a\right) \in \R_{+}^{*}$. Choissons $r \geq 0$ assez petit, tel que $M\left(r\right) = \sup_{\theta}\abs{f\left(a + re^{i\theta}\right)} \leq \abs{f\left(a\right)}$.
	D'après la propriété de la moyenne, 
	\begin{equation*}
		f\left(a\right) = \frac{1}{2\pi}\int_{0}^{2\pi}f\left(a + re^{i\theta}\right)\d \theta
	\end{equation*}
	donc $f\left(a\right) = M\left(r\right)$. Ainsi, $g\left(z\right) = \Re\left(f\left(a\right) - f\left(z\right)\right)\geq 0$ pour $\abs{z - a} = r$ et $g\left(z\right) = 0$ si et seulement si $f\left(z\right) = f\left(a\right)$. 
	Or, la valeur moyenne de $g\left(z\right)$ sur le cercle de centre $a$ et de rayon $r$ est nulle. Comme $g$ est continue positive, $g = 0$ sur ce cercle et donc $f = f\left(a\right)$ pour $\abs{z - a} = r$. \\
	La fonction harmonique $f$ coïncide avec une fonction constante à $f\left(a\right)$ sur $D\left(a, r\right)$. D'après le corollaire \ref{cor:harmo-inté}, $f = f\left(a\right)$ sur la composante connexe de $U$ contenant $a$. 
\end{proof}

\begin{corollaire}
	{Bornitude et Constance}{}
	Supposons $U$ connexe borné. Soit $f : \overline{U} \to \C$ continue possédant la propriété de la moyenne dans $U$. Soit $M = \sup_{z \in \partial U} \abs{f\left(z\right)}$. Alors, 
	\begin{enumerate}
		\item $\abs{f\left(z\right)} \leq M, z \in U$
		\item Si $\abs{f\left(a\right)} = M$ en $a \in U$ alors $f$ est constante.
	\end{enumerate}
\end{corollaire}
\begin{proof}
	Comme $z \mapsto \abs{f\left(z\right)}$ est continue sur $\overline{U}$ compact, il existe $a\in \overline{U}$ avec $\abs{f\left(a\right)} = \sup_{z \in \overline{U}}\abs{f\left(z\right)} = M'$. Si $a \in U$, par le principe du maximum, $f$ est constante sur un ouvert fermé non vide de $U$ connexe donc $f$ est constante sur $U$. Comme $f$ est continue, $f$ est constante sur $\overline{U}$. \\
	Si $\abs{f\left(a\right)} \neq M'$ pour tout $a \in U$, alors $M = M'$ et le corollaire suit. 
\end{proof}

\subsection{Formule de Poisson}
\begin{propositionfr}
	{Formule de Poisson}{poisson}
	Soit $g : U \to \R$, harmonique et $\overline{D\left(a, R\right)} \subseteq U$. Alors, pour tout $z \in D\left(a, R\right)$, 
	\begin{equation*}
		g\left(z\right) = \frac{1}{2\pi}\int_{0}^{2\pi}\frac{R^{2} - \abs{z - a}^{2}}{\abs{a + Re^{i\theta} - z}^{2}}g\left(a + Re^{i\theta}\right)\d \theta
	\end{equation*}
\end{propositionfr}
\begin{proof}
	Quitte à translater, il suffit de traiter le cas $a = 0$. D'après le corollaire \ref{propo:harmo-holo}, dans le disque $D\left(0, R\right)$, $g$ est la partie réelle d'une fonction holomorphe $f\left(z\right) = \sum_{n \geq 0} a_{n}z^{n}$ uniquement déterminée par la condition $f\left(0\right) \in \R$. 
	\begin{equation*}
		g\left(re^{i\theta}\right) = a_{0} + \frac{1}{2}\sum_{n \geq 1}r^{n}\left(a_{n}e^{in\theta} + \overline{a_{n}} e^{-in\theta}\right)
	\end{equation*}
	la convergence étant normale par rapport à $\theta \in \left[0, 2\pi\right]$. Les coefficients sont donnés par les formules intégrales: 
	\begin{equation*}
		\begin{aligned}
			a_{0} =& \frac{1}{2\pi}\int_{0}^{2\pi} g\left(Re^{i\theta}\right)\d \theta\\
			a_{n} =& \frac{1}{\pi}\int_{0}^{2\pi} \frac{g\left(Re^{i\theta}\right)\d\theta}{\left(Re^{i\theta}\right)^{n}}, n\in \N^{*}
		\end{aligned}
	\end{equation*}
	Ainsi, par convergence normale, nous pouvons échanger les signes somme et intégrale et 
	\begin{equation*}
		f\left(z\right) = \frac{1}{2\pi}\int_{0}^{2\pi}g\left(Re^{i\theta}\right) \left(1 + 2\sum_{n \geq 1}\left(\frac{z}{Re^{i\theta}}\right)^{n}\right)\d \theta
	\end{equation*}
	Or, $1 + 2 \sum_{n \geq 1} \left(\frac{z}{Re^{i\theta}}\right)^{n} = \frac{Re^{i\theta} + z}{Re^{i\theta} - z}$, d'où, pour $\abs{z} < R$, 
	\begin{equation*}
		f\left(z\right) = \frac{1}{2\pi}\int_{0}^{2\pi}g\left(Re^{i\theta}\right)\frac{Re^{i\theta} + z}{Re^{i\theta} - z}\d \theta
	\end{equation*}
	D'où le résultat en prenant la partie réelle des deux membres de l'égalité.
\end{proof}

\begin{corollaire}
	{Problème de Dirichlet Harmonique}{}
	Soit $u : \partial D\left(a, R\right) \to \R$ continue. Alors, il existe une unique fonction continue $f: \overline{D\left(a, R\right)} \to \R$ telle que $f_{\mid \partial D\left(a, R\right)} = u$ et $f$ est harmonique sur $D\left(a, R\right)$.
\end{corollaire}

\begin{propositionfr}
	{Inégalité d'Harnack}{}
	Soit $u : \overline{D\left(a, R\right)} \to \R_{+}$ continue harmonique dans $D\left(a, R\right)$. Si, $\abs{z - a} = r$, alors:
	\begin{equation*}
		u\left(a\right) \frac{R - r}{R + r} \leq u\left(z\right) \leq u\left(a\right) \frac{R + r}{R - r}
	\end{equation*}
\end{propositionfr}
\begin{proof}
	Les deux inégalités sont analogues. Montrons celle de droite:
	\begin{equation*}
		\begin{aligned}
			u\left(a + re^{i \theta}\right) =& \frac{1}{2\pi}\int_{0}^{2\pi}\frac{R^{2} - r^{2}}{R^{2} - 2Rr\cos\left(\theta - t\right) + r^{2}}u\left(a + Re^{it}\d t\right)\\
				\leq & \frac{1}{2\pi} \int_{0}^{2\pi}\frac{R + r}{R - r}u\left(a + Re^{it}\right)\d t = u\left(a\right)\frac{R + r}{R - r}
		\end{aligned}
	\end{equation*}
\end{proof}


\begin{théorème}
	{Harnack}{harnack}
	Soit $u_{n}$ une suite croissante de fonctions harmoniques sur $U$ connexe. Alors, on a l'une des deux propositions suivantes
	\begin{enumerate}
		\item $u_{n}$ converge uniformément sur tout compact de $U$ vers $u$ harmonique
		\item $u_{n}$ diverge vers $+ \infty$. 
	\end{enumerate}
\end{théorème}
\begin{proof}
	Quitte à changer $u_{n}$ en $u_{n} - u_{0}$, nous pouvons supposer que $\left(u_{n}\right)$ est une suite de fonctions positives. Pour tout $z \in U$, la suite $u_{n}\left(z\right)$ est croissante donc convergente vers $u\left(z\right)$ ou divergente vers $+ \infty$. 
	Notons
	\begin{equation*}
		A = \left\{z \in U \middle|, u_{n}\left(z\right) \to + \infty\right\}, B = \left\{z \in U\middle| u_{n}\left(z\right) \to u\left(z\right)\right\}
	\end{equation*}
	Soit $a \in U$ et $\overline{D}\left(a, R\right) \subseteq U$. Pour $z \in \overline{D}\left(a, R/2\right)$, $\frac{1}{3}u_{n}\left(a\right) \leq u_{n}\left(z\right) \leq 3u_{n}\left(a\right)$. 
	Donc $A$ et $B$ sont ouverts, $A \cap B = \emptyset$, $A \cup B = U$ connexe donc $A = U$ ou $B = U$.\\
	Si $a \in B$, soit $\varepsilon > 0$, il existe $N$ tel que 
	\begin{equation*}
		0 \leq u_{n}\left(a\right) - u_{m}\left(a\right) < \frac{\varepsilon}{3}, n \geq m \geq N
	\end{equation*}
	Donc, 
	\begin{equation*}
		0 \leq u\left(z\right) - u_{m}\left(z\right) \leq \varepsilon, m \geq N, z \in D\left(a, R/2\right)
	\end{equation*}
	donc la convergence est uniforme sur tout compact de $U$. \\
	Ainsi, $u$ est harmonique par théorème de dérivation des limites. 
\end{proof}

\subsection{Couronnes du Plan}
\begin{lemme}
	{Biholomorphisme de Couronnes}{biholo-couronne}
	Soit $R_{1}, R_{2} > 1$ et $\phi: A_{1} \to A_{2}$ un biholomorphisme entre les deux couronnes
	\begin{equation*}
		A_{1} = \left\{z \in \C \middle| 1< \abs{z} < R_{1}\right\} \text{ et } A_{2} = \left\{ z \in \C \middle| 1 < \abs{z} < R_{2}\right\}
	\end{equation*}
	Pour toutes suites $\omega_{j}, \omega_{j}'$ de $A_{1}$ qui convergent respectivement vers $\omega, \omega' \in \partial A_{1}$, avec $\abs{\omega} = 1$ et $\abs{\omega'} = R_{1}$, on a l'une des deux propositions suivantes:
	\begin{enumerate}
		\item $\abs{\phi\left(w_{j}\right)} \to 1$ et $\abs{\phi\left(w_{j}'\right)} \to R_{2}$
		\item $\abs{\phi\left(w_{j}\right)} \to R_{2}$ et $\abs{\phi\left(w_{j}'\right)} \to 1$
	\end{enumerate}
\end{lemme}

\begin{théorème}
	{Biholomorphisme de Couronnes}{biholo-couronnes}
	Soit $R_{1}, R_{2} > 1$. Les couronnes $A_{1} = C\left(0, 1, R_{1}\right)$ et $A_{2} = C\left(0, 1, R_{2}\right)$ sont biholomorphes si et seulement si $R_{1} = R_{2}$.
\end{théorème}

\part{Fonctions Méromorphes}
\section{Théorème des Résidus}
\subsection{Développement en Séries de Laurent}
\begin{définition}
	{Série de Laurent}{}
	Une série de Laurent est une série de la forme \begin{equation*}
		\sum_{n\in \Z}\left(z - a\right)^{n}
	\end{equation*}
	Pour qu'elle est un sens, il faut que $\sum_{n \geq 0}a_{n}\left(z - a\right)^{n}$ ait un rayon de convergence $R_{2} > 0$ et que $\sum_{n \leq 0} a_{n}\left(z - a\right)^{n}$ ait un rayon de convergence $\frac{1}{R_{1}}$ où on a \begin{equation*}
		0 \leq R_{1} < R_{2} \leq \infty
	\end{equation*}
\end{définition}

\begin{théorème}
	{Séries de Laurent et Fonctions Holomorphes}{laurent-holomorphe}
	Si $f \in \O\left(U\right)$, si $V = \left\{z \mid R_{1} < \abs{z - a} < R_{2}\right\}$ est une couronne inclue dans $U$, $f\left(z\right) = \sum_{n \in \Z} a_{n}\left(z - a\right)^{n}$ sur $V$, les coefficients $a_{n}$ étant uniques. 
\end{théorème}
\begin{proof}
	\begin{description}
		\item[Unicité] Pour $f = 0 = \sum_{n \in \Z} a_{n}\left(z - a\right)^{n}$. 
			Mais, si $R_{1}  < r < R_{2}$
			\begin{equation*}
				\int_{\partial D\left(0, r\right)}z^{n}\d z = \begin{cases}
					0 & n \neq -1\\ 2\pi & n = - 1
				\end{cases}
			\end{equation*}
			Comme de plus: 
			\begin{equation*}
				\int_{\partial D\left(0, r\right)}f\left(\omega\right)\omega^{- p - 1}\d \omega = a_{p}2\pi = 0
			\end{equation*}
			On a bien, pour tout $p$, $a_{p} = 0$.\\
			Donc on a bien l'unicité du développement en prenant la différence des deux développements.
		\item[Existence] Pour $R_{1} < r_1 < \abs{z - a} < r_{2} < R_{2}$, on a: 
			\begin{equation*}
				2i\pi f\left(z\right) = \int_{\partial D\left(a, r_{2}\right)}\frac{f\left(u\right)}{u - z}\d u - \int_{\partial D\left(a, r_{1}\right)} \frac{f\left(u\right)}{u - z} \d u
			\end{equation*}
			\begin{itemize}
				\item Sur $\partial D\left(a, r_{2}\right)$, \begin{equation*}
						\begin{aligned}
							\frac{1}{u - z} =& \frac{1}{u - a - \left(z - a\right)} = \sum_{n \geq 0}\frac{\left(z - a\right)^{n}}{\left(u - a\right)^{n + 1}}\\
								\int_{\partial D\left(a, r_{2}\right)} \frac{f\left(z\right)}{u - z} \d z = \sum_{n \geq 0} \left(z - a\right)^{n}\int_{\partial D\left(a, r_{2}\right)} \frac{f\left(u\right)}{\left(u-a\right)^{n + 1}}\d u
						\end{aligned}
					\end{equation*}
				\item De même, sur $\partial D\left(a, r_{1}\right)$, on a:
					\begin{equation*}
						\begin{aligned}
							- \int_{partial D\left(a, r_{1}\right)} \frac{f\left(u\right)}{u - z} \d u =& - \sum_{n \geq 0} \frac{1}{\left(z - a\right)^{n + 1}} \int_{\partial D\left(a, r_{1}\right)} f\left(u\right) \left(u - a\right)^{n} \d u\\
								=& - \sum_{n \leq -1} \left(z - a\right)^{n} \int_{\partial D\left(a, r_{1}\right)} \frac{f\left(u\right)}{\left(u - a\right)^{n + 1}}\d u
						\end{aligned}
					\end{equation*}
			\end{itemize}
			Finalement, on a:
			\begin{equation*}
				f\left(z\right) = \frac{1}{2i \pi}\left(\sum_{n \geq 0} \left(z - a\right)^{n}\int_{\partial D\left(a, r_{2}\right)}\frac{f\left(u\right)}{\left(u - a\right)^{n + 1}}\d u - \sum_{n \leq - 1}\left(z - a\right)^{n}\int_{\partial D\left(a, r_{1}\right)}\frac{f\left(u\right)}{\left(u- a\right)^{n + 1}}\d u\right)
			\end{equation*}
	\end{description}
\end{proof}

\begin{définition}
	{Singularité}{}
	Soit $f : U \setminus \{a\} \to \C$ une fonction holomorphe. Le point $a$ est dit singularité de $f$. 
	Il existe $r > 0$ tel que $f\left(z\right) = \sum_{n \in \Z}a_{n} \left(z- a\right)^{n}$ pour $0 < \abs{z} < r$.
	\begin{enumerate}
		\item Si $a_{n} = 0, n < 0$, la singularité en $a$ est dite \emph{singularité apparente} et $f$ se prolonge en une fonction holomorphe sur $U$. 
		\item S'il existe $k < 0$ tel que $a_{n} = 0, n < k$, le point $a$ est dit \emph{pôle} d'ordre $k$ ou \emph{singularité isolée} et la fonction $\left(z - a\right)^{k}f\left(z\right)$ est holomorphe sur $U$. 
		\item Si $a_{n} \neq 0$ pour une infinité de $n < 0$, alors, la singularité en $a$ est dite \emph{singularité essentielle}.
	\end{enumerate}
\end{définition}

\subsection{Théorème des Résidus}

\newcommand{\Res}{\mathrm{Res}}

\begin{définition}
	{Résidu d'une fonction}{}
	Soit $f : \left\{z \middle| R_{1}  < \abs{z - a} < R_{2}\right\} \longrightarrow \C$ une fonction holomorphe.
	Le résidu de $f$ au point $a$ est le coefficient $a_{- 1}$ de son développement de Laurent
	\begin{equation*}
		\Res\left(f, a\right) = \frac{1}{2i\pi}\int_{\gamma}f\left(u\right) \d u
	\end{equation*}
	avec $\gamma = \partial D\left(a, r\right)$ pour $R_{1} < r < R_{2}$.
\end{définition}


\begin{propositionfr}
	{Exemple Très Utile}{}
	\begin{enumerate}
		\item Si $f\left(z\right) = \frac{f_{1}\left(z\right)}{\left(z - a\right)^{m}}$ où $f_{1} \in \O\left(V_{a}\right), f_{1}\left(a\right) \neq 0$, 
			alors $\Res\left(f, a\right)$ est le coefficient $\left(z - a\right)^{m - 1}$ du développement en série de Taylor de $f_{1}$, 
			$\Res\left(f, a\right) = \frac{f_{1}^{\left(m - 1\right)}\left(a\right)}{\left(m - 1\right)!}$.
			Pour un pôle simple, c'est simplement $f_{1}\left(a\right)$.\\
			Par exemple, $f\left(z\right) = \frac{z}{\left(z^{2} + 1\right)^{2}}$ s'écrit $\frac{f_{1}\left(z\right)}{\left(z - i\right)^{2}}$ avec $f_{1}\left(z\right) = \frac{z}{\left(z + i\right)^{2}}$. On a donc $\Res\left(f, i\right) = f_{1}'\left(i\right) = 0$.
		\item Si $f = \frac{g}{h}$, avec $g\left(a\right) \neq 0$ et $h$ a une racine simple en $a$, 
			alors $\Res\left(f, a\right) = \frac{g\left(a\right)}{h'\left(a\right)}$.
	\end{enumerate}
\end{propositionfr}

\begin{théorème}
	{des Résidus}{residus}
	Soit $U$ un ouvert simplement connexe de $\C$, $n \in \N^{*}$ et $a_{i} \in U^{n}$ $n$ points distincts de $U$.
	Soit $f$ holomorphe sur $U \setminus \{a_{1}, \ldots, a_{n}\}$ et $\gamma$ un lacet de $U$ évitéant les $a_{i}$.
	Alors, 
	\begin{equation*}
		\int_{\gamma} f\left(z\right)\d z = 2i\pi \sum_{i = 1}^{n} \Res\left(f, a_{i}\right)j\left(a_{i}, \gamma\right)
	\end{equation*}
\end{théorème}
\begin{proof}
	La fonction $f$ admet un développement en série de Laurent dans un voisinage de chacun des points $a_{k}$: $f\left(z\right) = \sum_{l \in \Z}b_{l}^{\left(k\right)}\left(z - a_{k}\right)^{l}$. Posons
	\begin{equation*}
		u_{k}\left(z\right) = \sum_{l = 1}^{\infty}b_{- l}^{\left(k\right)}\frac{1}{\left(z - a_{k}\right)^{l}}, 1 \leq k \leq n 
		\text{ et } g\left(z\right) = f\left(z\right) - u_{1}\left(z\right) - \cdots - u_{n}\left(z\right)
	\end{equation*}
	La fonction $g$ est holomorphe sur $U \setminus \{a_{1}, \ldots, a_{n}\}$ et les $a_{i}$ sont des singularités apparentes de $g$.
	En effet, si $D\left(a_{k}, r\right) \subseteq U\setminus \{a_{1}, \ldots, a_{k - 1}, a_{k + 1}, \ldots, a_{n}\}$, sur $D\left(a_{k}, r\right) \setminus \{a_{k}\}$:
	\begin{equation*}
		g\left(z\right) = \left(f\left(z\right) - u_{k}\left(z\right)\right) - \sum_{j \neq k} u_{j}\left(z\right)
	\end{equation*}
	est prolongeable en une fonction holomorphe sur $D\left(a_{k}, r\right)$, donc $g$ est prolongeable en une fonction holomorphe sur $U$. 
	Comme $U$ est simplement connexe, le théorème de Cauchy \ref{thm:cauchy} sur le lacet $\gamma$ de $U$ implique \begin{equation*}
		\int_{\gamma}g\left(z\right) \d z = 0
	\end{equation*}
	Mais, pour $1 \leq k \leq n, \int_{\gamma}u_{k}\left(z\right)\d z = 2i\pi \Res\left(u_{k}, a_{k}\right)j\left(a_{k}, \gamma\right)$. D'où
	\begin{equation*}
		\int_{\gamma}f\left(z\right) \d z = 2i\pi\sum_{i = 1}^{n}\Res\left(f, a_{i}\right)j\left(a_{i}, \gamma\right)
	\end{equation*}
\end{proof}

\begin{propositionfr}
	{Quelques Applications - 1}{}
	\begin{description}
        
		\item[Intégrale Trigonométriques] Soit $a \in \left]1, + \infty\right[$, 
			\begin{equation*}
				I = \int_{0}^{2\pi} \frac{\d t}{a + \sin t}
			\end{equation*}
			En posant $z = e^{it}$, on a:
			\begin{equation*}
				I = \int_{\abs{z} = 1}\frac{2 \d z}{z^{2} - 1 + 2aiz}
			\end{equation*}
			Les pôles sont en $z_{\pm} = - ai \pm i\sqrt{a^{2} - 1}$ et seul $\abs{z_{+}} < 1$.
			Ainsi, $j\left(z_{+}, \partial D\left(0, 1\right)\right) = 1$ et $j\left(z_{-}, \partial D\left(0, 1\right)\right) = 0$.
			Par le théorème des résidus \ref{thm:residus}, on obtient
			\begin{equation*}
				I = 2i\pi Res\left(f, z_{+}\right) = 2i\pi\left(\frac{2}{\left(z^{2} - 1 + 2aiz\right)'}\right)\left(z_{+}\right) = 2i\pi \frac{2}{2z_{+} + 2ai} = \frac{2\pi}{\sqrt{a^{2} - 1}}
			\end{equation*}

		\item[Intégrale d'une Fonction Rationnelle] Soit $P, Q \in \R\left[X\right]$ avec $\deg Q \geq \deg P + 2$ et $Q$ sans racine réelle. Soit 
			\begin{equation*}
				f = \frac{P}{Q} \text{ et } I = \int_{- \infty}^{+ \infty} f\left(x\right) \d x
			\end{equation*}
		Les hypothèses garantissent que $I$ converge et \begin{equation*}I = \lim_{R \to +\infty}\int_{- R}^{R}f\left(x\right) \d x\end{equation*}
		Considérons la fonction de la variable complexe $f$ et le contour $\gamma_{R} = \left[-R, R\right] \cup S_{R}$ où $S_{R}$ désigne le demi-cercle de centre $0$ de rayon $R$ sur le demi-plan $\Im z \geq 0$ parcouru dans le sens trigonométrique. 
		Par théorème des Résidus \ref{thm:residus}, on a:
		\begin{equation*}
			\int_{\gamma_{R}} f\left(z\right) \d z = \int_{- R}^{R} f\left(x\right) \d x + \int_{S_{R}} f\left(z\right) \d z = 2i\pi \sum_{a_{j}, \abs{a_{j} < R}, \Im\left(a_{j}\right) > 0} \Res\left(f, a_{j}\right)
		\end{equation*}
		où les $a_{j}$ sont les racines de $Q$. Posons \begin{equation*}\alpha = \lim_{\abs{x} \to + \infty} x^{2}\abs{f\left(x\right)}\end{equation*}
		Quand $R \to \infty$, $\abs{\int_{S_{R}}f \d z}\leq \pi R\frac{\alpha + 1}{R^{2}}$ tend vers $0$.
		Donc, 
		\begin{equation*}
			I = 2i\pi \sum_{a_{j}, \Im a_{j} > 0} \Res\left(\frac{P}{Q}, a_{j}\right)
		\end{equation*}
		Par exemple, 
		\begin{equation*}
			\begin{aligned}
				\int_{- \infty}^{+\infty}\frac{\d x}{x^{4} + 1} =& 2i\pi \left(\Res\left(\frac{1}{z^{4} + 1}, e^{i\pi/4}\right) + \Res\left(\frac{1}{z^{4} + 1}, e^{3i\pi/4}\right)\right)\\
				=& 2i\pi \left(e^{- 3i\pi/4}/4 + e^{-i\pi /4}\right)\\
				\int_{-\infty}^{+\infty}\frac{\d x}{x^{4}+1} =& \frac{\pi}{\sqrt{2}}
			\end{aligned}
		\end{equation*}
	\end{description}
\end{propositionfr}

\begin{propositionfr}
	{Quelques Applications - 2}{}
	\begin{description}
		\item[Intégrales de Fourier] On calcule la transformée de Fourier pour $f$ une fonction holomorphe sur $\C \setminus \{a_{1}, \ldots, a_{n}\}$ sans singularité sur l'axe réel avec $\lim\limits_{\abs{z}\to \infty, \Im\left(z\right) \geq 0} f\left(z\right) = 0$ et telle que l'intégrale soit absolument convergente.\\
			On considère $\gamma_{R}$ comme pour l'intégrale d'une fonction rationnelle.
			On montre de même que 
			\begin{equation*}
				\int_{S_{R}} e^{iz}f\left(z\right)\d z \xrightarrow[R \to \infty]{} 0
			\end{equation*}
			On obtient alors $I = 2i\pi \sum_{a_{j}, \Im\left(a_{j}\right) > 0}\Res\left(f\left(z\right)e^{iz}, a_{j}\right)$.
            
		\item[Intégrale Logarithmique] On calcule 
            \begin{equation*}
                I = \int_{0}^{+ \infty} f\left(x\right)\log\left(x\right)\d x
            \end{equation*} où $f \in \R\left(X\right)$ de degré $\leq -2$.
			On introduit le contour $\gamma_{r, R} = C_{r} \cup L_{-} \cup C_{r} \cup L_{+}$ où $C_{r}, C_{R}$ sont les cercles de centre $0$ et de rayon $r$ (resp. $R$) privés d'un petit arc de cercle pour éviter le demi arc réel positif.
			Les notations $L_{-}$ et $L_{+}$ désignent 
			\begin{equation*}
				L_{+} = \left\{z = t + i\epsilon \mid r \leq t \leq R \right\} \text{ et } L_{-} = \mathrm{conj}\left(L_{+}\right)
			\end{equation*}
			pour $\varepsilon$ que l'on fait tendre vers $0$. 
			En calculant avec le théorème des résidus quand $r \to 0$ et $R \to \infty$, on trouve :
			\begin{equation*}
				\int_{0}^{+\infty} \ln\left(x\right)f\left(x\right)\d x = -\frac{1}{2}\sum_{a}\Res\left(\left(\log z - i\pi\right)^{2}f\left(z\right), a\right)
			\end{equation*}
			où la somme porte sur les zéros de $Q$. 

		\item[Transformation de Mellin] On calcule, pour $\alpha \in \left]0, 1\right[$ l'intégrale
			\begin{equation*}
				I = \int_{0}^{+\infty} \frac{F\left(x\right)}{x^{\alpha}}\d x
			\end{equation*}
			On prolonge $x^{\alpha}$ par la fonction $z^{\alpha} = \exp\left(\alpha \log\left(z\right)\right)$ où $\log\left(\rho e^{it}\right) = \log \rho + it$ pour $t \in \left]0, 2\pi\right[$.
			On reprend alors le contour utilisé pour les intégrales logarithmiques. 
			On obtient alors 
			\begin{equation*}
				I\left(1 - e^{-2i\pi \alpha}\right) = 2i\pi\sum_{a} \Res\left(\frac{F\left(z\right)}{z^{\alpha}}, a\right)
			\end{equation*}
			pour $a$ les singularités de $F$. 
	\end{description}
\end{propositionfr}

\subsection{Résidus à l'Infini}
\begin{définition}
	{Singularité Isolée}{}
	La fonction holomorphe $f : \left\{z \in \C\mid \abs{z} \geq R\right\} \to \C$ admet une singularité isolée en $\infty$ si la fonction 
	\begin{equation*}g:
		\begin{array}{rcl}
			D\left(0, \frac{1}{R}\right) & \longrightarrow & \C\\
			z & \longmapsto & f\left(\frac{1}{z}\right)
		\end{array}
	\end{equation*}
	admet $0$ comme singularité isolée.
\end{définition}

\begin{définition}
	{Résidu à l'Infini}{}
	Soit $f$ holomorphe admettant une singularité isolée en $\infty$. Le résidu à l'infini de $f$ est défini par
	\begin{equation*}
		\Res\left(f, \infty\right) = -\frac{1}{2i\pi} \int_{\abs{z} = \rho}f\d z, \rho > R
	\end{equation*}

\end{définition}

\begin{lemme}
	{Résidu de l'Inverse}{}
	Soit $f$ holomorphe admettant une singularité isolée en $\infty$ et soit $g\left(z\right) = f\left(1/z\right)$. Alors
	\begin{equation*}
		\Res\left(f, \infty\right) = \Res\left(-g\left(\omega\right)\frac{1}{\omega^{2}}, 0\right).
	\end{equation*}
\end{lemme}
\begin{proof}
	Pour $\rho > R$, on a:
	\begin{equation*}
		\Res(f, \infty) = - \frac{1}{2i\pi}\int_{\abs{z} = \rho} f\d z = \frac{1}{2i\pi}\int_{\abs{\omega} = \frac{1}{\rho}} \frac{-g\left(\omega\right)}{\omega^{2}}\d\omega = \Res(g(\omega)\frac{1}{\omega^{2}}, 0)
	\end{equation*}
\end{proof}

\begin{corollaire}
	{Singularité Isolées}{}
	Si $f$ n'a que des singularités isolées en $a_{1}, \ldots, a_{n} \in \C$ et éventuellement en $\infty$, alors 
	\begin{equation*}
		\sum_{a\in \C\cup\{\infty\}} \Res(f, a) = 0
	\end{equation*}
\end{corollaire}
\begin{proof}
	Soit $R$ tel que $R \geq \abs{a_{i}}$. On a: 
	\begin{equation*}
		\Res(f, \infty) = -\frac{1}{2i\pi}\int_{\abs{z} = R} f\d z = -\sum_{j = 1}^{n}\Res(f, a_{j})
	\end{equation*}
	D'où le résultat.
\end{proof}

\section{Singularités}
\subsection{Étude Locale}
\begin{théorème}
	{du Prolongement de Riemann}{prolong-riemann}
	Soit $f : U\setminus \{a\} \to \C$ holomorphe. Si $f$ est bornée au voisinage de $a$ alors $a$ est une singularité apparente. 
\end{théorème}
\begin{proof}
	Les coefficients de Laurent vérifient les inégalités de Cauchy: 
	\begin{equation*}
		\abs{a_{n}} \leq \frac{\sup_{\partial D(a, r)}\abs{f(z)}}{r^{n}}
	\end{equation*}
	Donc pour $n < 0$, si $f$ est bornée au voisinage de $a$ en faisant tendre $r$ vers $0$, nous obtenons $a_{n} = 0, n < 0$.
\end{proof}

\begin{théorème}
	{de Weierstrass}{weierstrass-sing-essentielle}
	Soit $U$ ouvert, $a \in U$, $f$ holomorphe sur $U \setminus \{a\}$ avec $a$ une singularité essentielle de $f$. 
	Alors, pour tout $\epsilon > 0$, $f\left(\left\{z \middle| 0 < \abs{z - a} < \epsilon\right\}\right)$ est dense dans $\C$. 
\end{théorème}
\begin{proof}
	Raisonnons par l'absurde en supposant 
	\begin{equation*}
		\exists b \in \C, r > 0, f\left(\left\{z \middle| 0 < \abs{z - a} < \epsilon\right\}\right) \cap D(b, r) = \emptyset
	\end{equation*}
	Alors la fonction $g : D(a, \epsilon) \setminus \{a\} \to \C$ qui à $z$ associe $\frac{1}{f(z) - b}$ est holomorphe bornée au voisinage de $a$ donc admet un prolongement holomorphe sur $D(a, \epsilon)$ et un développement de la forme $g(z) = \sum_{n \geq 0} b_{n} \left(z - a\right)^{n}$. 
	Notons $b_{k}$ le premier coefficient non nul de $g$. On a: 
	\begin{equation*}
		f(z) = b + \frac{1}{g(z)} = b + \frac{1}{b_{k}(z - a)^{k}\left(1 + \sum_{n \geq k + 1} \frac{b_{n}}{b_{k}}(z - a)^{n-k}\right)}
	\end{equation*}
	Donc $f$ admet un développement de Laurent avec un pôle d'ordre $k$ en $a$, ce qui est absurde.
\end{proof}

\subsection{Grand Théorème de Picard} 
\begin{lemme}
	{Domination Holomorphe}{domi-holo}
	Soit $z_{0} \in \C$, $R > 0$. Il existe une fonction
	\begin{equation*}
		M: \left[1, +\infty\right[ \times \left]0, R\right[ \to \R^{+}
	\end{equation*}
	croissante en les variables $A, r$ telle que toute application holomorphe $f: D(z_{0}, R) \to \C\setminus\{0, 1\}$ vérifiant $A^{-1}\leq \abs{f(z_{0})} \leq A$ satisfait
	\begin{equation*}
		\abs{f(z)} \leq M(A, r), z \in \overline{D}(z_{0}, r)
	\end{equation*}
\end{lemme}
\begin{proof}
	Comme $D(z_{0}, R)$ est simplement connexe et $f(D(z_{0}, R)) \subseteq \C^{*}$, il existe une détermination du logarithme $f_{1} = \frac{1}{2i\pi}\log f$ sur $D(z_{0}, R)$ telle que $0 \leq \Re f_{1}(z_{0}) \leq 1$. 
	Nous avons $\abs{\Im f_{1}(z_{0})} \leq \frac{1}{2\pi}\log A$, donc
	\begin{equation*}
		\abs{f_{1}(z_{0}) - j} \leq B = \sqrt{1 + \frac{1}{4\pi^{2}}\left(\log A\right)^{2}}, j = 0, 1
	\end{equation*}
	Comme $1 \notin f(D(z_{0}, R))$, $f_{1}(D(z_{0}, R)) \subseteq \C \setminus \Z$. 
	Il existe donc des fonctions $f_{2}, f_{3}$ holomorphes telles
	\begin{equation*}
		f_{2} = \sqrt{f_{1}} + \sqrt{f_{1} - 1}, f_{3} = \sqrt{f_{1}} - \sqrt{f_{1}-1}
	\end{equation*}
	Nous avons $f_{2}f_{3} = 1$, $\abs{f_{2}}\leq 2\sqrt{B}$ et $\abs{f_{3}}\leq 2\sqrt{B}$. Donc
	\begin{equation*}
		C(A)^{-1} \leq \abs{f_{2}(z_{0})} \leq C(A), C(A) = 2\sqrt{B} = 2\left(1 + \frac{1}{4\pi^{2}}\left(\ln A\right)^{2}\right)^{1/4}
	\end{equation*}
	Or $f_{2}$ n'atteint jamais $\sqrt{n} + \sqrt{n - 1}$ ni $\sqrt{n} - \sqrt{n - 1}$, $n \in \N^{*}$. 
	Comme $f_{2}$ ne s'annule pas et $D(z_{0}, R)$ est simplement connexe, il existe une détermination holomorphe $f_{4} = \log f_{2}$.
	Alors $f_{4}$ omet toutes les valeurs de l'ensemble 
	\begin{equation*}
		S = \left\{\pm \log\left(\sqrt{n} + \sqrt{n - 1}\right) + 2i\pi k\middle| n \in \N^{*}, k \in \Z\right\}
	\end{equation*}
	On peut choisir un argument $-\pi < \Im f_{4}(z_{0}) \leq \pi$ de telle sorte que $\abs{f_{4}(z_{0})} \leq\sqrt{\left(\ln C(A)\right)^{2} + \pi^{2}}$.
	Les nombres $\pm \ln\left(\sqrt{n} + \sqrt{n - 1}\right)$ sont de moins en moins espacés quand $n$ augmente, ce qui entraîne que le plus grand disque contenu dans $\C \setminus S$ a pour rayon
	\begin{equation*}
		R_{0} = \sqrt{\frac{1}{4}\left(\ln(\sqrt{2} + 1)\right)^{2} + \pi^{2}}
	\end{equation*}
	Pour $z \in D(z_{0}, R)$, le théorème de Bloch-Landau \ref{thm:bloch-landau} applliqué au disque $D(z, \rho)$ avec $\rho = R - \abs{z - z_{0}}$ montre que l'image de $f_{4}$ contient un disque de rayon supérieur à $\frac{1}{12}\rho\abs{f_{4}'(z)}$. 
	Or, $f_{4}\left(D(z_{0}, R)\right) \subseteq \C \setminus S$. Ainsi
	\begin{equation*}
		\abs{f_{4}'(z)}\leq 12\frac{R_{0}}{R - \abs{z - z_{0}}}
	\end{equation*}
	Or, $f_{4}(z) = f_{4}(z_{0}) + \int_{z_{0}}^{z}f_{4}'(t)\d t$ donc
	\begin{equation*}
		\abs{f_{4}(z)} \leq \sqrt{\left(\ln C(A)\right)^{2} + \pi^{2}} + 12R_{0}\ln\frac{R}{R - r} = T(A, r)
	\end{equation*}
	sur le disque $D(z_{0}, r)$.
	Ceci entraîne l'existence de majorations uniformes explicites pour $f_{2} = \exp f_{4}$ et $f_{2}^{-1} = \exp -f_{4}$, $\sqrt{f_{1}} = \frac{1}{2}\left(f_{2} + f_{2}^{-1}\right)$ et $f = \exp\left(2i\pi f_{1}\right)$:
	\begin{equation*}
		\begin{aligned}
			\abs{f_{2}(z)}^{\pm 1} \leq & \exp\left(T(A, r)\right)\\
			\abs{f_{1}(z)} \leq & \exp\left(2T(A, r)\right)\\
			\abs{f(z)} \leq& M(A, r) = \exp\left(2\pi \exp\left(2T(A, r)\right)\right)
		\end{aligned}
	\end{equation*}
\end{proof}

\begin{théorème}
	{Grand Théorème de Picard}{grand-picard}
	Soit $f \in \O(D(z_{0}, r_{0})\setminus \{z_{0}\})$ holomorphe avec une singularité essentielle en $z_{0}$. 
	Pour tout voisinage $V$ de $z_{0}$, $f(z)$ prend sur $V \setminus \{z_{0}\}$ une infinité de fois toutes les valeurs de $\C$ sauf au plus une. 
\end{théorème}
\begin{proof}
	On peut supposer $z_{0} = 0$ par translation. Soit $V \subseteq D(0, r_{0})$ un voisinage de $0$. 
	Supposons que $f(V \setminus \{0\}) \subseteq \C \setminus \{a, b\}$. 
	Quitte à remplacer $f$ par $\frac{\left(f(z) - a\right)}{b - a}$, on peut supposer $a = 0, b = 1$. 
	Pour $r \in ]0, r_{0}[$, posons
		\begin{equation*}
			\mu(r) = \inf_{\abs{z} = r}\abs{f(z)}, m(r) = \sup_{\abs{z} = r}\abs{f(z)}
		\end{equation*}
	Montrons par l'absurde que $\lim_{r \to 0} m(r) = +\infty$.
	Sinon, il existe une suite $r_{k} \to 0$ telle que $m(r_{k})$ est bornée. 
	Le principe du maximum \ref{thm:principe_max} sur les couronnes $\left\{z \middle| r_{k + 1} \leq \abs{z}\leq r_{k}\right\}$ implique que $f$ est bornée au voisinage de $z_{0} = 0$ ce qui est exclu car $f$ admet une singularité essentielle en $0$.
	Par le même raisonnement appliqué à $1/f$, $\lim_{r \to 0} \mu(r) = 0$.\\
	Pour $r$ assez petit, $\mu(r) < 2 < m(r)$ et le théorème des valeurs intermédiaires garantit l'existence de $\theta_{r} \in \left[0, 2\pi\right[$ tel que $\abs{f\left(re^{i\theta_{r}}\right)} = 2$.
	Pour $0 < r < r_{0}e^{-2\pi}$, définissons 
\begin{equation*}
	g_{r}: z \longmapsto f\left(re^{i\theta_{r}}e^{2i\pi z}\right)
\end{equation*}
	Alors $g_{r}\left(D(0, 1)\right) \subseteq \C\setminus \{0, 1\}$, $\abs{g_{r}(0)} = \abs{f\left(re^{i\theta_{r}}\right)} = 2$. 
	Le lemme précédent \ref{lemme:domi-holo} garantit l'existence d'une constance $M(2, 1/2)$ telle que les fonctions $g_{r}$ soient uniformément bornées par $M(2, 1/2)$ sur $\overline{D}(0, 1/2)$. 
	Or $g_{r}\left(\left[-1/2, 1/2\right]\right)$ décrit toutes les valeurs atteintes par $f$ pour $\abs{z} = r$.
	Donc $m(r) \leq M(2, 1/2)$ pour tout $r < r_{0}e^{-2\pi}$ ce qui est absurde.\\
	Par conséquent, pour tout voisinage $V \setminus \{0\}$ de $0$, $f(V \setminus \{0\})$ atteint toute valeur de $\C$ sauf au plus une. 
	En réduisant la taille des voisinages, $f(V \setminus \{0\})$ prend une infinité de fois toutes valeurs de $\C$ sauf au plus une. 
\end{proof}

\begin{corollaire}
	{Petit Théorème de Picard}{petit-picard}
	Soit $f \in \O(\C)$ non constante. Alors $f$ prend toutes les valeurs complexes sauf éventuellement une. 
\end{corollaire}
\begin{proof}
	Les polynômes non constants prennent toutes les valeurs complexes, mais seulement un nombre fini de fois. 
	Supposons que $f \in \O(\C)$ n'est pas un polynôme et notons $f(z) = \sum a_{n}z^{n}$. 
	La fonction $g(z) = f(1/z)$ est holomorphe sur $\C^{*}$ avec 
	\begin{equation*}
		g(z) = \sum_{n = 0}^{\infty} a_{n}z^{-n}
	\end{equation*}
	Comme $f$ n'est pas un polynôme, $g$ admet une singularité essentielle en $0$. 
	Donc $g$ et $f$ atteignent une infinité de fois toutes les valeurs compelxes sauf éventuellement une. 
\end{proof}

\begin{corollaire}
	{Points Fixes}{no-translate-point-fix}
	Soit $f$ holomorphe qui n'est pas une translation. Alors $f \circ f$ admet un point fixe.
\end{corollaire}
\begin{proof}
	Par l'absurde, si $f \circ f$ n'a pas de point fixe, $f$ n'a pas de point fixe et 
	\begin{equation*}
		g(z) = \frac{f(f(z)) - z}{f(z) - z} \in \O(\C)
	\end{equation*}
	évite les valeurs $0$ et $1$. 
	Par le petit théorème de Picard \ref{cor:petit-picard}, $g = c \neq 1$ est constante et 
	\begin{equation*}
		f(f(z)) - z = c(f(z) - z), z \in \C
	\end{equation*}
	Par conséquence $f'(z)(f'(z)(f(z)) - c) = 1 - c$. 
	Comme $c \neq 1$, la fonction holomorphe $f'\circ f$ ne prend pas la valeur $c$. S'il existe $z_{0} \in \C$
	avec $f'\circ f(z_{0} = 0$ alors pour $z = f(z_{0})$, on obtient $c = 1$ ce qui est exclu. 
		Donc $f' \circ f \in \O(\C)$ évite les valeurs $c$ et $1$ donc est constante par le petit théorème de Picard. 
	Donc $f'$ est constante et donc $f = az + b$. Comme $f$ n'est pas une translation, $a \neq 1$ et donc $f$ a un point fixe. 
\end{proof}


\section{Théorème de Rouché}
\subsection{Fonctions Méromorphes}
\begin{définition}
	{Corps des Fractions Méromorphes}{}
	Soit $U$ connexe. L'anneau $\O(U)$ est intègre. Nous notons $\mathcal{F}(U)$ son corps des fractions. 
\end{définition}

\begin{définition}
	{Fonction Méromorphe}{}
	Une fonction $f: U \to \C\cup\{\infty\}$ est dite méromorphe si $f$ est holomorphe sur un ouvert $U'$ de $U$ avec $U\setminus U'$ discret et si tout point $z_{0}$ de $U$ admet un voisinage ouvert connexe dans $U$ tel que $f_{\mid V_{z_{0}}} \in \mathcal{F}\left(V_{z_{0}}\right)$ 
\end{définition}

\begin{propositionfr}
	{Propriétés des fonctions Méromorphes}{}
	\begin{itemize}
	\item L'ensemble des fonctions méromorphes sur $U$ est une $\C$-algèbre. Si $U$ est connexe, $\mathcal{M}(U)$ est un corps. 
	\item Soit $f$ méromorphe au voisinage de $z_{0}$. Alors, il existe une fonction holomorphe $g$ au voisinage de $z_{0}$ telle que $f(z) = (z - z_{0})^{m}g(z)$ avec $g(z_{0}) \neq 0$. Si $m \neq 0$ la fonction $\log f(z)$ n'est pas définie au voisinage de $z_{0}$ mais sa dérivée logarithmique l'est.
	\item Une fonction $f: U \to \C\cup \{\infty\}$ est méromorphe si et seulement si elle holomorphe en dehors d'une suite de points singuliers localement finie dans $U$ et si tous ces points singuliers sont des pôles. 
		Les fonctions rationnelles sont méromorphes sur $\C$. La fonction $\left(\sin(1/\pi z)\right)^{-1}$ est méromorphe sur $\C^{*}$ mais pas sur $\C$. 
	\end{itemize}
\end{propositionfr}

\subsection{Principe de l'Argument}
\begin{définition}
	{Ordre}{}
	Soit $U$ un ouvert de $\C$, $a\in U$. Soit $f$ méromorphe sur $U$. On définit l'ordre de $f$ en $a$ par
	\begin{equation*}
			\mathrm{ord}(f, a) = 
			\begin{cases}
		n & \text{si } a \text{ est un zéro d'ordre } n\\
		-n & \text{si } a \text{ est un pôle d'ordre } n\\
		0 & \text{sinon}
	\end{cases}
	\end{equation*}
\end{définition}

\begin{théorème}
	{Principe de l'Argument}{principe-argument}
	Soit $f$ méromorphe sur $U$ et un $\gamma$ un lacet de $U$ ne contenant ni pôle ni zéro de $f$. 
	\begin{equation*}
		\frac{1}{2i\pi} \int_{\gamma}\frac{f'}{f}\d z = \sum_{a}j(a, \gamma)\mathrm{ord}(f, a)
	\end{equation*}
	où $a$ parcour l'ensemble des zéros et pôles de $f$ à l'intérieur de $\gamma$. 
\end{théorème}
\begin{proof}
	Si $a \in U$ est un pôle ou un zéro de $f$, alors $a$ est un pôle simple de $\frac{f'}{f}$ avec un résidu égal à $\mathrm{ord}(f, a)$ car
	\begin{equation*}
		f(z) = (z - a)^{m}f_{1}(z)
	\end{equation*}
	On déduit alors le résultat du théorème des résidus \ref{thm:residus}
\end{proof}

\begin{propositionfr}
	{Quelques Corollaires}{}
	\begin{itemize}
		\item Si $\gamma$ est une courbe simple fermée : 
			\begin{equation*}
				\frac{1}{2i\pi}\int_{\gamma}\frac{f'}{f} \d z = \text{\#zéros}f - \text{\#pôles}f
			\end{equation*}
		\item On a: 
			\begin{equation*}
				\Re\left(\frac{1}{2i\pi}\frac{f'}{f}\right) = \Re\left(\frac{1}{2i\pi}\frac{\d \log(f)}{\d z}\right) = \Re\left(\frac{1}{2i\pi}\frac{\d}{\d z}\left(\log \abs{f} + i\arg f\right)\right) = \frac{1}{2\pi}\frac{\d}{\d z}\arg f
			\end{equation*}
		\item Si $g$ est holomorphe:
			\begin{equation*}
				\frac{1}{2i\pi} \int_{\gamma}g\frac{f'}{f}\d z = \sum_{a} g(a)j(a, \gamma)\mathrm{ord}(f, a)
			\end{equation*}
	\end{itemize}
\end{propositionfr}

\begin{propositionfr}
	{Injectivité sur un Voisinage}{}
	Soit $f$ holomorphe sur $U$, $\overline{\Delta}\subseteq U$ tel que $f_{\mid \overline{\Delta}}$ est injective et $f\left(\Delta\right)$ est ouvert. Alors $f: \Delta \to V$ est bijective et 
	\begin{equation*}
		f^{-1}(q) = \int_{\partial \Delta}z\frac{f'(z)}{f(z) - q}\d z 
	\end{equation*}
	En particulier, $f^{-1}$ est holomorphe. 
\end{propositionfr}

\subsection{Fonctions Elliptiques}

Dans la suite, on pose $\omega_{1}, \omega_{2} \in \C\setminus \{0\}$ avec $\frac{\omega_{1}}{\omega_{2}} \notin \R$et on définit le réseau $\Lambda = \Z\omega_{1} + \Z\omega_{2}$

\begin{définition}
	{Fonction Elliptique}{}
	Une fonction $f$ est dite \emph{elliptique} si elle est méromorphe sur $\C$ et si $f$ est $\Lambda$-périodique: 
	\begin{equation*}
		f(z) = f(z + \omega_{1}) = f(z + \omega_{2}), \forall z \in \C
	\end{equation*}
\end{définition}

\begin{propositionfr}
	{Propriétés}{}
	\begin{itemize}
		\item L'ensemble des fonctions elliptiques est un corps stable par dérivation. 
		\item Soit $f$ elliptique holomorphe et $P_{A}$ un parallélogram fondamental de $\Lambda$. Comme $f$ est holomorphe sur $\C$, $f_{\mid\overline{P_{A}}}$ est continue. Comme $f$ est périodique, $f$ est bornée donc constante. 
	\end{itemize}
\end{propositionfr}

On rappelle qu'un paralléllogramme est de la forme 
\begin{equation*}
	P_{a} = \left\{a + t_{1}\omega_{1} + t_{2}\omega_{2}\middle| 0 \leq t_{1} < 1, 0 \leq t_{2} < 1\right\}
\end{equation*}

\begin{lemme}
	{Sans zéros}{}
	Soit $f$ elliptique. Il existe $a \in \C$ tel que $\partial P_{a}$ ne contient ni zéro ni pôle de $f$. 
\end{lemme}
\begin{proof}
	L'ensemble des pôles et zéros est discret et $\overline{P_{a}}$ est compact donc il n'en contient qu'un nombre fini. 
\end{proof}


\subsection{Théorème de Rouché}
\begin{théorème}
	{Rouché}{rouché}
	Soit $U$ ouvert, $D(a, r) \subseteq U$, $f, g \in \mathcal{M}(U)$ sans pôle ni zéro sur $\partial D(a, r)$ et telles que $\abs{f - g} < \abs{f}$ au bord. Alors, dans $D(a, r)$: 
	\begin{equation*}
		\text{\#zéros} f - \text{\#pôles}f = \text{\#zéros}g - \text{\#pôles}g
	\end{equation*}
\end{théorème}
\begin{proof}
	On pose $h \frac{g}{f}$ sur $D(a, r)$. Soit $\gamma$ le lacet bord de $D(a,r)$ et $\Gamma = h\circ \gamma$.
	La fonction $\abs{\frac{f(\gamma(t)) - g(\gamma(t))}{f(\gamma)(t)}}$ est définie et continue dans $I$ donc atteint son maximum $M$ en un point de $I$.
	Ainsi, $\abs{h(\gamma(t)) - 1} \leq M < 1$ sur $I$. 
	Le lacet $\Gamma$ est donc inclus dans $D(0, 1)$ simplement connexe et $\Gamma$ ne contient pas $0$. 
	La fonction $\log h$ est définie sur $\gamma(I)$. Donc: 
	\begin{equation*}
		\frac{1}{2i\pi}\int_{\gamma}\frac{h'(z)}{h(z)} = 0
	\end{equation*}
	Or, $\frac{h'}{h} = \frac{g'}{g} - \frac{f'}{f}$. Donc on a l'égalité attendu par principe de l'argument \ref{thm:principe-argument}.
\end{proof}

\begin{propositionfr}
	{Quelques Applications}{}
	\begin{itemize}
		\item Pour déterminer le nombre de racines de $f(z) = z^{6} + 21z^{4} + 2z + 3$ sur $D(0, 1)$, on introduit $g(z) = 21z^{4}$. Le théorème de Rouché nous indique que $f$ admet $4$ zéros (autant que $g$). 
		\item Plus généralement, pour $f(z) = z^{n} + a_{1}z^{n - 1} + \ldots + a_{n}, g(z) = z^{n}$ et $R$ assez grand, le théorème de Rouché \ref{thm:rouché} montre que $f$ admet $n$ zéros dans $\C$, i.e. que $\C$ est algébriquement clos. 
		\item Toutes les solutions de $z\sin z = 1$ sont réelles en appliquant le théorème de Rouché à $f(z) = z\sin z$ et $g(z) = z\sin z - 1$ sur le carré $\Delta = \left\{z \middle| \abs{x} + \abs{y}\leq \pi \left(n + \frac{1}{2}\right)\right\}$.
	\end{itemize}
\end{propositionfr}

\begin{corollaire}
	{Point Fixe Méromorphe}{}
	Soit $f \in \mathcal{M}(U)$ avec $\overline{D(0, 1)} \subseteq U$, $\abs{f(z)} < 1$ pour $z \in \partial D(0, 1)$.
	Alors $f$ a un unique point fixe dans $D(0, 1)$. 
\end{corollaire}
\begin{proof}
	Posons $h(z) = f(z) - z$, $g(z) = -z$. Alors $\abs{h - g} = \abs{f} < 1 = \abs{g}$ sur $\partial D(0, 1)$. On conclut par théorème de Rouché. 
\end{proof}

\begin{théorème}
	{Hurwitz}{hurwitz}
	Soit $U$ connexe et $f_{n} \in \O(U)$ convergeant uniformément sur tout compact de $U$ vers $f$.\\
	Si $\overline{D(z_{0}, R)}\subseteq U$ et $f$ ne s'annule pas au bord de ce disque, alors, pour $n$ assez grand $f_{n}$ ne s'annule pas sur $\partial D(z_{0}, R)$ et $f_{n}$ a autant de zéros que $f$ sur ce disque. 
\end{théorème}
\begin{proof}
	On prend $n$ assez grand pour que $\abs{f_{n} - f} < \inf_{\partial K}\abs{f}$ sur $\partial D(z_{0}, R)$. 
\end{proof}

\begin{corollaire}
	{}{}
	Soit $U$ un ouvert connexe, $f_{n} \in \O(U)$ convergeant vers $f$ uniformément sur tout compact avec $f_{n}$ injective. 
	Alors $f$ est injective ou constante. 
\end{corollaire}
\begin{proof}
	Supposons qu'il existe $z_{1}, z_{2}$ tels que $f(z_{1}) = f(z_{2}) = \omega$ et que $f(z) - \omega$ a ses zéros isolés. 
	Alors $f(z) - \omega$ ne s'annule pas en dehors de $z_{1}, z_{2}$ sur $\overline{D(z_{1}, r_{1})} \cup \overline{D(z_{2}, r_{2})}$. 
	Par théorème de Hurwitz, $f_{n} - \omega$ a le même nombre de zéros que $f - \omega$ à partir d'un certain rang sur $D(z_{1}, r_{1})$ et $D(z_{2}, r_{2})$. Donc $f_{n}$ n'est pas injective, contradiction!
\end{proof}

\subsection{Théorème de l'Application Conforme de Riemann}
\begin{lemme}
	{Isomorphisme d'ouverts}{iso-ouvert-dans-compact}
	Soit $U \subsetneq \C$ simplement connexe, il existe $U'$ inclus dans un compact isomorphe à $U$
\end{lemme}
\begin{proof}
	Soit $a \in \C \setminus U$. 
	Comme $U$ est simplement connexe il existe une détermination du logarithme $g(z) = \log(z - a)$ sur $U$ qui de plus est injective. 
	Soit $z_{0} \in U$. Par théorème d'inversion globale \ref{thm:global_inv}, il existe $D(g(z_{0}), r)\subset g(U)$. 
	Comme $g$ est une détermination du logarithme, le disque $D(g(z_{0}), r) + 2i\pi$ ne contient aucun point de $g(U)$. 
		Ainsi, pour tout $z \in U$, $\abs{g(z) - g(z_{0}) - 2i\pi} \geq r > 0$. 
		La fonction 
		\begin{equation*}
			h : z \longmapsto \frac{1}{g(z)  - g(z_{0}) - 2i\pi}
		\end{equation*}
		est holomorphe, injective et bornée sur $U$.
		Par théorème d'inversion globale \ref{thm:global_inv}, $h: U \to U'$ est un isomorphisme ou $U' \subseteq \overline{D\left(0, r^{-1}\right)}$
\end{proof}

\begin{lemme}
	{Isomorphisme et Image dans le Disque Unité}{iso-disque-unité}
	Soit $U$ un ouvert simplement connexe de $\C$ contenant $0$. Soit 
	\begin{equation*}
		A = \left\{f \in \O\left(U\right)\ \middle| \ f \text{ injective}, f(0) = 0, \forall z \in U \abs{f(z)} < 1\right\}
	\end{equation*}
	Alors, pour $f \in A$, les propriétés suivantes sont équivalentes: 
	\begin{enumerate}
		\item $f(U) = D(0, 1)$
		\item $\abs{f'(0)}$ est maximum sur $A$. 
	\end{enumerate}
\end{lemme}
\begin{proof}
	\begin{description}
		\item[$i \Rightarrow ii$] D'après le théorème d'inversion globale \ref{thm:global_inv}, les éléments de $A$ sont des isomorphismes sur leur image. 
			Soit $f \in A$ d'image $D(0, 1)$. Alors $h = g \circ f^{-1} : D(0, 1) \to g(U)$ est un isomorphisme tel que $h(0) = 0$. 
			Par le lemme de Schwarz \ref{thm:lemme_schwarz}, $\abs{h'(0)} \leq 1$ donc $\abs{g'(0)} \leq \abs{f'(0)}$ et $\abs{f'(0)}$ est maximum. 
		\item[$ii \Rightarrow i$] On montre $\lnot i \Rightarrow \lnot ii$. Soit $f \in A$ et $a \in D(0, 1)$ avec $a \notin f(U)$. 
			Soit $\phi_{a} \in \mathrm{Aut} D(0, 1): z\mapsto \frac{z - a}{1 - \bar{a}z}$. 
			Comme $f(U)$ est simplement connexe, $F(z) = \log \phi_{a}f(z)$ est définie holomorphe et injective dans $U$ et $\Re F(z) < 0$. 
			Pour $u, v \in \C$ avec $\Re u, \Re v < 0$, on a: $\abs{\frac{v - u}{v + \bar{u}}} < 1$.\\
			Ainsi, $g(z) = \frac{F(z) - F(0)}{F(z) + \bar{F(0)}}$ est holomorphe injective dans $U$, $g(0) = 0$, $\abs{g(z)} < 1$ et donc $g \in A$. 
			Enfin
			\begin{equation*}
				\abs{\frac{g'(0)}{f'(0)}} = \frac{1 - \abs{a}^{2}}{2\abs{a}\log\abs{\frac{1}{a}}} > 1
			\end{equation*}
			car $\abs{\phi_{a}'(0)} = 1 - \abs{a}^{2}, \abs{\phi_{a}(0)} = \abs{a}, \Re F(0) = \log\abs{1/A}$
	\end{description}
\end{proof}

\begin{théorème}
	{de l'Application Conforme de Riemann}{riemann-appli-conforme}
	Tout ouvert simplement connexe $U$ de $\C$, distinct de $\C$ est isomorphe à $D(0, 1)$. 
\end{théorème}
\begin{proof}
	D'après le lemme \ref{lemme:iso-ouvert-dans-compact}, quitte à effectuer une homothétie et une translation nous pouvons supposer $0 \in U \subseteq D(0, 1)$.\\
	D'après le lemme \ref{lemme:iso-disque-unité}, il suffit de prouver qu'il existe $f \in A$ telle que $\sup_{g \in A}\abs{g'(0)} = \abs{f'(0)}$. 
	On munit $\O(U)$ muni de la topologie de la convergence uniforme sur tout compact et on pose
	\begin{equation*}
		B = \left\{ f \in A \middle| \abs{f'(0)} \geq 1\right\}
	\end{equation*}
	Alors $B$ est non vide et bornée dans $\O(U)$ car pour tout $z \in U, \abs{f(z)} < 1$, donc pour tout compact $\sup_{K}\abs{f} < 1$. \\
	La partie $B$ est fermée dans $\O(U)$. En effet, pour $f \in \overline{B}$ et $f_{n}$ convergeant uniformément sur tout compact vers $f$, on a $f(0) = \lim f_{n}(0) = 0$ et de plus $f' = \lim f_{n}'$ d'où $\abs{f'(0)} = \lim \abs{f_{n}'(0)} \geq 1$. 
	Ainsi $f$ n'est pas constante et comme c'est une limite de fonctions holomorphes injectives, elle est injective. 
	Comme $\abs{f_{n}(z)} < 1$ sur $D(0, 1)$, $\abs{f(z)} \leq 1$ sur $U$ et d'après le principe du maximum, comme $f$ n'est pas constante, il n'existe pas de point $z \in U$ où $\abs{f(z)} = 1$. Donc $f \in B$. 
	Ainsi, $B$ est fermée bornée dans $\O(U)$ donc compacte, d'où l'existence d'un maximum à $g \mapsto \abs{g'(0)}$. 
\end{proof}

\begin{corollaire}
	{Homéomorphie d'Ouverts}{}
	Deux ouverts simplement connexes $U$ et $V$ de $\C$ distincts de $\C$ sont isomorphes.\\
	Deux ouverts simplement connexes $U$ et $V$ de $\C$ sont homéomorphes. 
\end{corollaire}

\end{document}

